\documentclass{book}
% to set font size use
% \documentclass[12pt]{book}

% https://www.overleaf.com/learn/latex/Tables
\usepackage{array}
\newcommand{\dilemmatablewidth}{5.5cm}

\usepackage{tikz} 

\usepackage{float}
% https://tex.stackexchange.com/a/20380/235813

\usepackage{graphicx}

\usepackage{lineno}

% https://www.overleaf.com/learn/latex/Indices
\usepackage{imakeidx}
\makeindex

% https://www.quora.com/What-font-are-most-books-printed-in
% to set font, use
%\usepackage{ebgaramond}

\usepackage{hyphenat} % http://www.ctex.org/documents/packages/special/hyphenat.pdf

\usepackage{setspace}\onehalfspacing\frenchspacing\flushbottom\sloppy

% https://www.scivision.dev/include-svg-vector-latex/
%\usepackage{svg}
% see https://tex.stackexchange.com/questions/442077/is-it-possible-to-use-svg-images-with-overleaf

% https://tex.stackexchange.com/a/8459/235813
\usepackage[nottoc]{tocbibind}

\usepackage{hyperref}
\hypersetup{
    colorlinks=true,
    linkcolor=blue,
    filecolor=magenta,      
    urlcolor=cyan,
    pdftitle={Bureaucracy Fieldguide},
    pdfsubject={How to be an Effective Bureaucrat},
    pdfauthor={Ben Payne},
    pdfkeywords={Bureaucrat, Bureaucracy},
    pdfpagemode=FullScreen,
    }

% https://en.wikibooks.org/wiki/LaTeX/Glossary says
% "\usepackage{glossaries} and \makeglossaries in your preamble (after \usepackage{hyperref} if present)"

% https://www.overleaf.com/learn/latex/Glossaries
\usepackage[toc]{glossaries}
%\usepackage{glossaries-extra}

\makeglossaries % The command must be before the first glossary entry.

% https://en.wikibooks.org/wiki/LaTeX/Glossary says
% "define any number of \newglossaryentry and \newacronym glossary and acronym entries in your preamble"
% see https://en.wikibooks.org/wiki/LaTeX/Glossary


\newglossaryentry{bureaucracy}{
    name=bureaucracy,
    description={definition here}
}

\newglossaryentry{visible bureaucracy}{
    name=visible bureaucracy,
    description={procedures and processes are written down and can be discovered by stakeholders}
}
\newglossaryentry{invisible bureaucracy}{
    name=invisible bureaucracy,
    description={procedures and processes are known to some stakeholders and are conveyed verbally to some of the other stakeholders.}
}

\newglossaryentry{process}{
name=process,
description={a task broken into a specified set of subtask dependencies.}
}

% https://graphthinking.blogspot.com/2021/07/bureaucracy-book-outline.html
\newglossaryentry{bureaucrat}{
    name=bureaucrat,
    description={a person responsible for subjective implementation of someone else's intent, with unquantifiable results. Examples of a bureaucrat role: teacher, police, government employee. Not bureaucrats: factory line worker, student}
}

\title{How to Be an Effective Bureaucrat:\\
A Field Guide}

% alternative titles:
% * "You are a bureaucrat" -- probabilistically true, but not always
% * "Everyone is a bureaucrat" -- false
% * "organization membership" -- bland
% * "distributed knowledge and distributed decision making"
% * "Distributed Decision Making"
% * "Be an Effective Bureaucrat: A Field Guide"
% * "How to be an Effective Bureaucrat: A Field Guide"
% * "Field Guide for Being An Effective Bureaucrat"
% * ""

\author{Ben Payne}
\date{\today}

\begin{document}
\pagenumbering{alph} % Fixes problem with glossary links; see https://tex.stackexchange.com/questions/119527/glossary-backlink-points-to-wrong-page

% https://texblog.org/2012/02/08/adding-line-numbers-to-documents/
%\linenumbers


%\setcounter{secnumdepth}{1}
% -2 results in no numbering for chapters, sections, or subsections
% 0  results in no numbering for sections or subsections
% 1  results in no numbering for subsections


\begin{titlepage}
\maketitle
\thispagestyle{empty}
\end{titlepage}
\clearpage

%\thispagestyle{empty}
\frontmatter % the front of the book has roman numerals

%\pagenumbering{gobble}
\thispagestyle{empty}

Copyright \copyright 2022 Ben Payne

Creative Commons \href{https://creativecommons.org/licenses/by-nc-nd/4.0/}{Attribution-NonCommercial-NoDerivs}

CC BY-NC-ND
\clearpage
\thispagestyle{empty}

Thank you to my coworkers. Our interactions helped me learn how to be a better bureaucrat.%\clearpage
%\pagenumbering{roman}

\hypertarget{contents}{}
\tableofcontents

% 2022-07-11: Brian
%\chapter*{Preface}% * excludes from Contents
\chapter{Preface}

A guide for situations you haven't been in before can decrease your anxiety by making the unfamiliar less intimidating. 
A guide labels patterns and provides context for your observations. 
Recognizing a familiar challenge can be preferable to the uncertainty of facing a novel scenario. An unfamiliar situation means not knowing how to start and not knowing the potential paths to success. In contrast, a challenge you understand can feel less burdensome. That doesn't make your life easy -- even a challenge that is understood requires creativity and effort.

Bureaucracy is a consistent and pervasive challenge in modern society. You have direct experience with bureaucracy and may be familiar with complaints about bureaucracy. You probably don't consider yourself a bureaucrat and are probably not an expert on bureaucracy.

%If you don't think of yourself as a bureaucrat, I hope to change your mind on this essential topic. 

Negative impressions of bureaucracy are common. The word
%bureaucracy 
% 2022-07-09: Jonah@upwork advocates removing this use
evokes reactions of disgust, primarily driven by how individuals have experienced bureaucracy.
The results of bureaucracy can appear idiotic.
Bureaucracy is even regarded as dangerous to society\footnote{``The only thing that saves us from the bureaucracy is inefficiency. An efficient bureaucracy is the greatest threat to liberty." -Eugene McCarthy (1979)}.
I have bad news for you if you have a negative impression of bureaucracy: bureaucracy is vital to society. 
% Effective bureaucracy is  to modern society.
Organizations have an insatiable need for members to perform a variety of tasks.

The good news is that you can learn to skillfully navigate bureaucracy. The consequence is that a skilled bureaucrat is more effective. That improvement enhances the organization you are a member of. Modern society is an interwoven fabric of organizations, with  bureaucracy being the crucial material.
Rather than ponder societal-scale concerns, this book focuses on the complex and typical experience of a bureaucrat. 


The framing of how you think about bureaucracy shapes what action you feel is relevant for changing your environment. 
Suppose your assumptions are wrong about how the system of bureaucracy works. Then you will be trying to convert the system to meet your incorrect expectations and your actions are likely to be wasted effort tackling irrelevant aspects. 
An accurate model of bureaucracy allows you to be more effective and your changes to the system are more likely to succeed. 

Having a precise definition of bureaucracy enables a more nuanced discussion. 
\begin{quote}
\textbf{Definition}: 
\ifglossarysubstitutionworks
\Gls{bureaucracy} 
\else
Bureaucracy
\fi
is the set of subjective decisions made by individual participants in an organization. The decisions originate from the management of shared resources. The resources can be either tangible or in the form of expertise. Bureaucracy relies on distributed knowledge and distributed decision-making. 
\end{quote}

The concepts and techniques in this book are derived from this definition. For more exploration of this definition, see the section ``\hyperref[sec:define-bureaucracy]{What is Bureaucracy?}''
\marginpar{See page~\pageref{sec:define-bureaucracy}.}

As an example of how the above definition is distinct from conventional definitions\footnote{\href{https://www.merriam-webster.com/dictionary/bureaucracy}{Merriam-Webster's definition} focuses on government; the \href{https://en.wikipedia.org/wiki/Bureaucracy}{Wikipedia entry}
\index{Wikipedia!\href{https://en.wikipedia.org/wiki/Bureaucracy}{bureaucracy}}
acknowledges non-governmental bureaucracy.}, I view government as an instance of where bureaucracy occurs. Bureaucracy is not limited to governments. I use the term ``organization'' to encompass governments, corporations, non-profits, or any collection of people working together. As a consequence, a vast majority of people are bureaucrats. 

The negative connotation of bureaucracy may lead you to reject the label of bureaucrat. Your recognition of your status as a bureaucrat matters even if you intentionally reject the label or do not recognize that you are a bureaucrat. How you behave and what you think your responsibilities are depend on how you label yourself.
% If you do not recognize that you are bureaucrat, you're less likely to be successful interacting with those around you. 
If you don't think of yourself as a bureaucrat, then you won't be able to distinguish what is your fault, what is the fault of your coworkers, what is the fault of management, and what is intrinsic to bureaucracy. 

Reasoning about bureaucracy is an accessible topic for every reader because we each have personal experiences.
A commonly held view is that the systems of bureaucracy we work within are broken: participants are harmed and there is inequality. 
The counterargument is that the bureaucratic behavior observed is how the system was designed to work. There is a third argument: bureaucratic systems work in a degraded mode -- they are less efficient than what is possible. 
In this book I set aside the macro-scale questions of bureaucracy like the three arguments above and focus on the human-scale interactions that give rise to bureaucracy.

Even people who are smart (e.g., they know history, they have memorized capital cities, they can do math) can struggle when faced with complex, large-scale distributed systems comprised of humans. Thinking about complex, large-scale systems is not part of the conventional education curriculum. This book will help you learn about bureaucracy, increase your ability to identify patterns, and apply relevant techniques.


The difficulties of bureaucracy are central to large societies that depend on complex interactions of many people. The challenge of bureaucracy is widespread across various governments and companies in different societies. And bureaucracy is durable -- it has existed for thousands of years. Because bureaucracy is pervasive, tackling the topic of bureaucracy is exciting and crucial. The excitement comes from the opportunities available for the many potential improvements.
%I enjoy pondering hard problems and then leveraging insights gained from reflection.
Being an effective bureaucrat is important because bureaucracy can be a \href{https://en.wikipedia.org/wiki/Force_multiplication}{force multiplier}
\index{Wikipedia!\href{https://en.wikipedia.org/wiki/Force_multiplication}{force multiplication}}
beyond what one person could accomplish on their own.

% Who this book is for

% from https://graphthinking.blogspot.com/2021/07/bureaucracy-book-outline.html
This book is for you if you are curious about the complex world we live in, or if you are thinking about how to contribute to society, or if you want impactful employment, or if your job is not what you expected. If you're wondering why innovation is suffocated within a bureaucratic organization, this book will to help you understand the challenges. Understanding the challenges means you can design ideas that are more likely to succeed.


% What you should expect reading this book: 
Every person in society is a participant in bureaucracy. The purpose of this book is to decrease the surprise of that experience and better prepare you emotionally and intellectually for the toil of being a bureaucrat. With focused reflection and a good guide, you can improve your skills as a bureaucrat. 

The topic of bureaucracy spans academic disciplines - \href{https://en.wikipedia.org/wiki/Psychology}{psychology},
\index{Wikipedia!\href{https://en.wikipedia.org/wiki/Psychology}{psychology}}
\marginpar{[Wikipedia] psychology}
\href{https://en.wikipedia.org/wiki/Sociology}{sociology},
\index{Wikipedia!\href{https://en.wikipedia.org/wiki/Sociology}{sociology}}
\href{https://en.wikipedia.org/wiki/Organizational_behavior}{organizational dynamics},
\index{Wikipedia!\href{https://en.wikipedia.org/wiki/Organizational_behavior}{organizational behavior}}
\href{https://en.wikipedia.org/wiki/Public_administration}{public administration}, 
\index{Wikipedia!\href{https://en.wikipedia.org/wiki/Public_administration}{public administration}}
\href{https://en.wikipedia.org/wiki/Anthropology}{anthropology},
\index{Wikipedia!\href{https://en.wikipedia.org/wiki/Anthropology}{anthropology}}
and  
\href{https://en.wikipedia.org/wiki/Political_science}{political science}. 
\index{Wikipedia!\href{https://en.wikipedia.org/wiki/Political_science}{political science}}
Most bureaucrats have no formal training in any of these domains. This book is not written from an academic perspective; it is written by a practicing bureaucrat to serve as a guide to be read by bureaucrats. 

This book is not a defense of bureaucracy, nor is it intended to disparage bureaucrats or the system of bureaucracy. Instead, this book is a guide to bureaucracy-as-it-is. There are no hacks presented for circumventing bureaucracy. Instead, this book provides advice on being an effective member of a bureaucratic organization.

This book does not focus on leadership, managing a team, being a team member, planning, time management, project management, advancing your career, promotions, or self-improvement. In the process of being a better bureaucrat some lessons may apply in those domains.


Nothing in this book is domain-specific, nothing is tied to engineering of products, and nothing is applicable solely in science research or policy development. However, skills associated with being an effective bureaucrat do translate into other domains.

This book doesn't address personal stress caused by bureaucracy, how to decrease stress, or how to balance the activities of work and life-outside-work.   However, understanding bureaucracy can enable a positive sense of self based on how your efforts integrate with other people's tasks. 

This book does not address work conditions, pay, benefits, or retirement plans. If you are seeking a competitive advantage that might result in improved outcomes, learning to navigate bureaucracy helps.


This book doesn't focus on citizens, political policymakers, oversight (e.g., Congress -- state or federal), competing organizations, contracts, the budget of an organization, fines and other regulatory punishments resulting from bureaucracy, policy enforcement, specific policies and regulations, formal arbitration or judicial resolution of disputes. For observations on those topics, see Wilson's Bureaucracy~\cite{1991_Wilson}. 


This book doesn't address discrimination or harassment. I neglect differences in experience based on race, gender, or socioeconomic status. This book doesn't address bad coworkers, abusive bosses, psychological defects of individuals, or malicious intent. I assume no traitors, spies, 
\marginpar{[Wikipedia] saboteurs}
or \href{https://www.hsdl.org/?abstract&did=750070}{saboteurs}\footnote{See also Wikipedia's entry for the \href{https://en.wikisource.org/wiki/Simple_Sabotage_Field_Manual}{Simple Sabotage Field Manual}
\index{Wikipedia!\href{https://en.wikisource.org/wiki/Simple_Sabotage_Field_Manual}{Simple Sabotage Field Manual}}
}.
I do not address \href{https://en.wikipedia.org/wiki/Dark_pattern}{dark patterns}; 
\index{Wikipedia!\href{https://en.wikipedia.org/wiki/Dark_pattern}{dark patterns}}
issues like lying, bribery~\cite{2021_Ang}, and criminal behavior are not discussed. All of those aspects happen whenever people interact; those challenges are not specific to bureaucracy. 
Bureaucracy is not a manifestation of incompetence, malicious actors, or mistakes. Sources of friction in a well-run bureaucracy stem from perpetual challenges like ambiguity, finite attention, and inadequate resources.


In this book I assume you are honest, and that other people are honest.  
I assume bureaucrats strive for fairness, though the particular definition of fairness may vary from bureaucrat to bureaucrat. 
Even with all these simplifying assumptions the complexities of bureaucracy arise. My guidance on how to be an effective bureaucrat is in the context of benign bureaucracy.


% What is the benefit of reading this book?
As a result of reading this book, you will be better able to recognize and navigate complex human interactions, both within your career and outside of work. Perspectives offered in this book can benefit you directly, whether by promotion of rank or title, increase in pay, successful completion of a project, or through decreased stress derived from an improved understanding of how bureaucracy works. Being a more effective bureaucrat can also positively impact the causes you care about and the people you engage with.





% there's no avoiding the issue
Everyone is a bureaucrat because there are few alternatives to bureaucracy for a society constrained by finite resources. Gaining skills in navigating bureaucracy is helpful both for your happiness and the improved functioning of society. That may seem abstract and lofty, but the situation is inescapable: you are a member of society and benefit from participating in society. 

Hoping that modern technology will eliminate or reduce bureaucracy is not a solution. Automation and computers merely obfuscate processes and make negotiation more challenging. Even with faster decisions and fewer humans, there is still reliance on humans to make decisions and design processes.

Another approach to avoiding bureaucracy is to characterize interactions with other people merely as personal relationships. That is easier than considering bureaucracy as a complex system.
% this is also stated on page 24 of~\cite{1991_Wilson}
This ``just relationships'' view misses emergent phenomena and over-simplifies the situation. As a result your effectiveness is decreased.


% Caveats

%In my reflections and trying to draw lessons there is a risk of overanalysis. Sometimes a situation is merely happenstance, and sometimes trying to extract lessons from randomness is folly. Avoiding conjecture about conspiracy and malice is a fuzzy boundary when insufficient information is available. 

%My experiences cannot be generalized to every situation. Some of the observations here may be analogous to your context if you squint. 


%While this material is intended to be timeless and generic, it is culturally specific to the United States of America in the early twenty first century. 
%There are cultural blindspots not addressed in this book because I did not encounter systemic hurdles in my career as financially privileged white male. 
Focusing on the morals of organizations or bureaucrats\footnote{Jackall's Moral Mazes~\cite{2009_Jackall}} is less compelling than studying bureaucracy because norms (i.e., expectations for individuals and organizations) shift over time, while bureaucracy is persistent. 
In this book I focus on culturally invariant aspects of bureaucracy. 

% Source of this content: 
This book is based on personal experience, reading published books and papers, and anecdotes from other bureaucrats. Surveys are not the foundation supporting the claims made. 
\marginpar{[Wikipedia] blinded-\\experiment}
No \href{https://en.wikipedia.org/wiki/Blinded_experiment}{double-blind} 
\index{Wikipedia!\href{https://en.wikipedia.org/wiki/Blinded_experiment}{blinded experiment}}
experiments were conducted. 
% my experience
% I wrote this book for a younger version of me.
% When I first started my job in a large organization I recognized differences between the expectations of the education system I had left and the challenges of a professional environment. 
I have attempted to learn from my mistakes by reflecting on my (in)actions and the consequences. This approach has been an expensive education. My mistakes delayed progress and damaged relationships. This book aims to provide generalizations from my experiences that might benefit the reader.


% How the book should be read: 
Reading this book front-to-back is the default option, but not essential. The dependencies of topics are not sequential -- see Figure~\ref{fig:core-concepts}.

\begin{figure}[ht]
    \centering
    \includegraphics[width=1\textwidth]{images/core_concepts_map.pdf}
    \caption{Relation of the concepts related to bureaucracy discussed in this book. }
    % TODO: include chapter/section in diagram
    % TODO: generate diagram from Latex
    \label{fig:core-concepts}
\end{figure}

% 2022-03-27: BHP was not impressed by the options on https://texample.net//tikz/examples/tag/graphs/ including the "mindmap" option
% this looks better: http://randomresearchdata.blogspot.com/2013/12/tikz-simple-decision-tree-or-cart-tree.html
% see https://texample.net/tikz/examples/feature/trees/

%\begin{tikzpicture}[main/.style = {draw, circle}] 
%\node[main] (1) {Core Concepts}; 
%\node[main] (2) [above right of=1] {Decision-making}; 
%\draw[->] (1) -- (2);
%\end{tikzpicture} 

The approach in this book is to view bureaucracy from a few perspectives. The first is a conceptual view focused on the \hyperref[sec:define-bureaucracy]{definition of bureaucracy} 
\marginpar{See page~\pageref{sec:define-bureaucracy}.}
in terms of managing shared resources. The second view is structural (e.g., hierarchy, meetings). The \hyperref[sec:fundamentals-of-b]{essentials of bureaucracy} 
\marginpar{See page~\pageref{sec:fundamentals-of-b}.}
will be familiar to experienced bureaucrats. 
The third view is based on the experience of a practicing bureaucrat -- 
\hyperref[sec:dilemma-trilemma]{dilemmas}
\ifhaspagenumbers
(page~\pageref{sec:dilemma-trilemma}), 
\else
, 
\fi
\hyperref[sec:unavoidable-hazards]{unavoidable hazards}\ifhaspagenumbers
(page~\pageref{sec:unavoidable-hazards}), 
\else
, 
\fi
and 
\hyperref[sec:effective-bureaucrat]{tips on how to be effective}\ifhaspagenumbers
(page~\pageref{sec:effective-bureaucrat}).  
\else
. 
\fi

%The three productive views of bureaucracy (concept, structure, experience) are contrasted with~\hyperref[sec:alternative-views-from-within]{how other bureaucrats view bureaucracy} and~\hyperref[sec:models-of-bureaucracy]{other models of bureaucracy}.

Chapters~\ref{sec:introduction} 
\ifhaspagenumbers
(page~\pageref{sec:introduction}) 
\fi
through~\ref{sec:why-bur-hard} 
\ifhaspagenumbers
(page~\pageref{sec:why-bur-hard}) 
\fi
provide background context for bureaucracy. 
Chapters~\ref{sec:individual-in-org} 
\ifhaspagenumbers
(page~\pageref{sec:individual-in-org}) 
\fi
through~\ref{sec:process} 
\ifhaspagenumbers
(page~\pageref{sec:process}) 
\fi
are the field guide. 
Each chapter of the field guide can be read independently. 



\ \\

% as per https://tex.stackexchange.com/q/393238/235813
\begin{flushright}
Ben Payne\\
\today\\
United States of America\\
Earth
\end{flushright}


\clearpage

\mainmatter % the main part of the book will have standard pages

% https://tex.stackexchange.com/questions/2958/why-is-newpage-ignored-sometimes

\chapter{Introduction to Bureaucracy}
% essentials

%  Section~\ref{sec:define-bureaucracy} defines bureaucracy and associated terminology for specific roles. That is contrasted in section~\ref{sec:alternative-views-from-within} with how bureaucrats view themselves and their interactions. In section~\ref{sec:models-of-bureaucracy} a summary of external perspectives of bureaucracy is provided. 
  
  

  \section{What is bureaucracy?\label{sec:define_bureaucracy}}

Bureaucracy is the specialization of roles necessitated by scaling and complexity for the distribution of common resources or widespread policy

In the course of carrying out someone else's subjectively defined policy, you have to make your own subjective decisions in the execution and enforcement in the course of carrying out someone else's subjectively defined policy, you have to make your own subjective decisions in the execution and enforcement


While you may know it when you see it or experience it, for this book definitions are useful. There are three distinct roles in bureaucracy: policy creator, policy enforcer, and the person upon whom policy is inflicted.

The existence of bureaucracy is independent of an organization's purpose.


The policy creator is either a politician or a bureaucrat. 

A \gls{bureaucrat} is a person subjectively interpreting policies on behalf of an organization and has discretionary enforcement to facilitate coordination of stakeholders. 

Let's break that down piece-by-piece. First, ``subjective interpretation'' means there is a person making a decision about how to do something. The subjectivity arises from different reasons one might choose an option over a competing option.  ``Policies" is a set of actions in a given circumstance. ``An \gls{organization}" is the collection of people for who the policy is made. ``Discretionary enforcement'' means the person is choosing how to apply the policy in the specific circumstances. ``Facilitating coordination'' means bureaucracy is about getting multiple people (or sometimes a person at different instances in time) to work together. The ``stakeholders'' is a group of people who care about the application of the action in each circumstance.  That's still pretty dense, so the rest of the book is spent expanding on the nuances and implications of this definition.

Bureaucracy is neither good nor bad. Bureaucracy is not tied to politics, or any specific institution (corporations, governments, academics). Bureaucracy is not defined to be efficient nor, does it have to be inefficient. Bureaucracy is not restricted to paperwork, or record keeping, or quantification, or gathering metrics. 

Bureaucracy is about delegation of control, communication, decision making, coordination, and processes. Involves negotiation, primarily informal. 

An organization comprised of bureaucrats is a \gls{bureaucracy}. The definition of bureaucracy used in this book is independent of government. Nothing in this definition involves paperwork or an office building. Definitions that limit the concept of bureaucracy to specific contexts result in a decreased ability to describe complex large-scale human organizations. 

The protagonist within a \gls{bureaucracy} is the \gls{bureaucrat} -- the person who is a member of an organization and is responsible for subjective implementation of policy for the organization. The person that a bureaucrat's decisions are inflicted on a \gls{subject}.  Depending on context, a subject may be a student (when the bureaucrat is a teacher) or a subject may be a citizen if the bureaucrat is a police officer or government official. Sometimes a bureaucrat's decisions are inflicted on other bureaucrats-as-subjects, such as when a Chief of Police creates guidelines for police in their district, or when a senior diplomat sets policy for embassy employees. 

A critical aspect of bureaucracy is that everything is made up, specifically by other humans. The consequence is that everything is negotiable. You (in the role of either a subject or a bureaucrat) need to know who to negotiate with and how to negotiate the desired changes. The only actual rules are mathematical physics that describe nature. Everything else is either naturally occurring macroscopic emergent phenomena (e.g., chemistry, biology) or humans making up labels and norms. 

Bureaucracy arises when there is no common objectively quantifiable feedback mechanism for individual participants in the organization. This aspect is why governments, schools, and prisons are characterized as bureaucratic. The military doesn't rank soldiers by ``number of enemies killed'' and is bureaucratic. Even profit-driven commercial organizations are bureaucratic when the impacts of individual employees are not coupled to sales metrics. 

Profit-based feedback makes some roles in a business context slightly more predictable and understandable, though there are still trade-offs like long-term profit versus short-term profit and externalization of harm. 

The concept of bureaucracy is most visible for complex, long lasting, and recurring situations involving many people. The apparent friction can be lower when there are only a few people involved (``I'm just talking to my collaborator" or ``I'm just buying groceries from a clerk at the store'' or ``I'm using a website for a government service''), but there is a continuous gradient. 

There is the external resource (mail delivery for USPS, public safety for FBI, environment for EPA) and there are resources internal to the bureaucracy. The focus of this book is on internal resources. In that context, bureaucracy is for the disseminated responsibility for use of resources: attention, skill, expertise. Time, money, staffing are proxy measures.



A useful way to think about bureaucracy is as a system for distributed knowledge and distributed decision making. That is in contrast to easier-to-understand concepts like centralized knowledge and centralized decision making. A government run by dictatorship is easy to conceptualize compared to democracies because there is a central character around which a narrative can be formed. Similarly, telling stories about the \href{https://en.wikipedia.org/wiki/Chief_executive_officer}{CEO} of a company is much easier than capturing the thousands of interactions conducted by the many employees of that company. Linear story-telling with a small number of protagonists does not map well to the complexities of bureaucracy. 
% are there alternatives to Bureaucracy that accomplish the same non-centralized non-consensus approach to complexity?


% https://graphthinking.blogspot.com/2017/09/market-friction-and-bureaucratic.html
Distributed knowledge and distributed decision making are hindered by
\begin{itemize}
    \item limited bandwidth between people, specifically the bureaucrats involved
    \item non-zero latency of information between people, specifically the bureaucrats involved
    \item the cost of measurement (getting data)
    \item the cost of analysis of the data
    \item making decisions that are suboptimal
\end{itemize}


\subsection{Bureaucracy as an economics model}
Firms exist in a market because negotiating contracts and prices for every interaction is burdensome. 
% https://www.kellogg.northwestern.edu/faculty/hubbard/htm/research/ec174/lectures/3coase.htm

Doesn't address small vs large companies, and doesn't distinguish between profit-oriented and non-profit and government. 

\subsection{Bureaucracy as emergent phenomenon}
Bureaucracy as a set of many bilateral interactions may not need to invoke emergence. However, there's a universality that hints at emergence. 

Above the threshold for emergence, there is scale-free behavior. The same patterns are observable at large organizations and extremely large organizations.

All those choices faced by the individual are not independent choices with respect to other bureaucrats in their environment. There is a flocking behavior of my choices are informed by the choices of those around me. Not necessarily in space.

Everyone is playing by different rules and has different objectives and everything is dynamic (both individuals and the conditions). 

Bureaucracy as a macroscopic phenomenon is emergent at sufficient scale. The scale is important because there is no longer dependence on individual relationships (beyond \href{https://en.wikipedia.org/wiki/Dunbar\%27s_number}{Dunbar's number}. There are people in the organization that you don't know and for which there is no common accountability. An organization subdivided into team recursively until there is local person-to-person accountability.  

The local rules bureaucrats employ to enable distributed decisions using distributed knowledge is meetings, processes, and communications. 

The relevance of making a claim that something is emergent is that there is behavior occurring at the macroscopic scale, and Knowing that individual motives and actions of every player at the microscopic level is not relevant.

% https://www.preposterousuniverse.com/podcast/2021/10/11/168-anil-seth-on-emergence-information-and-consciousness/
What does ``emergent'' mean? Nominal emergence example: a circle is emergent from a collection of points. Weak emergence is measurable using \href{https://en.wikipedia.org/wiki/Granger_causality}{Granger causality} or, equivalently\footnote{https://arxiv.org/abs/0910.4514}, \href{https://en.wikipedia.org/wiki/Transfer_entropy}{transfer entropy} (information theory). 


\subsection{Bureaucracy in Game Theory}
Bureaucracy does not fit cleanly into game theory categories of cooperative or competitive.

Maybe all the interactions within a bureaucracy are a bunch of small games?

Bureaucracy is self-modifying. 

Bureaucracy is in constant flux due to external conditions, externally imposed constraints, staff turn-over, internal dilemmas, disagreements of individuals. 


\subsection{Bureaucracy resists characterization}
Actually, bureaucracy is worse than emergent - the system rules can be altered or ignored by the stakeholders. \href{https://en.wikipedia.org/wiki/Wicked_problem}{Wicked problem}. This is why coming up with a holistic theory of bureaucracy is difficult. 

As soon as a claim is made, then a group can respond to that claim by behaving in an opposing manner. 

\subsection{Money as a fitness function}
Commercial businesses have a different accountability -- money. Common across all participants within the organization, and common with external stakeholders. The goal of a company is to generate profit. Commercial businesses have people who make subjective decisions and enforce policies, but there is a common metric for feedback. The feedback mechanism is not perfect. Being a good commercial bureaucrat does not necessarily result in monetary success.

Prisons, schools, medical, government, military all consume and spend money, but money isn't the goal. When faced with a decision, choice is not guided by which will generate more profit. 


\subsection{Bureaucracy as evolutionary outcome}

Biological, Genetic -- individual level
Biological, Genetic -- Group selection
Memetic

\subsection{Bureaucracy as Psychological Phenomenon}

Are you doing what's best for you, the group you're in, or everyone?
Altruistic or reciprocal? Retaliation
The answer changes time the time and situation of situation and person to person

Just a mixture of pathologies? % section
  \clearpage 
  \section{Scope of Bureaucracy}
  Bureaucracy is not a distant faceless concept. Bureaucracy permeates your environment and is the basis for many interactions you have with other people. 
  %\begin{flushright}
  %Back to the \hyperlink{contents}{Table of Contents}
  %\end{flushright}
    \subsection{Who is a bureaucrat?}

A cashier in a gas station is a bureaucrat. The ``policy'' might simply be ``take money from customer in exchange for items and gas,'' but the subjective application of that policy leaves a lot of room for the cashier to shape the customer's experience. Does the cashier greet the customer when the customer enters the store? Does the cashier look at the customer to acknowledge the customer? Smile? How quickly does the cashier engage the customer? Minor nuances that are left to the cashier in the execution of the store policy means there is room for subjective application of the policy. 

This same discretionary application of policy applies to commercial bureaucrats like sandwich makers, car salespeople, grocery clerks, retail clerks. Teachers, police, tax collectors, and other state works are government bureaucrats. 

% bureaucracy is not limited to white collard office workers
Factory line workers subjectively apply policies. Enforcement of quality standards is the most obvious area. Pacing of work is a negotiation.

Sometimes bureaucrats do not work directly with customers or citizens or products. Then the bureaucratic process is inflicted on fellow bureaucrats. In this scenario, a bureaucrat is subjectively applying a policy to other bureaucrats. 

Identifying yourself as a bureaucrat matters, both to the employee and to the business. The risk of not self-identifying as a bureaucrat is that you won't grasp how much control you have in implementing and enforcing policy. If you think of yourself as having to blindly follow rules, you will harm the people you are applying the rules to and you will harm the business/institution you are applying the rules for. Adapting policies to circumstances is the value of having judgement capacity. 

In a similar sense from the consumer/citizen perspective, if you don't think you are interacting with a bureaucracy, you won't perceive the opportunity to negotiate.  If you view rules as fixed and inflexible, you will harm your ability to make progress. If a rule was made by a human, then that rule is flexible. Who made the rule? Who enforces the rule? If you can talk to them, could they be convinced to make a modification or an exception?
 % subsection
    \subsection{Number of people in a bureaucracy}
Although bureaucracy can be present for one person, and bureaucracy is often apparent on teams (e.g., 3 to 20 people), this book primarily focuses on the situation of multiple teams comprising an organization. This might be a few hundred people (above \href{https://en.wikipedia.org/wiki/Dunbar's_number}{Dunbar's number}) up to millions of people. 

There are tens of companies that employ more than a million people\footnote{see \href{https://en.wikipedia.org/wiki/List_of_largest_employers}{Wikipedia's list of largest employers}}, including Walmart, Amazon, and McDonald's. Bureaucracy emerges at a far smaller scale, is scale invariant, and is generic across sectors. 

Small companies with a few people incur bureaucracy. Non-profit organizations encounter bureaucracy. The complexity of the tasks may be different, but the same scale-independent patterns can emerge because of a common factor: human behavior.

Size of bureaucracy scales with the complexity of the problem. 

Scaling a bureaucracy up is easier than scaling down. % subsection
  %\section{From Startup to Bureaucracy}

In a novel domain, people show up in order: nerds, then commoners, then jackasses


  \clearpage
  \section{Fundamentals of Bureaucracy\label{fundamentals_of_b}}
  
Bureaucracy as a system of distributed knowledge and distributed decision making for management of shared resources relies on particular facets of human interaction. I outline the roles of decision making, hierarchy, and communication as central to bureaucracy. My list is smaller than that of Weber \cite{2015_Weber}\footnote{\href{https://en.wikipedia.org/wiki/Bureaucracy\#Max_Weber}{Bureaucracy: Max Weber}}.

Decision making (\S\ref{sec:decision-making}) is central to bureaucracy. The management of resources requires policies to be created. 

Coordination of decision making within an organization of people motivates formal and informal hierarchies. Who gets to make which decision is managed using hierarchy; see \S\ref{sec:hierarchy_of_roles}.

Meetings and written communication help establish consensus among bureaucrats.
Once a decision is made, the choice selected by a person have to propagate throughout the organization to achieve consistency; see \S\ref{sec:meetings-for-coordination} on meetings and \S\ref{sec:written-communication} on written communication.
 % section
    % define the core concepts 
    \subsection{Decision making in a bureaucracy\label{sec:decision-making}}


Decision are not the only source of change in an organization. Occasionally events unfold without decisions being made. This might be because the decision makers are not informed, or there is an \href{https://en.wikipedia.org/wiki/Willful_blindness}{intentional neglect}. This section focuses on situations where bureaucrats recognize the need for a decision and want to make the best decision.

There are multiple types of decisions. 
A \gls{simple decision} has one correct or beneficial choice and one or more wrong or harmful choices. The work of decision making is then to gather information that identifies which is the correct or beneficial choice and select that option.

The best case scenario for any decision making is one person making a well-informed simple decision that has immediate consequence and the consequence is to the decision maker. Examples from elementary formal education include arithmetic math problems, multiple choice quizzes, spelling tests, and memorization tests. A bureaucrat's \href{https://en.wikipedia.org/wiki/Moral_injury}{moral injury} can come from decision making that involves multiple people, weak feedback loops, and complex decisions.

\subsubsection{Example Decision Method: Pareto Frontier\label{sec:pareto}}

A complex decision may have many choices, and there are might not be a best option. Then a \href{https://en.wikipedia.org/wiki/Pareto_front}{Pareto frontier} might exist where trade-offs can be made. 

As an example of a complex decision made by one person with immediate consequence and direct relevance to the decision maker, suppose you want to purchase a car. You car about only two aspects: fuel efficiency and cost. 

\begin{figure}[ht]
    \centering
    \includegraphics[width=1\textwidth]{images/pareto_frontier_car_options.pdf}
    \caption{Four cars, L, M, N, and P. The goal of the buyer is to spend less money and get better fuel efficiency. Choices not on the frontier should be avoided, but that doesn't yield a single result.}
    \label{fig:pareto_frontier_cars}
\end{figure}

Visualzing a Pareto frontier for two quantitative variables is easy, but typically decisions involve more factors. For example, evaluating the trade-off of three quantitative variables like speed, accuracy, and cost creates a surface. With more than three variables visualization are less useful, though the analysis technique still applies. 

Another constraint on using Pareto frontier analysis is that it works well when there are many options relative to the number of variables being optimized for. 
The assessment does not work as well when there are few choices relative to the number of variables. For example, suppose there are 10 choices of car and you want high fuel efficiency, sufficient cargo capacity, maximum number of passengers, stylish, low cost, low maintenance, good durability, and high resale value. 

For a set of quantitative variables, a Pareto frontier does not account for relative importance of different variables. Assign weights to each of these factors merely stretches one axis relative to the other axes. 

There are many possible decision making frameworks besides Pareto frontiers, but in practice a typical bureaucratic decision is ill-informed, has diffuse consequences, delayed impact, and does not affect the decision maker. In bureaucratic processes there is rarely a formal assessment of options. 
Decisions are rarely recorded. 
Even afterwards a decision can be difficult to evaluate for correctness because there are multiple stakeholders.


\subsubsection{Risks of Using Decision Frameworks}

Decision making frameworks can be attractive to bureaucrats intending to formalize processes (see \S\ref{sec:process}) and encourage predictability. There are potential risks worth being aware of. 

% https://graphthinking.blogspot.com/2019/01/political-decisions-versus-science.html
A decision is political when the basis is historical relationships, maintenance or creation of a relationship, or to enable future relationships. A decision is subjective when someone else faced with the same scenario would have come to a different conclusion.
A decision is quantitative when it is based on measurements. To avoid the appearance of subjective decision making or political decision making, a decision may be framed as ``data driven." 
% https://graphthinking.blogspot.com/2018/06/data-driven-decisions-versus-data.html
A good approach for data driven quantitative analysis involves coming up with a testable hypothesis, then performing experiments and collecting data to evaluate the hypothesis. More commonly, a decision is made, then data is gathered which supports the desired outcome. Forming an opinion and then looking for evidence to back the outcome yields suboptimal results for the organization.

Even if a bureaucrat is not intentionally biased towards an outcome, there are many ways to gather evidence and some approaches have biased sampling and produce biased results.

With valid and representative data measurement, decision makers can be led astray by poor modeling. A model may use inapplicable techniques, or may have implementation bugs.

For all the dangers of decision making methods, there are worse approaches that do not rely on measurement. People rely on history (if they are aware of it) and perpetuate bad ideas, or take action based on what is best for their career, or decide based on how to accumulate more power, or simply choose based on what someone else says to do.  


\subsubsection{Decision Making Delay}

What appears from the outside as ``organizational inertia'' is internal delay of decision making and the delay of dissemination. 
Delay comes from
\begin{itemize}
    \item It takes time for each decision maker to gather information, arrive at a decision, modify processes, disseminate their selection, justify their selection. 
    \item forcing a continuous variable into a discrete set of choices. Typically the number of choices is small. Discrete choices for a continuous variable is a loss of effectiveness.
    % https://dynomight.net/teaching/
    \item processes designed to account for cheaters and people with malicious intent, whether that means a malicious bureaucrat or malicious subject. 
\item Analysis paralysis, due to {insufficient information, too much information, which framing is unclear}
\item When other people who are needed to carry out the action push back, either in disagreement or seeking clarification. The fact that the organization is not profit driven is important because the justification for the action isn't quantitatively obvious. Therefore there's a higher burden for communication.
\end{itemize}


% The following characterization has no consequence
% Decision making by bureaucrats can be informal or formal, consensus-based or solo. 


\subsubsection{A Bureaucratic Decision involves many Decisions}

A decision is actually a collection of dependent choices. After recognizing the need for a decision, follow-on decisions include identifying the stakeholders (who to include in an impact analysis) and identifying options. Who is a stakeholder and what are the options are interrelated. Involving more people expands the number of options and the complexity. Additional choices associated with reducing uncertainty are how much time to spend on the decision, how much information to gather for the decision, whether the make the decision or push the decision to someone with more expertise, whether to push the decision to someone with more exposure to the consequences.

Most decisions you make as a bureaucrat do not have hard deadlines. Instead, there are trade-offs in allocation of your time. Sooner is preferable since the consequence of the decision benefits the organization and allows you to focus on other tasks, but delaying allows for more information gathering for a better informed decision. (See Dilemma~\ref{table:gather_data_lots-vs-little} and other related Dilemmas in that section.)


If a bureaucrat is going to rely on expert consultation (see \S\ref{sec:expertise} on expertise), the decision maker needs to be confident the expert is not of straying outside their area of expertise. For example, I don't rely on a botanist with many published papers to tell me how to change the oil in my car. 
Besides knowing their own limitations, the expert should be clear about whether their input is a factual summary, a predictive assessment, or a value judgement. This nuance complicates what you as a bureaucrat are interested in (what's the best choice?).

 % Process and policy
    \subsection*{Hierarchy of Roles\label{sec:hierarchy-of-roles}}


In an ideal situation, sufficient depth and breath for decision making would be embodied in one person. That might not be possible in every situation. One way to resolve this is to identify distinct scopes of responsibility and then assign different members of an organization separate scopes for decision making. Within a decision making scope there may be more work than one person can handle, so a team is formed. That team may have some members focused on tactical work and other members focused on strategy and coordination. Hierarchy within an organization is the formalization of separate decision-making scopes and associated specialization. 

Partitioning knowledge and decision making enables complex work beyond what one person can accomplish and causes friction among members. An expert reporting to a manager knows things the manager does not, and the manager may have context that the expert lacks. Both bureaucrats (the expert and the manager) need to convey their respective understanding and seek the holistic view.

Hierarchical decision making is one option for coordination among alternatives (like consensus), so why is hierarchy so common? Members of an organization gravitate towards hierarchy because it helps define task scope, assigns responsibility, and obviates a need for building consensus. Reaching consensus for every decision would take time and be more burdensome than appointing a person as the decision maker.

\ \\

A hierarchical organization with partitioned knowledge introduces a challenge: the order in which you share information with others matters. Your choices for who to first describe an idea to are your peers, your management, and your subordinates. 
\marginpar{[Tag] Trilemma} 
\index{trilemma!communication priority: peers, management, subordinates}
The people subordinate to you know more about the topic and are exposed to the consequences. Giving them a chance to vet the idea results in a more robust idea and validates their value in the organization. Alternatively, first sharing your idea with management  allows your superiors to provide context you might not be aware off. And choosing to first start the conversation with your peers indicates you value the relationship and decreases the risk of duplicating work.

\ \\

I've included hierarchy in the section on Fundamentals of Bureaucracy, but that does not imply that hierarchy is a required feature of bureaucracy. Hierarchies of bureaucrats are a common \href{https://en.wikipedia.org/wiki/Organizational_structure}{organizational structure} and are worth studying even if not essential to bureaucracy. The relevance of understanding hierarchy is to identify recurring behavior and patterns to leverage.
Organizations of bureaucrats can intentionally work against the use of hierarchy for decisions, but the amount of effort needed to enable alternatives results in hierarchy being a common approach.

The benefits of formal hierarchy include improved capacity for the number of policy decisions made, enabling consistency of decisions, and leveraging specialization of knowledge. 
Hierarchical decision making has costs: higher latency (compared to a signle decider), inconsistency among bureaucrats (dissemination isn't perfect), waste due to inefficiency, and 
\hyperref[sec:unavoidable-hazards]{many others}.
\ifsectionref
described in section~\ref{sec:unavoidable-hazards}. 
\fi
As another example of the potential for harm, hierarchy enables strategic ignorance. Bureaucrats in positions of power can deny having knowledge of improper activity\footnote{L.~McGoey, ``The Unknowers: How Strategic Ignorance Rules the World" (2019)
% review: https://www.tandfonline.com/doi/abs/10.1080/19460171.2020.1768422?journalCode=rcps20
and 
L.~McGoey, ``The logic of strategic ignorance" (2012). DOI 
10.1111/j.1468-4446.2012.01424.x
}. 



The structure of an organization is dynamic, but at each point in time an organization typically has a defined set of roles. Each role is distinguished by different scopes of decision authority. 
Roles are often confused with titles. What matters is the role (scope of decisions) and who reports to whom. The names of teams can be similarly not descriptive.




Roles in an organization are defined by the boundaries of responsibility. The purpose of a role is to minimize conflict, decrease the need for coordination, reduce redundancy, and allow for control of resources. Clear responsibility boundaries enables effective bureaucracy. 


A conventional characterization of an organization's hierarchy involves two criteria: the depth and breadth of the org chart.
The more people a supervisor oversees, the flatter the organization -- that's the breadth of the organization. The depth of the hierarchy is how many layers there are. See the Valve handbook~\cite{2012_Valve} and Joreen's essay~\cite{1972_Joreen} for contrasting views on the merits of an organization's hierarchy. 

A more practical view of an organization's hierarchy also involves two criteria. The two choices of how a hierarchy is shaped are 
1) how many people a supervisor oversees and 
2) how many supervisors a person has. 
Though you might naively expect that an employee has one boss, but that is \href{https://en.wikipedia.org/wiki/Matrix_management}{not a requirement}. A supervisor for a given topic may have many people reporting to them, and a bureaucrat with multiple roles may report to more than one supervisor.

\ \\

Acting as part of a group means ceding part of your autonomy. Hierarchy cedes more of your responsibility and adds expectations about relationships.
The consequence of hierarchy in an organization is that, as a member of the bureaucracy, you do not have full autonomy -- otherwise you would not be a member of the hierarchy. At the same time, you are not under strict control of the organization -- you still have some subjective decision making authority as a bureaucrat.

The person at the top of the hierarchy does not know everything. The person at the top of the hierarchy does not have input on every decision made in the organization. Some autonomy is retained by all members of the bureaucracy.

Independent of the defined roles and titles in an organization's hierarchy, there are a set of implicit roles and a separate social hierarchy of informal influencers and decision makers. Informal influencers in a bureaucracy usually have long relationships with the decision maker or relevant credentials or both. The credentials can be formal (e.g., a \href{https://en.wikipedia.org/wiki/Doctor_of_Philosophy}{PhD}) or informal (demonstrated success on a project). In either case, the decision maker is relying on another person's expertise. 

Another set of informal relations within an organization is mentors and mentees. These relations allow mentors to share institutional knowledge to mentees, and allows people in senior positions to access the novice perspective. 


\ \\

One consequence of hierarchy is a sense of fear felt by people who report to other people. This fear stems from the loss of control (less autonomy) that leaves the person feeling disempowered. 

For example, consider the following relationship. Sue, is perceived to have power over another person, Amy, because Amy gave up some control to Sue. Amy not having full control over decisions triggers the feeling of fear in Amy, regardless of how Sue behaves. Having complete responsibility for decisions also induces anxiety.

If Sue is aware of the potential for this emotional experience, Sue can compensate for Amy's fear by being friendly and receptive towards Amy. Alternatively Sue may exploit or rely on the fear felt by subordinates. Sue not noticing or accounting for Amy's fear does not invalidate Amy's emotional experience.


% Active bystander when the person doing wrong is in a position of authority
% PACT (Probe, Alert, Challenge, Take Action)
% https://mobile.twitter.com/GeorgetownABLE/status/1408498438203969541


% Mintzberg's Coordination Mechanisms
% https://www.youtube.com/watch?v=IZET8VjSifQ


    \subsubsection*{Organizational chart as a Guide and a Lie\label{sec:org-chart-as-guide-and-lie}}

An \href{https://en.wikipedia.org/wiki/Organizational_chart}{organizational chart} 
\index{Wikipedia!\href{https://en.wikipedia.org/wiki/Organizational_chart}{organizational chart}}\iftoggle{WPinmargin}{\marginpar{[Wikipedia] Organizational\\chart}}{}
(hereafter an ``\gls{org chart}'') identifies formal roles and the formal relations among roles. An org chart is at best a snapshot in time, and more often aspirational than descriptive. Despite possible deficiencies, an org chart helps outsiders and newcomers understand the scope of responsibilities and interactions.\footnote{The organization chart hasn't always existed. The \href{https://en.wikipedia.org/wiki/George_Holt_Henshaw\#First_organization_chart}{first known org chart} 
\index{Wikipedia!\href{https://en.wikipedia.org/wiki/George_Holt_Henshaw}{George Holt Henshaw}}
was created in the 1850s.}

Org charts are a lie because undocumented relationships can matter more than official roles. Org charts fail to capture the informal roles and network of relations that facilitate progress in any organization. Org charts document titles instead of describing roles.

Org charts foster a second separate lie by creating a sense of power dynamics based on visual orientation. For more on this issue see the discussion of~\hyperref[sec:org-chart-orientation]{org chart orientation}\iftoggle{haspagenumbers}{ on page~\pageref{sec:org-chart-orientation}}{}.

    \subsection{Meetings for coordination\label{sec:meetings-for-coordination}}
In an organization comprised of more than one person, meetings are a necessary artifact for facilitating coordination. The coordination accomplished by a meeting can be explicit (verbal or written), or it can be indirect through signaling (who attended the meeting, when the meeting was held, where the meeting was held, how much notice was provided). 

Meetings are a vital aspect of coordination in a bureaucracy. Tips are in \S~\ref{well-run_meeting}.


% https://graphthinking.blogspot.com/2021/07/thought-terminating-concepts-in.html
\textit{No plan; I just do what you tell me.}\\
Employee: I just do what you tell me to do. With that approach, either I will be successful because I worked hard on what you directed, or I will fail because I was directed to do the wrong thing by you.\\
Manager: Is that how you want your career to go? I think you're smart, and I think you're capable of shaping your career.

\ \\

\textit{No point in making a plan}\\
Employee: There is no point in making a plan, because everything changes so frequently.\\
Manager: But the software has an end goal, right?\\

Employee: Yes, so then the plan is to get from where we are now to that end goal.\\
Manager: And there are no intermediary steps? Milestones?

Manager: Is it better to have no plans and just put up fires reactively, or to have a plan that is subject to change?

\ \\

\textit{What is a plan anyways?}\\
Manager: I think there is value in creating a goal, enumerating tasks that would support the goal, identifying the dependencies among the sub-tasks, and time-binning the dependencies with defined milestones and deliverables. That is my definition of a plan. And having that is more useful than merely reacting.

\ \\

\textit{Who's plan?}\\
Employee: I don't need to come up with that plan, you already have a plan. Just tell me what the plan is.\\

Manager: That's not as effective as coming up with independent plans and then resolving the differences. There's value in resolving the differences, Even though that will cost time and frustration and displace time to implement.
    \subsection{Written Communication\label{sec:written-communication}}

Reports, memos, emails are artifacts of bureaucracy in an organization. A written record creates evidence about policies and decisions. Existence of a record can be used for good or for harm.


% https://graphthinking.blogspot.com/2020/09/identifying-and-eliminating.html




%***************************

\subsubsection{Why written communication does not happen\label{sec:written-comm-does-not-happen}}
\begin{itemize}
    \item Too many reports/emails/memos/messages to read and process and respond to. Emotionally overwhelmed.
\item The person may read slowly.
\item The person may type slowly, or their hand writing is poor.
\item Fear of imperfect communication. What if the email is incomplete or inaccurate or ambiguous?
\item The person views written communication as ``official" or ``plan of record'' and does not feel comfortable brainstorming or creating contingency plans
\item The person wants to avoid accountability for their statements
\item The person may not be confident in their writing ability -- spelling, grammar, sentence composition, structuring content. \footnote{For tips on writing, see 
\href{https://en.wikipedia.org/wiki/The_Elements_of_Style}{The Elements of Style}
and
\href{https://www.google.com/search?q=dodm+5110.04}{DoDM 5110.04, Manual for Written Material}}
\end{itemize}



%***************************

\subsubsection{Email: A Piece of Art, A Form, and A Game\label{sec:}}

Originality is not a requirement in bureaucratic writing. Plagiarism is acceptable. The consistency of madlibs-based forms that follow decision trees is efficient. 

Readers need both context and conciseness. 

Using visual cues to improve readability

Email is a game of documenting decisions and responses so that the sequence of interactions is clear for audits and recrimination. 

\subsubsection{Email Responsiveness\label{sec:}}

three tiers of email responsiveness:
1) read all emails and reply to all emails
2) read all emails and reply to some
3) unable to read all incoming emails

The following scenarios are focused on case 2 -- able to triage (skim) emails but not enough time to respond to all.
****************************
Given insufficient time, which email do you reply to? An email from your boss, an email from your peer, an email from a person who reports to you, or a person who you do not know?
The ratio of these emails is not one to one to one to one.

1) the email from your boss is likely the top priority
of the remaining emails (peer, subordinate, unknown), 
2) peer is likely the next priority
3,4) the subordinate and unknown person are likely last

When there's insufficient time available, your transparency to subordinates and unknown people is likely to decrease

****************************
Three emails come in. One has no action, one is easy to reply to, and the third is difficult and takes time. Which one gets the response?
The ratio of these emails is not one to one to one.

1) the email that is easy to reply to gets answered first
2) the email that is difficult is second

If you only have time for one of the emails, it's likely to be the easy one. 
\href{https://en.wikipedia.org/wiki/Ambiguity_aversion}{ambiguity aversion}

As a consequence, outsiders see this as you are bikeshedding
\href{https://en.wikipedia.org/wiki/Law_of_triviality}{law of triviality}

\subsection{Email Structure}

% from https://graphthinking.blogspot.com/2021/10/structuring-email-content-for.html

If implementing all these tips sounds like a lot of work, that's because it is. Effective written communication requires intentional effort because of the lack of augmenting channels (compared to voice or video or in-person). 



Use consistent design and structure for your emails. Emails are part of your professional reputation.

Emails start with a greeting: Hi, Hello, Good morning, Good afternoon, Good evening. 
Email greetings include the name of the targeted recipient(s). 

Emails terminate with a professional closing, e.g., "Kindly", "Regards", etc

Emails contain a signature block with contact information -- phone number, normal hours of response, which timezone you're in if your team spans timezones, how long to wait for a response before asking again, which communication channel I prefer, etc.

\begin{figure}
\includegraphics[width=1\textwidth]{images/email_template.pdf}
\caption{Template for new email messages. Greeting has a space after the comma -- that is where the recipient's name will go. Signature block uses smaller after the name.}
\label{fig:email_template}
\end{figure}

Email signature blocks do not include unnecessary images, as that uses more storage for recipients. 
Email threads focused on a specific instance of a recurring event include the date (YYYY-MM-DD) in the subject line. 

Based on the purpose of the email, example key phrases for subject lines include: "meeting notes" vs "agenda" vs "question about"

Revising an existing subject line can disrupt the ability of email software to thread conversations. However, sometimes the revision is worth breaking threading.

When replying to an ongoing thread, retain the original message as part of the thread to provide readers historical context.

When replying to threads with sensitive messages, sanitize the included content as appropriate by removing name or identifying details.

If an email contains multiple requests or questions, at the top of the email (after the greeting) explicitly say how many of each type. Then, in the body of the message, number them.

\begin{figure}
\includegraphics[width=1\textwidth]{images/email_two_questions.pdf}
\caption{Distinct items the recipient should address in a reply.}
\label{fig:email_two_questions}
\end{figure}

If an item corresponds to a requested action, separately highlight the action and indicate who is supposed to take the action and what the deadline for response is.

\begin{figure}
\includegraphics[width=1\textwidth]{images/email_meeting_notes.pdf}
\caption{Who has what action due when?}
\label{fig:email_meeting_notes}
\end{figure}

Computer commands should be distinct separate fixed-width font. This distinguishes the text from the rest of the narrative. 


\begin{figure}
\includegraphics[width=1\textwidth]{images/email_computer_font.pdf}
\caption{The computer commands use fixed width font. The author distinguished input from output through the use of bold vs non-bold respectively. The author highlighted the error message using red. Inline text like cat in the last line is also fixed with.}
\label{fig:email_computer_font}
\end{figure}

References to documents include a direct full path.

If referring to a previous separate thread, include the subject and the date+time that email was sent

For bullet points, explicitly specify that items are joined by one of the following: OR, XOR, AND

If you have an unordered list, explicitly state that order is irrelevant.

If you have a sequence of steps, number them appropriately and indicate which steps are required versus optional

Use visual sketches to illustrate concepts rather than always relying on text. Don't use pictures all the time, and don't have too many pictures in an email. 

Know how to both embed pictures inline and how to attach files and when to use which. 
Email replies should preferentially be at the top of the thread. 
If replying to multiple points in the previous email, embed replies inline, mark the distinction, and highlight the authorship. 

\begin{figure}
\includegraphics[width=1\textwidth]{images/email_reply.pdf}
\caption{Bob's reply to Sue's questions. The third question is not shown in this illustration.}
\label{fig:email_reply}
\end{figure}

If replying inline, explicitly state that at the top of the thread.

The email trilemma is to balance the amount of detail against providing sufficient context and being concise. 

If the email requires a \href{https://en.wikipedia.org/wiki/BLUF_(communication)}{B.L.U.F} or \href{https://en.wikipedia.org/wiki/Wikipedia:Too_long;_didn\%27t_read}{tl;dr} or summary, then it's too long. Inclusion of a BLUF or tl;dr risks resulting in the reader skipping the content. 

Emails convey both emotional tone and facts. Your intent is practically irrelevant; the reader's perception is paramount. 

Every email should have a purpose. What are you asking the recipient to do? How do you want them to feel? How should they respond?

When replying, starting your email with an expression of gratitude for the work the recipient has done so far sets a positive tone by acknowledging their investment.

 
    \subsection*{Feedback loops and Ripples\label{sec:feedback-loop-and-ripples}}

%When there is a lack of input from customers or users, then someone has to decide.


A feedback loop exists when a decision-maker experiences the harms and benefits of their decision. Decisions that lack a feedback loop still have consequences for the decision-maker, in that potential future decisions are altered or limited \footnote{In~\cite{1983_Lipsky} Lipsky discusses feedback loops in chapter~4 on page~40.}.


Virtuous cycles and vicious cycles are rare in bureaucratic organizations because there are rarely mechanisms for feedback loops. Instead, there are \glspl{ripple} -- propagation of consequences for other people's schedules, altering what is possible for other people, and creating dependent tasks to be carried out by people other than the decision-maker.


The weak feedback mechanisms for a bureaucrat are reputation (subject to spin)
\index{reputation}
and rarely enforced retroactive accountability in the form of the question, ``What did you know and when did you know it?''
%Reputation is social.
Retroactive accountability depends on written records from meeting notes, emails, agendas, and reports. Reducing the potential for retroactive accountability is one motive for bureaucrats avoiding written records for decisions and policies.


Because of the lack of quantitative feedback loops for bureaucrats, there is significant interest in documenting justification, proactive monitoring, reporting, and retroactive assessment. Each of those activities creates more administrivia.


When there are multiple competing objectives among stakeholders in a 
\href{https://en.wikipedia.org/wiki/Zero-sum_game}{zero-sum}
\index{Wikipedia!\href{https://en.wikipedia.org/wiki/Zero-sum_game}{zero-sum game}}
\marginpar{[Wikipedia] zero-sum\\game}
use of resources, how can you determine what's best? In some domains there are feedback loops to guide progress. When feedback loops are weak or not present,  the most powerful stakeholder (potentially distinct from the biggest or loudest) will dominate. 


\ \\

%****************************************
% https://graphthinking.blogspot.com/2020/01/hierarchy-of-justification.html

Because feedback loops are weak for decision-makers, alternative mechanisms are needed in bureaucracy. A standard approach is the use of approvals and justifications. The strength of justification needed for a given action depends on your relationship with the approver and the potential ripples associated with the action. 

The amount of justification for an approval process should be just enough so that when something goes wrong there's a clear reason why. 

Ordering the quality of a justification
\index{decrease surprise!quality of justification}
\begin{enumerate}
    \item I have no explanation.
    \item This is my opinion.
    \item We've always done it that way (a \href{https://en.wikipedia.org/wiki/Cargo_cult}{cargo cult}).
    \index{Wikipedia!\href{https://en.wikipedia.org/wiki/Cargo_cult}{cargo cult}}
    \marginpar{[Wikipedia] Cargo\\cult}
    \item Based on my experience.
    \item I was told to do it this way.
    \item I think this is the best way (no reasoning, but a desire to optimize; optimization criterion undefined).
    \item This way is most effective because X (where X is not quantified, there is a desire to optimize, and there is an optimization criterion).
    \item This way is most effective because X (where X is quantified, there is a desire to optimize, and there is an optimization criterion).
    \item This way is most effective because X compared to other options (where X is quantified, there is a desire to optimize, and there is an optimization criterion).
\end{enumerate}

%****************************************

\ \\

% TODO: need a transition between topics


\subsubsection*{Special interest groups}

Within bureaucratic organizations special interest groups care about aspects of the shared resource central to the organization. 
Inefficiency can occur when there is a benefit to a small group and the cost is to a larger group.
% TODO? tie in with Public choice theory

The \href{https://en.wikipedia.org/wiki/Social_trap}{social trap}
\index{Wikipedia!\href{https://en.wikipedia.org/wiki/Social_trap}{social trap}}
\marginpar{[Wikipedia] social\\trap}
is ``a conflict of interest or perverse incentive where individuals or a group of people act to obtain short-term individual gains, which in the long run leads to a loss for the group as a whole."
The feedback loop for the diffused value relevant to the organization is weaker than the feedback loop for value to the interest group. 

\subsubsection*{Spending taxpayer dollars}

As an example of a weak feedback loop, consider the scenario of a government employee deciding how to spend government money.

%\marginpar{[Tag] Story Time; Math}
\index{story time!taxes and spending}
\begin{storytime}{Taxes and Spending}
Suppose you are a government bureaucrat and earn \$100,000 with a tax rate of 30\%. Then your tax money sent to the government is \$30,000.
How does that compare to what the government collects in taxes?

For the United States, ``in 2021 the government collected \$4.05 trillion in revenue."
\footnote{source: Government Revenue | \href{https://datalab.usaspending.gov/americas-finance-guide/revenue/}{U.S. Treasury Data Lab}}

That means your taxes of \$30,000 would be
30000/4050000000000 = 0.00000074\% of the tax base for the country.

If you are a federal government bureaucrat and you do not maximize the effectiveness of spending \$1,000,000 of government money, of that misallocated money only \$0.0074, or about one penny, was taxes you paid. The financial feedback loop is weak.

Some federal government bureaucrats earn slightly more, so the cost of wasting \$1,000,000 is more significant.
Federal pay is limited to about \$220,000\footnote{Wikipedia entry on \href{https://en.wikipedia.org/wiki/Executive\_Schedule}{Executive Schedule}
\index{Wikipedia!\href{https://en.wikipedia.org/wiki/Executive\_Schedule}{executive schedule}}
%%%CANTDOINFOOTNOT\marginpar{[Wikipedia] Executive\\schedule}
}, raising the feedback to 2 pennies.
\end{storytime}

The lack of feedback allows waste to go unfelt. There's no immediate consequence for the decision-maker.
%Waste is indistinguishable from not enough funding or insufficient skills.

\ \\

Example feedback loop:
\begin{center}
\begin{figure}[ht]
    \centering
    \includegraphics[width=0.8\textwidth]{images/feedback_loop_complexity_and_staffing}
    \caption{Increased complexity requires more staffing to enable specialization. More staff means more skills are available; under-utilized staff skills make room for more scope; more scope adds to complexity.}
    \label{fig:complexity_and_staff_growth}
\end{figure}
\end{center}

\subsubsection{Example feedback loop: Security agent making a fair decision}
% https://graphthinking.blogspot.com/2017/09/a-simple-illustration-of-bureaucracy.html
%\marginpar{[Tag] Story Time}
\index{story time!airport security}
\index{exemplar!\href{https://en.wikipedia.org/wiki/Transportation_Security_Administration}{Transportation Security Administration (TSA)}}
\begin{storytime}{Airport Security Line}
A coworker and I were going through passport control at an airport. 
A \href{https://en.wikipedia.org/wiki/Transportation_Security_Administration}{TSA officer} 
%\index{exemplar!\href{https://en.wikipedia.org/wiki/Transportation_Security_Administration}{Transportation Security Administration (TSA)}}
\index{Wikipedia!\href{https://en.wikipedia.org/wiki/Transportation_Security_Administration}{Transportation Security Administration}}
%CANTDOWITHINMDFRAMED\marginpar{[Wikipedia] TSA}
directing passengers into one of two lines. Both lines were long. The length of each line was not equal even though the entrance to the lines was at the same location. Each line terminated at a row of passport-checking agents. Each passport check took a minute. Passport-checking TSA officers operate independently and concurrently.


The TSA officer's perspective is that there are two long lines. Her procedure is two balance the two lines (for fairness). He does this by pointing people into one of the two lines, with her choice driven by which line appears to have room available.

My coworker and I enter the controlled area and are directed into lines by the TSA officer. The officer directs my coworker left and directs me to go right. The TSA officer completed her job.

My coworker in the left line finishes 5 minutes before me. This difference in completion time is frustrating for me.

Because the lines are not of equal length, balancing the start of the line is a suboptimal method. The consequence is that what seems fair to the TSA officer ends up not being fair for people going through the line. The TSA officer had incomplete information -- the lines are not of equal length. Because the TSA officer isn't exposed to the consequence of her approach, she doesn't get feedback on whether it is suboptimal or not.
\end{storytime}

Lesson: if the people making decisions do not experience the consequences of those decisions, then they have no incentive to improve decision-making.

The person in the longer line feels frustrated. The negative feeling is due to a sense of powerlessness, the situation is recurring, and a better solution is available.

The optimal solution in this situation is to have a single line feeding the multiple TSA passport checkers. A single line eliminates the need for the decision-maker but incurs work to change the status quo.


\subsubsection{Parking garage}
% source: 
% https://graphthinking.blogspot.com/2019/07/altering-feedback-loops-to-change.html

%\marginpar{[Tag] Story Time}
\index{story time!parking garage}
\begin{storytime}{Parking Garage}
Bob parks his car in a parking garage every day. 
The parking garage owner charges \$20 per day for people to park their car.

Bob recently found that one of the exit gates for the parking garage is broken. If Bob uses that gate to leave the parking garage, the gate does not function and Bob cannot exit. Then Bob has to call the parking gate operator to request an exception (which is granted) and Bob can then exit that gate, avoiding the \$20 per day charge.

This action (go to broken gate, request exception, avoid charge) has been repeated for a long time (months). Bob's motive is to avoid the \$20 parking charge; the cost is a mere phone call and a minor delay. This cheating behavior harms the parking garage owner's income. However, the parking gate operator, serving as intermediary, insulates the parking garage owner from interaction with the cheater. The cheating behavior is small enough that the parking garage owner may not notice.
\end{storytime}

Having an intermediary policy enforcer alleviates the need for the garage owner to interact with customers. If the incentives of stakeholders are not aligned then inefficiency goes unchecked. If the parking garage operator's profits are a fixed rate rather than tied to parking charges, then the operator lacks the motive to fix problems.  % subsection
  \clearpage
  \section{History of Bureaucracy\label{sec:history}}


Bureaucracy has repeatedly arisen independently in various cultures
\footnote{See the Wikipedia entry on the \href{https://en.wikipedia.org/wiki/Bureaucracy\%23History}{history of bureaucracy}.
\index{Wikipedia!\href{https://en.wikipedia.org/wiki/Bureaucracy\%23History}{history of bureaucracy}}
}
lasting for timescales that exceed the lifespan of one person.\footnote{See the YouTube video on the \href{https://www.youtube.com/watch?v=B_nsZlcC12g}{History of bureaucracy}.} That indicates the current situation is not a fluke or coincidence. There is some utility (or pathology) that is consistently recurring. 


Bureaucracy predates writing and language and even humans! Subjective policy enforcement in support of an organization arises in pre-human tribes, visible in groups of modern apes who have to manage access to shared resources~\cite{2016_Suchak}. 



Though bureaucracy is not new, the pervasiveness is. Before the industrial revolution the scale of employment and government were small, with organizations limited by the speed of communication. For the past 100+ years the size of organizations (commercial, governmental, and academia) have grown beyond \href{https://en.wikipedia.org/wiki/Dunbar\%27s_number}{Dunbar's number} -- the number of human relationships you can maintain, about 150. \iftoggle{WPinmargin}{\marginpar{[Wikipedia] Dunbar's\\number}}{}
\index{Wikipedia!\href{https://en.wikipedia.org/wiki/Dunbar\%27s_number}{Dunbar's number}}
More people participate in more organizations that are more bureaucratic. Driving this increase is the support for more complex products and processes. 

% claim: bureaucracy grow faster than the growth of human population?


\clearpage % section

\chapter{Bureaucracy in general\label{sec:bureaucracy-in-general}}
  \section{Each Phase of Life Involves Bureaucracy}

\href{https://en.wikipedia.org/wiki/TL;DR}{TL;DR}: 
\index{Wikipedia!\href{https://en.wikipedia.org/wiki/TL;DR}{TL;DR}}
You have experience with bureaucracy, though you may not have framed the experiences as bureaucratic. Viewing relationships and roles through the lens of bureaucracy explains the limitations of your expectations.

\ \\

In each person's life there are standard milestones: birth, education, work, death. 
%The relationships associated with each phase are distinct.
Each of these milestones and phases involves bureaucracy. Each phase is a different experience of bureaucracy because the bureaucratic roles change.

Bureaucracy is composed of three roles: the policy creator, the policy enforcer, and the subject (upon whom the policy is inflicted). 
\index{bureaucracy!roles}
These roles can be confounded by oversubscribing the label ``bureaucrat." Who is the bureaucrat and who is the subject depends on the relationship in a scenario. 

For example, a store manager creates a policy, a store clerk enforces the policy, 
\index{exemplar!store clerk}
and the policy is inflicted on the customer. Both the manager and the store clerk would be bureaucrats, while the customer is the subject. In a separate example, someone at corporate headquarters sets a policy, the store manager enforces it, and the policy is inflicted on the store clerk. Then the clerk was the subject of bureaucracy. 


\subsection*{Bureaucracy of Birth\label{sec:bureacracy-of-birth}}
Your birth was marked by getting a name, registering with the state, and initiating medical records. These tasks were administered by bureaucrats (doctors, nurses, and other hospital staff) on your behalf, and you were the subject of the bureaucracy. You had no autonomy or decision-making authority. 

\subsection*{Bureaucracy in Early Childhood\label{sec:bureaucracy-early-childhood}}
Before starting formal education, the bureaucracy of early childhood is inflicted primarily by family members setting and carrying out policies. The organization of bureaucrats is the family, with the shared resources being housing, food, and experience with survival. Other community members or caretakers may also be involved in carrying out the policies of taking care of you. Your decision-making authority as a subject in this bureaucracy was minimal. 

\subsection*{Bureaucrats at School\label{sec:bureaucracy-of-school}}
Once you started the formal education process, new bureaucrats got involved.  The community of bureaucrats could be a public school, a private school, or homeschooling. In any of those cases, the frontline bureaucrat is the teacher. You're not responsible for making policies that other people follow; you are still the subject of bureaucracy.


The expectations of each phase of school (high school, undergraduate, graduate school) are distinct, and they are different from working in a large organization. Your autonomy increases throughout the duration of school. %, and your ability to make policy that effects others grows. 
Your family and teachers are the bureaucrats. You start building informal organizations of friends, and you start to explore policies around social bonding.

Schooling sets a pattern that most students will fall into for the rest of their lives: you were handed a textbook and told to solve a set of problems. That pattern of taking direction can persist for a long time. However, you are not constrained to persist with that limitation. You have the autonomy to do more than what is required. You can find other textbooks that match your interests or are written from a different view. 

You get to choose the book you want, even if you don't get to choose the topic you will be evaluated on. You don't  need to pick just one reference book -- you can review lots of books and figure out which author style best fits you. You can also choose the level of difficulty -- basic and beginner level, or more advanced. You can choose what's best based on your understanding rather than defer to what a class in school is supposed to cover.

As a subject of the education bureaucracy, you can discover how you learn best. This extra effort requires self-reflection: How do you learn? What works best? What didn't work, and why not? What did you learn? What do you wish you had learned?

Another pattern that schooling relies on is single-question decisions with only one right answer. Examples include math problems and multiple-choice tests. Schooling tends to avoid setting up dilemmas or paradoxes for students. Academic problems in the education process are designed to be independent of the people involved or the history of the situation. 


% https://graphthinking.blogspot.com/2013/02/all-little-things.html
%Where you sit in the classroom matters. Being in the front means you will be exposed to fewer distractions, more likely to pay attention.


% https://graphthinking.blogspot.com/2012/09/how-to-not-be-average.html
%How were you taught? Did you have any input on the method?
%Who taught? Did you like them? Were they friendly, knowledgeable, and approachable?


% https://graphthinking.blogspot.com/2012/09/how-to-not-be-average.html
%What resources are available now for you to learn from? Do you like learning?
%you have the freedom to pursue what ever intellectual endeavors you want.


% https://graphthinking.blogspot.com/2011/09/which-skills-are-useful-after.html

% did not prepare me for addressing challenges at work. 
In my schooling I learned how to approach technical issues and develop solutions. That problem-solution paradigm neglects crucial steps of discovering the problem, isolating the challenge, identifying stakeholders, learning the history of the challenge, and negotiating with stakeholders before trying to address the challenge. 

My schooling led me to emphasize academics over socializing. When I transitioned to professional work, I found social skills and political savvy useful when trying to change organizations and policies. 

Academic problems are intended to be solvable and answers submitted get evaluated. In contrast, working in a bureaucratic organization the challenges are ill-defined, there's no known solution, and the topic is sufficiently complex that you have to collaborate.

% https://graphthinking.blogspot.com/2018/07/the-difference-between-problems-at.html


%\subsection*{undergrad vs graduate}
% education process roles and expectations vary over time

%college, graduate school: friends, teachers, advisors



\subsection*{Military Service Bureaucracy\label{bureaucracy-of-military}}
\index{exemplar!military}
Somewhere between 13 percent\footnote{\href{https://www.cbpp.org/research/federal-budget/where-do-our-federal-tax-dollars-go}{https://www.cbpp.org/research/federal-budget/where-do-our-federal-tax-dollars-go}} and 20 percent\footnote{\href{https://www.nationalpriorities.org/analysis/2019/tax-day-2019/where-your-tax-dollar-was-spent-2018/}{https://www.nationalpriorities.org/analysis/2019/tax-day-2019/where-your-tax-dollar-was-spent-2018/}} of the United States federal budget is spent on military, and 
less than 0.5\% of the United States population serves in the military.\footnote{\href{https://www.cfr.org/backgrounder/demographics-us-military}{Council on Foreign Relations}\iftoggle{boundbook}{, https://www.cfr.org/.}{.}}
For those who serve, the military's rigid hierarchy and defined protocols are a distinct experience compared to school or work. Transitioning from military to civilian life can present a dissonance for service members used to the chain of command and clearly defined orders. 

While you are in the military the hierarchy and use of inefficient policies feel stifling. Once out of the military you may reflect fondly on the clarity of orders compared to the vagaries of social politics.  

\subsection*{Working in a Bureaucratic Organization\label{sec:bureaucracy-of-work}}
%Within the employment phase of life, there are pairs of events which may apply: hiring or getting hired; firing, getting fired, or quitting. 

%Employment: managers above you, peer employees, people you manage are all members of an organization. 


Organizations with communal workspaces have \glspl{shared resource}: bathrooms, conference rooms, kitchen areas with fridges and microwaves, storage areas. Each of these incurs policies of use. Unlike being a student at school, you may find yourself responsible for developing and enforcing policy. 

As an example of workplace policies, consider the following scenario. 
%\marginpar{[Tag] Story Time}
\index{story time!bathroom smells}
%\begin{storytime}{The Bathroom Stinks}
\begin{mdframed}[frametitle={The Bathroom Stinks},frametitlerule=true,frametitlealignment=\centering]
The bathroom at work sometimes smells, so I'm a nice person and bring a scented air freshener. Unbeknownst to me, that triggers asthma or an allergic reaction in one of my coworkers. Because of this a policy gets created so this mistake doesn't happen again. 

Signs are posted. Violations are reported to management even if no one has a physical reaction to the air freshener.
%\end{storytime}
\end{mdframed}

Policies are often created in response to specific incidents. This intent can be helpful (promulgating lessons learned is efficient) or unhelpful (when policies are an overreaction). 

\subsection*{Healthcare and Death\label{sec:bureaucracy-of-death}}
Medical care alters our life, and a lot of money is spent on medical care: 25\% of the federal budget in the United States.\footnote{\href{https://www.cbpp.org/research/federal-budget/where-do-our-federal-tax-dollars-go}{https://www.cbpp.org/research/federal-budget/where-do-our-federal-tax-dollars-go}} Understanding the bureaucracy of healthcare is outside the scope of this book, but your role in the bureaucracy is helpful to understand.

Doctors, nurses, and other staff are bureaucrats; you are the subject; the hospital or clinic is the organization. 

\ \\
To illustrate bureaucracy of a large organization, consider the importance of toilet paper at a hospital. 
%\marginpar{[Tag] Story Time}
\index{story time!toilet paper at hospital}
%\begin{storytime}{Restocking Toilet Paper}
\begin{mdframed}[frametitle={Restocking Toilet Paper},frametitlerule=true,frametitlealignment=\centering]
If you go to a hospital, use the bathroom, and find there is no toilet paper, that would indicate a deficiency of the hospital.

The \textit{holistic view} is that someone didn't refill the toilet paper. Since the person who usually restocks toilet paper wasn't also a user, they aren't directly affected by the lack of toilet paper.

A routine is needed for checking the availability of toilet paper in bathrooms. 

The \textit{perspective of the purchasing manager} is that money spent checking the status of toilet paper is money not spent on the hospital's primary mission: improving the health of community members.

Minimizing checking of toilet paper is important for the organization's reputation, so a feedback mechanism is instituted: a phone number is posted in the bathroom so users can send a text regarding bathroom status at the hospital.

The \textit{perspective of the janitor} is that my routine used to be to go to each bathroom after normal business hours and refill toilet paper. Now I have to do that and be on-call when someone alerts management that service is needed. My responsibilities increased but my pay did not.

%\end{storytime}
\end{mdframed}

%\ \\

%Death can invoke both the medical system and the government. 

%\ \\


%All of these roles and relationships involve 

Bureaucracy is a set of subjectively administered policies within an organization. By recognizing the role of bureaucrats you can identify what is negotiable. 
 
  \clearpage
  \section{Avoiding Bureaucracy is Nearly Impossible}
\iftoggle{shortsectiontitle}{\sectionmark{Avoiding Bureaucracy}}{}

The only situation where bureaucracy might not exist is if you live entirely independently and have no interaction with other people. That means completely disengaging from society. Even then, personal routines are a self-imposed form of bureaucracy, with the roles of policy maker, bureaucrat, and subject collapsed to a single person -- you.

Self-sufficiency and autonomy are attractive alternatives to bureaucracy. The way participants in modern society strive for self-sufficiency is by denying their dependence on modern society. That's a relabeling of selfishness which feels better. 

For the rest of us who operate as members of  society, bureaucracy is necessary for our rights. You validate your name using paperwork, forms, and records. These artifacts are used, in cooperation with other people, to determine your claim of citizenship and associated rights. 
That's a subjective policy that \iftoggle{glossarysubstitutionworks}{\glspl{stakeholder}}{stakeholders} in society agree to. 


Bureaucracy is necessary because it is a response to the \iftoggle{WPinmargin}{\marginpar{$>$Wikipedia: Collective action problem}}{}%
\href{https://en.wikipedia.org/wiki/Collective_action_problem}{collective action problem} -- everyone would benefit from cooperation, but each person has to sacrifice their self-interests.  
\index{Wikipedia!collective action problem@\href{https://en.wikipedia.org/wiki/Collective_action_problem}{collective action problem}}
As long as humans form communities, we will address the challenge of \iftoggle{glossarysubstitutionworks}{\glspl{shared resource}}{shared resources}
 -- whether tangible (e.g., water, land, air) or intangible (e.g., expertise, information). 
That is why bureaucracy is culturally invariant and persistent across time.
Learning to be an effective bureaucrat improves your chances of success in society. 

The specific way society is constructed (democratic, authoritarian, dictatorship, anarchy) is irrelevant -- bureaucracy is still present. Even the libertarian approach of relying on contract enforcement implies some bureaucracy (e.g., forums for resolving contract disputes like a court system, enforcing decisions through violence). 


Not all bureaucracy is due to the state, nor is bureaucracy confined to companies. Parenting involves creating situation-specific requirements for children, with the organization being the family as mentioned in the section 
on \hyperref[sec:bureaucracy-early-childhood]{bureaucracy in childhood}.
\marginpar{See page~\pageref{sec:bureaucracy-early-childhood}.}%
Dress codes for sports teams are arbitrary standards. 
Store clerks are bureaucrats, as are website forum moderators.  Content moderation is the process of (inconsistently) enforcing arbitrary standards. This mindset even permeates individuals as internalized expectations of policy and enforcement when no one else is present. 

Recognizing instances of bureaucracy enables more skillful interaction, whether as a bureaucrat or subject. The rest of this section illustrates both the view of a person interacting with bureaucracy as a \gls{subject} and the perspective of bureaucrats working within organizations. 




 
  \clearpage
  \section{Alternatives to Bureaucracy don't work\label{sec:alternatives-to-bureaucracy}}
A common refrain for participants in an organization is ``Can't we just do the work?'' In this section we'll explore a few alternatives to bureaucracy. These thought experiments illustrate the naive intent and the practical deficiencies of each alternative.

The models of bureaucracy below are all imprecise because they are analogies. The solution is a precise definition of bureaucracy.


% https://graphthinking.blogspot.com/2021/02/how-to-have-efficient-bureaucracy.html

\subsection*{Efficient Bureaucracy}

In an ideal scenario with no bureaucracy, everyone comes to the same conclusion when presented with the same information regarding management of shared resources. Efficient bureaucracy requires each person to know the skills of everyone else so that anyone can act as an expert in any field. 

In this scenario the management process of building consensus becomes unnecessary. There is no need to fight over resources (money, staffing) and no need to fight over direction. A potential problem is that this model of efficient bureaucracy based on everyone having the same view may lack the diversity of views that enable robustness.


While that idealized scenario is not going to happen, it points to how to improve bureaucratic efficiency. Each person has or should have access to the same information. 
Each person should apply the same decision making process consistently.
Every person should have the same incentives.

Reasons bureaucracy is inefficient include
\begin{itemize}
    \item Not everyone has the same information.
    \item Processes are inconsistent.
    \item \hyperref[sec:motivations]{Incentives vary} among bureaucrats.
    %(see section~\ref{sec:motivations}).
% https://graphthinking.blogspot.com/2017/03/what-slows-down-bureaucracy.html
    \item Bureaucrats and Subjects use imprecise language.
    \item Decision Makers use opinions and experience rather than collect data and analyze data.
    \item People prioritize putting out fires rather than attacking largest issues.
    \item People do not respond to communications (email, phone, meetings); see the section on \hyperref[sec:slow-deployment]{slow deployment}. %section~\ref{sec:slow-deployment}.
    \item Participants are late to meetings.
    \item \href{sec:scope-creep}{Scope creep} for projects.
    \item Planned scope and actual scope are mismatched  (due to staffing or skills).
    \item Teams work in silos.
% https://graphthinking.blogspot.com/2017/03/progress-in-spite-of-better-ways.html
    \item Bureaucrats and Subjects lie and employ other dark patterns (not addressed in this book).
    \item People make mistakes.
    \item Each person's reference experiences are unique.
\end{itemize}
Consider the following \href{https://en.wikipedia.org/wiki/Thought_experiment}{thought experiment}. 
%What if everybody in a bureaucracy were the same?
What if everybody in a bureaucracy had a different opinion? How would consensus be arrived at? Can an organization operate without having to form consensus for every decision? 
\marginpar{[Tag] Exercise}



\subsection*{Perfect Bureaucracy}

The definition of ``perfect bureaucracy'' depends on who is providing the perspective. Here I illustrate both the view of the subject and the view of the bureaucrat. 

From the perspective of the subject of a bureaucratic organization, perfect bureaucracy means many things: minimizing the time the subject waits on a decision, getting perfectly correct decisions, decisions are consistent across subjects and circumstances, decisions take into account all relevant factors, and there is zero cost to the subject. Any process deviating from those expectations is a noticeable detriment to the subject, leading to a negative reference experience of bureaucracy. Never mind that the desires are unreasonable and conflicting. In a real bureaucracy, subjects have a negative or neutral experience. Positive interactions with bureaucracy are rare and are regarded as anomalous.

Perfect bureaucracy for a bureaucrat means all information is available, information is immediately available, and there is no moral ambiguity (each answer is objective). Trade-offs become trivial and the emotional toll of the work goes to zero. Then the bureaucrat is able to serve subjects (an emotionally rewarding prospect). Note that this job might feel hollow if all decisions are obvious. In a real bureaucracy, the typical experience is negative or neutral, punctuated by glimpses of satisfaction. 

\subsection*{Everyone do their own thing -- No Coordination, No Bureaucracy}
The minimal scenario to start from is to imagine a single person working on a single task that does not last long (a few minutes), is relatively easy (cognitively and physically and emotionally), and does not recur. In that situation, building consensus is irrelevant and no process is required. Even then, it is often the case that this simple task involves the use of shared limited resources -- basic essentials like water, air, land. If your task involves use of those things, then how is fair use determined among consumers?

Even for simplistic tasks, the concept of self-serve shared resources for each person is not trivial. independent allocation of shared resources breaks down and access is determined by violence.

Most of what you do occurs beyond the limits of simplicity and thus incurs some concept of \gls{process} (breaking a task into subtasks). Staying with the one-person constraint, a complex task can benefit from being broken into subtasks. Sometimes the order of the subtasks matters, so we need to track the dependencies. A recurring multi-step process with documentation is starting to have features of bureaucracy, but lacks the need for consensus. 


If one person lacks the skills relevant to a multi-step process, they may engage another person to help. The interaction may be informal (anarchy) or formalized in a contract (\href{https://en.wikipedia.org/wiki/Libertarianism}{libertarian}). If the parties working on the task fail to reach consensus, what is the recourse? Choices include physical violence, threats, or involving a third party (e.g., a court with lawyers and judges). 


To summarize, two aspects necessitate bureaucracy: shared resources and shared responsibility. 
If you want to manage resources and responsibilities, how is that policy decided?  Building consensus is relevant, but what is the process for establishing consensus? \href{https://en.wikipedia.org/wiki/Nepotism}{Nepotism} and cultural norms and \href{https://en.wikipedia.org/wiki/Religion}{religious practices} served this role prior to bureaucracy. 


\subsection*{Limit Bureaucracy to a Single Decider\label{sec:single-decider}}

Since bureaucracy is distributed knowledge and distributed decision making, it could be replaced by centralized knowledge and centralized decision making. If we are going to live in a society and coordinate shared resources, what if we had a single person deciding? How fast could that be? Good decisions can't be instantaneous -- the \href{https://en.wikipedia.org/wiki/OODA_loop}{OODA loop} is necessary. 
% https://graphthinking.blogspot.com/2016/03/how-to-evolve-organization-community-or.html
``Observe" is the input, ``act" is the output. These are measurable; ``orient" and ``decide" are not as easy to measure. The ``orient" is the process of labeling data (received in the ``observe" phase) and connecting that information to what is already known.

Three minutes is not a lot of time to gather information and share it to the relevant people, but let's set that as a lower bound.
If a single decision takes three minutes, then in a ten hour work day that's a max of (10*60)/3 = 200 decisions per day. If this decider works 300 days out of the year, that gets us to 60000 decisions per year. Everyone has to talk to this decider directly since there's no bureaucracy to gather or disseminate the information. The diversity of questions regarding shared resources would be challenging to answer well.

A bureaucracy of one person doesn't scale. Within this constraint we can't just automate everything because carrying out the automation requires a staff to implement. 
\index{automation!challenges of}
Automation isn't free -- it requires creation and maintenance. Unless the same person making the decisions is also the person creating and maintaining the system, there will be multiple people in an organization.


\href{https://en.wikipedia.org/wiki/Monarchy}{Monarchies} and \href{https://en.wikipedia.org/wiki/Dictator}{dictatorships} rely on a single decider. This is a simple model to understand, but addressing all the edge cases for large society can be difficult. A single decider doesn't scale for the number of decisions needed, so the decider then appoints bureaucrats to subjectively enforce policies. 

%\subsection*{Complexities of more than one Bureaucrat}

% https://graphthinking.blogspot.com/2019/07/first-principles-analysis-of-creating.html

If there's more than one person in an organization, then communication for coordination takes time. That's the ``orient'' phase of the OODA loop. Time spent orienting decreases the decision throughput of the organization.



\subsection*{Just use Common Sense}
\textit{Claim}: Everything would go smoothly if each person used common sense.
\marginpar{[Tag] Fallacy}

Common sense relies on reference experiences, cultural norms, incentives, emotional state, psychological defects. These aspects are particular to each person in an organization. 

In the same paradigm, suppose we start with an existing hierarchical bureaucracy and remove all the non-workers. That may sound like a worker's paradise, but then who would coordinates the activities with the number of people is above \href{https://en.wikipedia.org/wiki/Dunbar\%27s_number}{Dunbar's number}?

% https://graphthinking.blogspot.com/2021/09/why-is-everything-so-hard-in-large.html

\subsection*{Completely Avoid Bureaucracy}
That phrasing is more precisely worded by replacing ``bureaucracy" with ``coordination of stakeholders". If you avoid coordination of stakeholders, you either are constrained to only work on tasks that involve one person, or you get random (uncoordinated) interactions. 

\subsection*{Minimize Bureaucracy}
Again, try replacing ``bureaucracy" in with ``coordination of stakeholders". The goal of ``minimizing coordination" probably isn't the real objective. To be more precise, a specific objective might be ``minimize time spent executing the task" (which takes a lot of coordination prior to the task execution) or ``minimize the level of distraction to stakeholders" (chunk the coordination time, e.g. a meeting). Another strategy for minimizing bureaucracy is to reduce the number of stakeholders involved. For a given task complexity, this means having smarter people who have more skills. See figure~\ref{fig:complexity-and-size} for a quantitative trade-off. 


\begin{figure}
\includegraphics[width=0.8\textwidth]{images/people-per-task-for-skill-level.pdf}
\caption{A task that one smart person can do might take two dumb people. This concept applies to any task size and any population of workers. In this diagram three levels of task complexity are shown. As task complexity increases, the size of the team needs to grow. The growth may be less if the team members are brilliant. Those brilliant people cost more and there are fewer of them available.}
\label{fig:complexity-and-size}
\end{figure}


\subsection*{Automation of Processes\label{sec:automation}}

The role of automation is to make interactions more predictable, faster, and to handle more of them. Automation does not eliminate bureaucracy; automation obfuscates the subjective decisions and limits the ability to negotiate with decision makers.

% https://graphthinking.blogspot.com/2019/07/bureaucracy-is-social-process-executed.html
There are benefits to automation, and automation can be implemented to minimize harm to bureaucratic subjects.  Automated systems can be made more transparent by including documentation about what is happening, why, and what's next.
\index{automation!benefits of}
\marginpar{[Tag] Actionable Advice}
This helps both bureaucrats and subjects identify when assumptions made in the automation are incorrect. 


There are indicators of when to transition a human-based bureaucratic process to a more automated approach:
% https://graphthinking.blogspot.com/2019/07/cost-of-creating-and-supporting.html
\footnote{\href{https://xkcd.com/1319/}{https://xkcd.com/1319/} and \href{https://xkcd.com/1205/}{https://xkcd.com/1205/}}
\index{automation!when to}
\index{xkcd!\href{https://xkcd.com/1319/}{1319}}
\index{xkcd!\href{https://xkcd.com/1205/}{1205}}
\begin{itemize}
    \item The process to be automated is stable.
    \item The expected lifespan for the process -- how long the process will be needed -- is much longer than the time to implement the automation.
\item The process logic relies on objective evaluation criteria.  
\item The process is frequent.
\item The cost of automating (both the initial creation cost as well as the maintenance cost) is a savings over the manual implementation.
\end{itemize}


\subsection*{Market-based Approach}
% https://graphthinking.blogspot.com/2017/09/market-friction-and-bureaucratic.html

There is an alternative to bureaucracy that features decentralized approach to complex tasks and avoids reliance on consensus: \href{https://en.wikipedia.org/wiki/Market_(economics)}{markets}.

An over-simplified definition\footnote{\href{http://www.investopedia.com/terms/m/market.asp}{http://www.investopedia.com/terms/m/market.asp}} is ``A market is a medium that allows buyers and sellers of a specific good or service to interact to facilitate an exchange." 
In practice, there is market friction: ``trading is always associated with certain costs or restraints."
\footnote{\href{http://www.investopedia.com/terms/f/frictionlessmarket.asp}{Frictionless market};
and \href{https://insight.kellogg.northwestern.edu/article/adding_friction_to_the_market}{Adding friction to the market}.}

The sources of friction in a market include
\begin{itemize}
    \item commissions on trades
    \item taxes (needed to fund contract enforcement organizations)
    \item uncertainty in hiring
    \item constraints on firing labor
    \item search cost in exploring the available opportunities
    \item risk uncertainty
    \item cost of creating contracts
    \item cost of insurance
\end{itemize}
Market friction and bureaucratic inefficiency are similar. \footnote{  \href{http://www.economist.com/node/17730360}{Why do firms exist?} and \href{https://en.wikipedia.org/wiki/The_Nature_of_the_Firm}{The Nature of the Firm}.}

Complex tasks at a societal scale range from manufacturing advanced computer chips, to creating and distributing critical medicine, to immigration and border enforcement. 
If a society wants to carry out a complex task involving multiple people, then coordination of effort is required. The coordination can be explicit (in the form of a bureaucracy) or implicit (via a market).  Bureaucracy is one response to the complexity of a problem being solved.


\ \\

Regardless of whether a bureaucratic or market-based approach is used, 
% https://graphthinking.blogspot.com/2017/09/market-friction-and-bureaucratic.html
distributed knowledge and distributed decision making are hindered by aspects like
%\begin{itemize}
limited bandwidth between people participating in the coordination,
non-zero latency of information between people,
the cost of getting data,
and
the cost of analysis of data.
%\end{itemize}
Due to these factors, decisions that are suboptimal get made. See  %section~\ref{sec:dilemma-trilemma} 
the discussion of 
\hyperref[sec:dilemma-trilemma]{dilemmas}
for specific examples.

\ \\

The objective and quantified concept of money creates accountability that distinguishes commercial businesses from bureaucracy. Money as a metric is mutual to all participants within the business and money is in common with external stakeholders. 

Commercial businesses have employees who make subjective decisions and enforce policies, but unlike bureaucracy in business there is a common metric for feedback -- money. This distinction between business and bureaucracy is not sharp since the feedback mechanism is not perfectly applicable to all members of a business. Being an effective commercial bureaucrat may rely on the success of other employees rather than a direct interaction with customers.

Companies generally try to create profit. 
Bureaucracies like prisons, schools, hospitals, governments, and militaries all consume and spend money, but money isn't the goal. When faced with a decision, decisions are not guided by which will generate more profit~\cite{2012_Wilson}.

\subsection*{Adhocracy}
% https://graphthinking.blogspot.com/2017/09/complex-tasks-necessitate-bureaucracies.html

Suppose I have a team of diverse experts assembled to tackle a complex challenge.
\footnote{\href{https://en.wikipedia.org/wiki/Adhocracy}{Adhocracy}}
While this may be enough for short duration tasks, if the challenge lasts more than a day or two there will be new issues beyond the original challenge:

\begin{itemize}
    \item Who pays the salary for this labor?
    \item Who calculates the payroll?
    \item Who pays the rent for facilities used by the team?
    \item Who does the maintenance on the facilities and equipment used?
    \item Who cleans the facility?
    \item Does risk incurred mean insurance is needed?
\end{itemize}
To address these questions that are out-of-scope for the complex challenge, you can either use a market-based approach or build a bureaucracy. 

\subsection*{Consensus Algorithms}

In the description of a \hyperref[sec:single-decider]{single decider}
%section~\ref{sec:single-decider} 
the solution was centralization of decision making and knowledge. Another approach to distributed decision making is to use consensus algorithms like \href{https://en.wikipedia.org/wiki/Paxos_(computer_science)}{Paxos} and \href{https://en.wikipedia.org/wiki/Raft_(algorithm)}{Raft}\footnote{For a tutorial on Raft see \href{https://raft.github.io/}{https://raft.github.io/}}.


\subsection*{Elected Representatives}
Governments are composed of politicians and bureaucrats. (Government isn't the only place bureaucrats appear, but for this section we'll focus on government.)

Political representatives are an easier to understand concept because it's just one person acting in that behalf of other people, and people get a vote in who that representative is.
In contrast, the emergent behavior of bureaucracy is more difficult to understand: many people are involved, and subjects of bureaucracy do not appoint the bureaucrats. 

% why not make the entire system out of politicians?
Suppose that instead of a bureaucracy all members of an organization charged with managing shared resources were elected rather than being selected for their technical skills. In this thought experiment we are eliminating one of Weber's characteristics of bureaucracy -- competence for job appointments. 

In the United States there are bureaucratic positions that feature a mix of election versus appointment. For example, the method of selecting judges varies widely by state \footnote{\href{https://ballotpedia.org/Judicial_selection_in_the_states}{Ballotpedia's breakdown of Judicial Selection by State in the USA}.}. In a 2017 survey, 63\% of more than 1000 judges favored appointment over being elected. \footnote{\href{https://www.judges.org/news-and-info/the-age-old-question-should-judges-be-appointed-or-elected-heres-what-you-said/}{2017 survey} and \href{https://www.judges.org/wp-content/uploads/2020/03/Q1_Text.pdf}{comments}}
% fascinating set of survey results:
% https://www.judges.org/news-and-info/five-years-of-question-of-the-month/

Attorney Generals in the United States similarly a selected by election and appointment.
\footnote{\href{https://ballotpedia.org/Attorney_General_office_comparison}{Ballotpedia's breakdown of Attorney General offices}.} Positions that benefit from expertise (e.g., Attorney Generals, Coroner) sometimes lack minimum qualifications when selected by election. Having an informed electorate capable of selecting competing candidates is challenging the more positions there are to vote on.

\subsection*{Small Organizations}

Suppose you try to limit bureaucracy by imposing that teams or organizations be small.

Given three people in a team, options are
\begin{itemize}
    \item split the labor among the participants; each has the same workload and same tasks. Consensus decision making; no exploitation of unique skills.
    \item split the labor by specialization; each becomes dependent on the other. Specialization enabled by scale can improve effectiveness but also incurs lobbying which decreases effectiveness 
    \item make one the manager to oversee the other two -- hierarchy. This is a specific form of specialization where one person is not directly involved in the labor.
\end{itemize}


\subsection*{Defending against Malicious Participants}

Suppose that you are part of an organization that doesn't have oversight processes for finances and the resources your organization is responsible for. A small percentage of the population, say 0.1\%, will take advantage of the lack oversight for their own gain. Those malicious actors can be either subjects of bureaucracy or bureaucrats within the organization. \href{https://en.wikipedia.org/wiki/Bruce_Schneier}{Bruce Schneier} calls these people defectors~\cite{2012_Schneier}.

Similarly, suppose a small percentage of the population, say 0.1\%, is stupid and makes mistakes. Without a review process, these mistakes will negatively impact your organization's finances and the resources your organization manages. Preventing, detecting, and correcting mistakes can be costly investments. 

If you think bureaucracy sucks and it should be removed or replaced, then the parable of Chesterton's fence applies. 
\begin{quote}
There exists ... a fence or gate erected across a road. The more modern type of reformer goes gaily up to it and says, “I don’t see the use of this; let us clear it away.” To which the more intelligent type of reformer will do well to answer: “If you don’t see the use of it, I certainly won’t let you clear it away. Go away and think. Then, when you can come back and tell me that you do see the use of it, I may allow you to destroy it.\footnote{``The Thing'' (1929) Chesterton, in the chapter, ``The Drift from Domesticity''; \href{https://www.chesterton.org/taking-a-fence-down/}{https://www.chesterton.org/taking-a-fence-down/}}
\end{quote}
You need to understand the historical progression of each process that your organization relies on so that you understand why bureaucracy exists.

 % section
  \clearpage
  \subsection{Alternative views of bureaucracy within a bureaucratic organization\label{sec:alternative_views_from_within}}

Bureaucracy as I have defined it in \S\ref{sec:define_bureaucracy} is not the only way that bureaucrats perceive their environment. 

The perspectives below are archetypical; an individual's perspective might be a mixture of these views.

\ \\

\textbf{As a bureaucrat, what matters is what I can accomplish with my skills and the resources I have access to.} \\
\textit{Assessment}: This person is task oriented. Results are what matters. The intricacies of bureaucracy are a distraction to getting the work done. The emotional reward for this person is accomplishment of the task. This person is likely to say to their manager, "Tell me what I need to do to be successful" rather than identify collaborations.

\ \\

\textbf{As a bureaucrat, what matters is how I feel.} \\
\textit{Assessment}: Your feelings are real. They have consequence, in that your emotions impact motivation and enthusiasm. However, a feelings-centric perspective may not be productive for you or your team or the organization. Balancing those competing needs is challenging.

\ \\ 

\textbf{What matters is how others feel.}\\
Depending on the emotional state of those around you is unhealthy and can be unproductive. Working for the happiness or satisfaction of other people is risky -- they may not know what's best, or they may not have your interests in mind.

\ \\

\textbf{What matters is my immediate coworkers.}\\
In this scenario everything else is personified or ignored. This can be positive (I collaborate with those around me) or negative (I am in competition with those around me). \\
\textit{Assessment}: Your do relationships matter. However, they are not all that matters. Being able to explain what is happening outside the immediately observable realm is what is missing from this view. 

 % section
  \clearpage
  \section{Models of Bureaucracy that are Incomplete\label{sec:models-of-bureaucracy}}

Coming up with a holistic theory of bureaucracy is desirable but difficult. Bureaucracy exists in every society, so having an explanatory theory would help identify what aspects are essential and what is accidental. Having a theory of bureaucracy could help identify what should be improved and what should be discarded. An expectation for the existence of a theory stems from the repeated independent creation of bureaucracy in diverse societies in different time periods. 

Characterizing bureaucracy is difficult because organizations comprised of unique humans are aware of attempts to be characterized and respond to stories told about bureaucracy; see the \href{https://en.wikipedia.org/wiki/Hawthorne_effect}{Hawthorne effect}. 
\index{Wikipedia!\href{https://en.wikipedia.org/wiki/Hawthorne_effect}{Hawthorne effect}}
\marginpar{[Wikipedia] Hawthorne\\effect}
As soon as a claim about aspects that characterize bureaucracy is made, then an individual can respond to that claim by behaving in an opposing manner. Worse still for the theory, bureaucrats can coordinate amongst themselves to provide counterexamples. 

In the following section I outline a few conventional ways of modeling bureaucracy to point out the shortcomings of each model. The relevance to the practicing bureaucrat of familiarity with these models is so that you know the boundaries of each model. Awareness of the limitations of each model enables you to know when the model is an applicable story and when the model is not explanatory. 

\subsection*{Bureaucracy as a Machine}

\textit{Narrative}: Bureaucracy is a machine that has throughput and latency and dependencies and mechanisms. Bureaucrats are cogs in that machine.\\
\textit{Why this model feels true}: Bureaucracy can be complicated and feel mechanical with standardization, top-down dictates, and interlocking processes. There is a risk of \href{https://en.wikipedia.org/wiki/Deindividuation}{deindividuation} 
\index{Wikipedia!\href{https://en.wikipedia.org/wiki/Deindividuation}{deindividuation}}
\marginpar{[Wikipedia] deindividuation}
for bureaucrats. \\
\textit{What this model is missing}: Perceiving yourself as a cog in the machine implies a loss of agency. This is a \href{https://en.wikipedia.org/wiki/Self-fulfilling_prophecy}{self-fulfilling prophecy}
\index{Wikipedia!\href{https://en.wikipedia.org/wiki/Self-fulfilling_prophecy}{self-fulfilling prophecy}}
\marginpar{[Wikipedia] self-\\fulfilling prophecy}
-- if you think you don't have agency, then you may stop acting as though you have agency. 


\subsection*{Bureaucracy as an Economic model}

\textit{Narrative 1}: Bureaucracy is a collection of individual rational actors. \\
% https://graphthinking.blogspot.com/2019/05/cooperation-and-competition.html
\textit{Why this model feels true}: Individual bureaucrats cooperate or compete to promote their own self-interests.
The culture of an organization is a result of individual self-interests of each bureaucrat.
Individuals make decisions that promote their own self-interests. \\
\textit{What this model is missing}: This view is hard to distinguish from a marketplace. 

\ \\
\textit{Narrative 2}: Bureaucracy is comprised of competing special-interest groups. \\
\textit{Why this model feels true}: The observation of \href{https://en.wikipedia.org/wiki/Public_choice}{Public Choice theory} 
\index{Wikipedia!\href{https://en.wikipedia.org/wiki/Public_choice}{Public Choice theory}}
\marginpar{[Wikipedia] Public\\Choice theory}
is that a concentrated minority that stands to gain a disproportionate benefit will act in their own self-interest. \\
\textit{What this model is missing}: In a generic bureaucratic organization it is not clear which team would be more concentrated than any other team, nor is it clear what the disproportionate benefit might be. There is a concentrated minority at the top of any hierarchy, and the top of the hierarchy does act in their own self-interest, though this is not particular to bureaucracy.

Bureaucratic organizations may have specific teams with access to disproportionate benefits, but this book focuses on generic bureaucratic features. Effective bureaucrats are all alike; every bureaucrat is ineffective in their own way
\footnote{A modified version of the first sentence of Leo Tolstoy's novel ``Anna Karenina.''}.
% Anna Karenina is also referenced on page 142 of "Making of a Manager." 
% I came up with my adaptation before reading "Making of a Manager"



\ \\
\textit{Narrative 3}: A Bureaucracy is a subcategory of a Firm. \\
% https://graphthinking.blogspot.com/2017/09/market-friction-and-bureaucratic.html
Firms exist in a market because negotiating contracts and prices for every interaction is burdensome. 
% https://www.kellogg.northwestern.edu/faculty/hubbard/htm/research/ec174/lectures/3coase.htm
A bureaucracy could be considered as a type of firm that specializes at the organizational level in policy administration or resource management. \\
\textit{Assessment}: This is correct.

%Doesn't address small vs large companies, and doesn't distinguish between profit-oriented and non-profit and government. 


\subsection*{Bureaucracy as Emergent phenomenon}

\textit{Narrative}: There is a universality to bureaucracy in both the diversity of scenarios and persistence across time that hints at \href{https://en.wikipedia.org/wiki/Emergence}{emergence}.\\
\index{Wikipedia!\href{https://en.wikipedia.org/wiki/Emergence}{emergence}}
\marginpar{[Wikipedia] emergence}
\textit{Why this model feels true}: Bureaucracy as a macroscopic phenomenon is emergent when there are a sufficient number of people involved. The size of the organization is important because there is no longer dependence on individual relationships (i.e., a size above \href{https://en.wikipedia.org/wiki/Dunbar\%27s_number}{Dunbar's number}). 
\index{Wikipedia!\href{https://en.wikipedia.org/wiki/Dunbar\%27s_number}{Dunbar's number}}
\marginpar{[Wikipedia] Dunbar's\\numbers}
There are people in the organization that you don't know and therefore there is a lack of personal accountability. An organization is subdivided into teams recursively until there is local person-to-person accountability.  The underlying behaviors that enable emergence are bilateral interactions among humans and a lack of feedback mechanisms. 

At the scale of individual bureaucrats, every person is playing by different rules and has different goals and everything is changing -- both the individuals and the conditions. 
Above the threshold for emergence of bureaucracy there is scale-free behavior. The same patterns are observable at large organizations and extremely large organizations. A bureaucrat in one organization recognizes patterns of professional life experienced by bureaucrats at another organization. The local mechanisms bureaucrats employ to enable distributed decisions using distributed knowledge include meetings, processes, and communications. While local nuances differ, a generic pattern is apparent. 

The choices faced by an individual bureaucrat are interdependent with the choices made by other bureaucrats in their environment. There is a \href{https://en.wikipedia.org/wiki/Flocking_(behavior)}{flocking behavior} 
\index{Wikipedia!\href{https://en.wikipedia.org/wiki/Flocking_(behavior)}{flocking behavior}}
\marginpar{[Wikipedia] flocking\\behavior}
where my choices are informed by the choices of those around me. Unlike flocking of birds, the adjacency metric for bureaucrats is not necessarily spatial distance. Instead, visibility of the decisions and consequences inform adjacency.


The relevance of claiming bureaucracy is emergent is that there is behavior occurring at the macroscopic scale. Knowing  motives and actions of every individual bureaucrat at the microscopic level is not relevant. Treating organizations as complex and adaptive systems gives insight on how to work within the dynamic environment~\cite{2011_Eisenhardt}.


A colloquial interpretation of emergent behavior from complex phenomena is treating the system as an entity -- personification.
\marginpar{[Tag] Fallacy} 
When an organization is assigned behaviors~\cite{2002_Gall}, it is useful to remember that the organization is comprised of individual bureaucrats. 


% https://www.preposterousuniverse.com/podcast/2021/10/11/168-anil-seth-on-emergence-information-and-consciousness/
More formally, there are distinct categories of emergence~\cite{2002_Bedau, 2021_Carroll_168}. Nominal emergence names the phenomena of a thing being distinct from its constituents: a circle is emergent from a collection of points; a pile of sand emerges from grains of sand. Merely putting many bureaucrats in a room is insufficient to create bureaucracy; there's more to bureaucracy. 

Weak emergence occurs when there are phenomena that are independent of the underlying interactions. For example, \href{https://en.wikipedia.org/wiki/Glider_(Conway\%27s_Life)}{gliders} 
\index{Wikipedia!\href{https://en.wikipedia.org/wiki/Glider_(Conway\%27s_Life)}{gliders}}
in \href{https://en.wikipedia.org/wiki/Conway\%27s_Game_of_Life}{Conway's Game of Life}.
\index{Wikipedia!\href{https://en.wikipedia.org/wiki/Conway\%27s_Game_of_Life}{Conway's Game of Life}}
Weak emergence is measurable using \href{https://en.wikipedia.org/wiki/Granger_causality}{Granger causality} 
\index{Wikipedia!\href{https://en.wikipedia.org/wiki/Granger_causality}{Granger causality}}
or, equivalently\footnote{\href{https://arxiv.org/abs/0910.4514}{https://arxiv.org/abs/0910.4514}}, \href{https://en.wikipedia.org/wiki/Transfer_entropy}{transfer entropy} 
\index{Wikipedia!\href{https://en.wikipedia.org/wiki/Transfer_entropy}{transfer entropy}}
(information theory). I don't know how to apply these measures to bureaucratic organizations. 

\textit{What this model is missing}: The problem with treating an organization as an entity is that apparent behavior is counter-intuitive~\cite{2002_Gall}. Breaking the organization into individual people with motives helps clarify causes of observed behavior. 
%This merely improve the post hoc rationalization. 

In practice, bureaucracy is worse than emergent because the rules can be altered or ignored by  stakeholders. Bureaucracy is a \href{https://en.wikipedia.org/wiki/Wicked_problem}{wicked problem}~\cite{1973_Rittel} 
\index{Wikipedia!\href{https://en.wikipedia.org/wiki/Wicked_problem}{wicked problem}}
\marginpar{[Wikipedia] wicked\\problem}
which resists mathematical models. 

\subsection*{Bureaucracy in \href{https://en.wikipedia.org/wiki/Game_theory}{Game Theory}}
\index{Wikipedia!\href{https://en.wikipedia.org/wiki/Game_theory}{Game Theory}}
\textit{Narrative}: Bureaucracy is comprised of one or more \href{https://en.wikipedia.org/wiki/List_of_games_in_game_theory}{games} 
\index{Wikipedia!\href{https://en.wikipedia.org/wiki/List_of_games_in_game_theory}{list of games in game theory}}
played by bureaucrats. Which game is applicable depends on the specifics of the situation. \\
\textit{Why this model feels true}: Different bureaucrats have different motives in distinct situations. \\
\textit{What this model is missing}: While interactions among bureaucrats may be describable in terms of games, that doesn't provide an underlying motive for bureaucracy as an identifiable experience.  

Defining bureaucracy as distributed knowledge and distributed decision-making for the subjective management of access to shared resources does not fit as a \href{https://en.wikipedia.org/wiki/Coordination_game}{coordination game} 
\index{Wikipedia!\href{https://en.wikipedia.org/wiki/Coordination_game}{coordination game}}
or \href{https://en.wikipedia.org/wiki/Non-cooperative_game_theory}{competitive} game. 
\index{Wikipedia!\href{https://en.wikipedia.org/wiki/Non-cooperative_game_theory}{non-cooperative game theory}}
%Bureaucracy is self-modifying. 
Bureaucracy is in constant flux due to external conditions, externally imposed constraints, staff turn-over, internal dilemmas, and disagreements among individuals. 
% https://en.wikipedia.org/wiki/Evolutionary_game_theory




\subsection*{Bureaucracy as Evolutionary Outcome}


\textit{Narrative 1}: Biological, Genetic evolution -- Individual level. \\
There might be biological arguments for a genetic basis for bureaucracy. 
\textit{Why this model feels true}: For example, Zebra fish and Hyenas have a gene for dominance\footnote{``Society, demography and genetic structure in the spotted hyena'' (2012); doi:10.1111/j.1365-294X.2011.05240.x}. \\
%Breeding animals for aggressiveness?
\textit{What this model is missing}: ``Genes code for proteins, so there are no `genes for' phenotypes per se, including behavioral phenotypes."~\cite{2015_Lilienfeld}

Non-human animals make subjective decisions and do not get labeled as bureaucrats. Human make decisions that are not explained by reproductive fitness.

\ \\
\textit{Narrative 2}: Biological, Genetic evolution -- Group selection. \\
\textit{Why this model feels true}: \hyperref[sec:hierarchy-of-roles]{Hierarchy}
%(section~\ref{sec:hierarchy-of-roles}) 
is not unique to humans. Primates form into social hierarchies based on dominance over a shared resources like mates and food. \\
\textit{What this model is missing}: a gene for bureaucracy. A genetic model does not explain what behaviors an individual bureaucrat can take to be more effective. 

\ \\
\textit{Narrative 3}: \href{https://en.wikipedia.org/wiki/Memetics}{Memetic}
\index{Wikipedia!\href{https://en.wikipedia.org/wiki/Memetics}{Memetic}}
-- bureaucracy as method for coordination is a better idea than \href{https://en.wikipedia.org/wiki/Nepotism}{nepotism} 
\index{Wikipedia!\href{https://en.wikipedia.org/wiki/Nepotism}{nepotism}}
or religion. \\
\textit{What this model is missing}: Bureaucracy persists along with nepotism and religion, and bureaucracy occurs within religion. 

\subsection*{Bureaucracy as Product-focused Narrative}
\textit{Narrative}: Ignore the bureaucrats and instead focus on how a product progresses through a process.\\
\textit{Why this feels true}: Ignoring bureaucrats involved in processes simplifies the narrative (the product is the main character). 
\textit{Example}: \href{https://www.youtube.com/watch?v=OgVKvqTItto}{School House Rock: I'm Just A Bill}\\
\textit{What this model is missing}: Decision-makers involved the process exert subjective control. Outcomes depend on who participates. 

\subsection*{Bureaucracy as Subject-focused Narrative}
\textit{Narrative}: Ignore the bureaucrats and instead focus on the person subjected to bureaucracy. \\
\textit{Why this feels true}: The person experiencing bureaucracy as a subject is confused by ``why isn't this easier?''  \\
\textit{What this model is missing}: History of why the process exists and how it evolved (i.e., legacy), protection against malicious subjects, and protection against malicious bureaucrats. 


\subsection*{Bureaucracy as Psychological Phenomenon}

\textit{Narrative}: Bureaucracy as a pure power struggle. Or the interplay of individual pathologies. 
%Are you doing what's best for you, the group you're in, or everyone?
%Altruistic or reciprocal? Retaliation.
%The answer changes time the time and situation of situation and person to person.
%\href{https://en.wikipedia.org/wiki/Spiral_of_silence}{Spiral of silence}
% and
%\href{https://en.wikipedia.org/wiki/Social_proof}{social proof}

Organizations are composed of individuals with personalities, and inefficiency is attributable to a mashing together of distinct individuals with conflicting desires.
While this is true, it isn't complete. \\
\textit{What this model is missing}: Personality-focused narrative neglects the history of processes (legacy), protection against malicious subjects, and protection against malicious bureaucrats. Analysis that stops at personalities misses emergent phenomena. %Analysis that stops at personalities of individuals typically explains larger scale phenomena using personification of teams and organizations. 

\ \\

Each of the above models of bureaucracy has shortcomings. Knowing why each model is not a complete description helps you avoid the trap of thinking only in terms of a single model. 

The rest of this book ignores these partial characterizations of bureaucracy. 
Rather than take the external (\href{https://en.wikipedia.org/wiki/Emic_and_etic}{emic}) 
\index{Wikipedia!\href{https://en.wikipedia.org/wiki/Emic_and_etic}{emic and etic}}
view of bureaucracy, this book takes the internal (etic) perspective of a bureaucrat operating within an organization. 
The description of bureaucracy in this book is used to contextualize advice for how to be an effective bureaucrat.  % section
  \clearpage 
  \clearpage
  \section{Not Many Options for Organizations}
% https://graphthinking.blogspot.com/2021/07/patterns-anti-patterns-in-bureaucracy.html

The purpose of an organization or a team is to obfuscate the bureaucracy of enacting policies for managing access to shared resources. The value of organizations and teams is that they are a simplifying assumption that depersonalize the function of inputs and outputs. Without organizations or teams you would be directly exposed to the complexity of knowing every person's title and responsibility. Organizations and teams are abstractions to hide complexity.

When complaining about the ineptitude of organizations (or more specifically, leaders, managers, bureaucrats you manage, and coworkers), consider the variables available to be modified. As a thought exercise, what would it take to rebuild the organization you are in from scratch? 
\marginpar{$>>$ Thought Exercise}

Listing the levers available for enacting change in an organization is intended to emphasize the importance of personal interactions within bureaucracy. When you get frustrated with inefficiency, consider the systemic changes available -- there aren't that many. Options for restructuring the organization are limited, so the importance of your creativity and informal influence is dominant in determining the effectiveness of bureaucracy. Those factors relevant to Process Empathy are separate from having the right staff and being able to have an adequate number of staff. 

Organizations comprised of bureaucrats have fewer options for change than individual bureaucrats. The choices an individual bureaucrat faces are described
\iftoggle{haspagenumbers}{on page~\pageref{sec:dilemma-trilemma}}{ }%
 in the section on 
\hyperref[sec:dilemma-trilemma]{dilemmas}.
In comparison, the choices faced by the designers of an organization include:
\begin{itemize}
    \item Flatness of organizational hierarchy -- how many layers of oversight are there? Another way of arriving at the same structure is to ask how many employees there are per supervisor.
    \item Hierarchy split by function (all the lawyers in one team, all the engineers in one team, all the sales people in one team) or by product (mixing experts to solve a customer's problem)?

% https://graphthinking.blogspot.com/2023/07/matrixed-organizations-best-practices.html
    
    \item Size of organization. Having more people allows for more specialization and requires more support. 
    
    \item Organizations create processes for recurring tasks like hiring, promotion (pay or title), awards, compensation, recognition, professional training, and firing. Each of those has a set of design choices that inform the organization's culture.
\end{itemize}
In a government bureaucracy, the constraints of some incentives like pay and financial awards are set outside the organization. That further constrains the ability of bureaucrats to shape their organization's culture.

The organization may have policies and processes regarding hiring, promotion, training, and firing, but the decision may be made by team managers rather than top-level managers. 

%Listing how an organization could change is relevant 
Individual bureaucrats can enact change, and subjects of bureaucracy can request improvement, but those are weak inputs. 
Organizations are accountable to their source of funding, not the subjects of bureaucracy or members of the organization. For example, the \href{https://en.wikipedia.org/wiki/Internal_Revenue_Service}{Internal Revenue Service (IRS)} 
\index{exemplar!Internal Revenue Service@\href{https://en.wikipedia.org/wiki/Internal_Revenue_Service}{Internal Revenue Service}}
\index{Wikipedia!Internal Revenue Service@\href{https://en.wikipedia.org/wiki/Internal_Revenue_Service}{Internal Revenue Service}}%
is accountable to Congress, not taxpayers or IRS employees. 



A flat organization (more employees per manager) means more people doing work instead of managing, but sufficiently large organizations are typically broken into teams to make coordination less chaotic. 
%TODO:~\cite{2015_Katzenbach} (Wisdom of Teams book).


Organizations are segmented into teams and there are not many options: create a new team, merge existing teams, or dissolve a team. 


Shuffling the structure of teams (referred to as re-organization) is a favorite hobby of executive leadership. The alignment of hierarchical structure with an organization's purpose is a change that is easy to point to as an accomplishment. 


For a given set of teams, the lateral interactions are competitive or cooperative. Coordination is required (or conflict will occur) for money, staffing, and resources. Examples of resources include access to or control of data, computer equipment, hardware, floor space in a building, prestige, and ownership of products.

% ==== hierarchy - expand/contract ====
% Ways to expand and contract hierarchy as the number of employees in the workforce grows and shrinks
% * Add a new organization below: person who was an employee is now a manager and hires a bunch of people
% * Add a new peer team,
% * Add a new parent oversight organization above the current one
% * Unflatten an organization by inserting a new layer of management


Not all changes have to be structural; there are choices to be made when operating within the existing hierarchy.


\subsection*{Not Many Options Within Teams}



Similar to the limited number of options available for shaping organizations, the structural choices faced by a team are constrained. 


The purpose of pointing out the variables that can be changed is to show that there aren't many. As a consequence of the limited options, relationships and effective communication are more significant for effective bureaucracy. 


Team managers might decide who gets hired, who gets promoted, who goes to what training, and who gets fired. Though in some environments even that control is relegated to an external team. A team manager usually has decision-making authority regarding tasks the team works on. 

% intra- and inter- team dynamics

Accountability in the context of teams comes from person-to-person interactions. These can be either lateral (sideways), parent-child (top-down), or child-parent (bottom-up)~\cite{2014_Jorgensen}. Each of these three categories has associated constraints.

The upward child-parent (bottom-up) communication is either inadequate (too few updates, not enough information, or insufficient context), relevant, or excessive. For example, a weekly or monthly report to multiple superiors may be inadequate. 

Finding the balance depends on the presenter knowing the individual audience members so that a tailored message is provided and then adapting to the specifics of the situation. 
% tips on managing up: 
% https://svpg.com/managing-up/

The downward parent-child communication either is inadequate (no direction provided or imprecise direction provided), provides actionable vision, or micromanagement. Finding the right balance and specificity requires insight into the communication needs on both sides of the relationship. The manager of a team member has a variety of tools available beyond setting tasks and deadlines. A manager can create an environment that provides psychological safety, promotion based on accomplishments, recognition of contributions, encouragement, and constructive feedback.

For interactions among team members at the same hierarchical level the conflict of interests between cooperation and competition manifests in struggles over money, staffing, products (output), and resources (inputs). Here ``resources''  refers to constraints like access to data, control of data, technology resources, hardware, floor space, and expertise. 


% https://graphthinking.blogspot.com/2021/07/patterns-anti-patterns-in-bureaucracy.html

%\begin{itemize}
%    \item money
%    \item staffing
%    \item prestige
%    \item products (output)
%    \item resources (inputs)
%    \begin{itemize}
%        \item access to or control of data
%        \item technology resources
%        \item hardware
%        \item floor space
%        \item expertise
%    \end{itemize}
%\end{itemize}


Understanding the options available for organizations, teams, and relationships informs your Process Empathy. There are constraints bureaucrats operate within and there are a limited number of ways to evolve out of a situation. By listing possible ways to change you can identify which options are best to invest in. 
  %\section{Folk Wisdom}
While most of the entries on
 https://github.com/dwmkerr/hacker-laws
don't apply to bureaucracy, the list is useful to review. 

Folk wisdom is an attempt to explain bureaucratic features but ends up being \gls{thought terminating}.

\href{https://en.wikipedia.org/wiki/Murphy\%27s_law}{Murphy's law}

\href{https://en.wikipedia.org/wiki/Wiio\%27s_laws}{Wiio's laws}

\href{https://en.wikipedia.org/wiki/Hanlon\%27s_razor}{Hanlon's razor}

\href{https://en.wikipedia.org/wiki/Parkinson\%27s_law}{Parkinson's law}

\href{https://en.wikipedia.org/wiki/Putt\%27s_Law_and_the_Successful_Technocrat}{Putt's Law}

\href{https://en.wikipedia.org/wiki/Hick\%27s_law}{Hick's law}

\href{https://en.wikipedia.org/wiki/Allen_curve}{Allen curve}

\href{https://en.wikipedia.org/wiki/Peter_principle}{Peter principle}

\href{https://en.wikipedia.org/wiki/Dilbert_principle}{Dilbert principle}

\href{https://en.wikipedia.org/wiki/Law_of_triviality}{Law of Triviality}\clearpage
  \section{Bureaucracy as Accidental, Legacy, or Essential}
\gls{essential bureaucracy} is the minimum processes and staffing and skills necessary to address the complexity of the managing a community's access to a \gls{shared resource}. Achieving this minimum is tricky since the optimization can be with respect to resilience to change, resilience to edge cases, staff turn-over, speed experienced by consumer, financial cost, time spent by the organization, and number of staff. Miss any one of those and the bureaucracy is deemed inefficient.

Undesirable bureaucracy is categorized as either accidental or legacy. Accidental bureaucracy arises when someone misunderstands what is needed, or when skills of the bureaucrats involved are insufficient for the complexity a problem. Legacy bureaucracy occurs when the situation changes but the processes do not. Resolving each of these suboptimal conditions may seem easy: have better knowledge of the problem, assign the right people to the problem, and change processes as problems evolve. 

Having enough knowledge is often infeasible, especially for complex problems at large scale. Having the people with the right skills assumes that a pipeline of people with relevant talents exists and that people in the pipeline won't be poached to work on other challenges. Keeping up with evolving problems depends on having resources to change (beyond the maintenance baseline), and having a defined approach for changing the process. 
  \clearpage
  \section{Bureaucratic Fallacies\label{sec:fallacies}}
\index{bureaucratic fallacy}

Discussions about bureaucracy by non-experts often rely on common conceptions that are \gls{thought-terminating}. Identifying these enables you to understand both why the fallacy is attractive during a discussion with fellow bureaucrats and how each idea is incomplete.

These fallacies may at first feel right but are in fact misleading. In contrast, there are  
\hyperref[sec:unavoidable-hazards]{unavoidable hazards} \iftoggle{haspagenumbers}{(see page~\pageref{sec:unavoidable-hazards})}{} that may feel bad but reveal underlying truths.

\ \\
\begin{samepage}
\textit{Bureaucratic fallacy}: \textbf{Bureaucracy is bad}. \\
\index{bureaucratic fallacy!bureaucracy is bad}
\textit{Why this feels true}: When a person subjected to bureaucracy has a negative experience, the easiest attribution is to the least-understood aspect -- the bureaucracy.\\
\textit{What this is missing}: 
\iftoggle{glossarysubstitutionworks}{\Gls{bureaucracy}}{Bureaucracy}
 as defined in this guide is neither good nor bad. Bureaucracy is merely a way of managing  resources shared amongst a community. 
 \end{samepage}

\ \\
\begin{samepage}
\textit{Bureaucratic fallacy}: 
\textbf{There is no point in planning ahead since everything (staffing, funding, purpose, scope) is always changing.}\\
\index{bureaucratic fallacy!no point in planning}
\textit{Why this feels true}: Change can feel disorienting, especially when it is unexpected. A change of the assumptions for a plan may make the plan less relevant. \\
\textit{What this is missing}: Preparing for change and thinking ahead about contingencies enables effective use of resources. Have a vision and work towards it while accounting for and adapting to change. This approach requires extra work, some of which will be left unused. See Dilemma \ref{table:dilemma-personal-emergencies-vs-ignore}\iftoggle{haspagenumbers}{ on page~\pageref{table:dilemma-personal-emergencies-vs-ignore}.}{.}
\end{samepage}

\ \\
\begin{samepage}
\textit{Bureaucratic fallacy}: \textbf{Bureaucracy is an aberration, a mistake, due to poor planning or incompetent participants}. \\
\index{bureaucratic fallacy!bureaucracy is a mistake}
\textit{Why this feels true}: Mistakes are made, poor or insufficient planning does happen, and some participants are incompetent.\\
\textit{What this is missing}: Bureaucracy occurs even if no mistakes are made, effort is spent on effective planning, and participants are competent. That's because, in the use of distributed knowledge and distributed decision-making, bureaucrats face \hyperref[sec:dilemma-trilemma]{dilemmas}.
\marginpar{See page~\pageref{sec:dilemma-trilemma}.}
%\ifsectionref
%(see section~\ref{sec:dilemma-trilemma}).
%\fi
\end{samepage}

\ \\
\begin{samepage}
\textit{Bureaucratic fallacy}: \textbf{Bureaucracy is inefficient}. \\
\index{bureaucratic fallacy!bureaucracy is inefficient}
\textit{Why this feels true}: Expressed by both subjects and bureaucrats who observe seemingly wasteful processes.\\
\textit{What this is missing}: If bureaucracy were truly inefficient (not allocating resources in the most efficient way), then in a competitive environment it would be replaced by a more efficient approach. The key is to ask, ``efficient with respect to what metric?'' The metric of money, or time, or number of people, or stability, or robustness to perturbation?  Second, what would motivate improved efficiency? Without incentives, change is less likely. 
\end{samepage}

\ \\
\begin{samepage}
\textit{Bureaucratic fallacy}: \textbf{Bureaucracy is due to malfeasance.}\\
\index{bureaucratic fallacy!bureaucracy is due to malfeasance}
The specific number of malicious bureaucrats in variations on this fallacy ranges from ``all of the participants'' to ``just enough to be problematic.'' \\
\textit{Why this feels true}: There are bad actors present in any system comprised of humans. \\
\textit{What this is missing}: Most participants are earnestly trying to help make a positive contribution, even though that can be hard to see from the view of subjects or even other bureaucrats. Processes like isolation or promotion exist within bureaucracy to deal with malicious bureaucrats.
\end{samepage}

\ \\
\begin{samepage}
\textit{Bureaucratic fallacy}: \textbf{Bureaucracy is a sign of decay from within the organization.} \\
\index{bureaucratic fallacy!bureaucracy indicates decay}
\textit{Why this feels true}: Relationships within an organization have a half-life and require ongoing investment to renew. At the same time, new bureaucratic processes are constantly being developed by other bureaucrats. The number of processes increases as the organization ages. Bureaucracy seems to arise without effort and countering it takes effort.  \\
\textit{What this is missing}: Bureaucracy unavoidably emerges in every organization because coordination is required. The negative connotation of decay should be replaced with a sense of neutral evolution.
\end{samepage}

\ \\
\begin{samepage}
\textit{Bureaucratic Fallacy}: \textbf{If the response to a request I make can't be expedited, my request must not be important}.  \\
\index{bureaucratic fallacy!importance is measured by response latency}
\textit{Why this feels true}: Other people would demonstrate they care about what I am working on by prioritizing things I am dependent on.\\
\textit{What this is missing}: When everything gets prioritized, that's the same as nothing getting priority.
\end{samepage}

\ \\
\begin{samepage}
\textit{Bureaucratic Fallacy}: \textbf{The expected duration of a task is how long it would take one person to accomplish}.  \\
\index{bureaucratic fallacy!task duration for one person}
\textit{Why this feels true}: When I imagine carrying out a task, the default is a story with one character. \\
\textit{What this is missing}: This narrative fails to account for the overhead of interaction among participants and delays due to asynchronous dependencies.~\cite{1975_brooks}
\end{samepage}

\ \\
% https://graphthinking.blogspot.com/2019/08/two-misleading-simplifications-when.html
\begin{samepage}
\textit{Bureaucratic Fallacy}: \textbf{When developing or altering policy, focus on the average or majority (to the exclusion of outliers)}. \\
\index{bureaucratic fallacy!ignore outliers}
\textit{Why this simplification is misleading}: Sometimes outliers are not just more of the same; they alter the outcome. During the transition from horses-for-transportation to cars, cars could initially have been considered outliers. 
\end{samepage}

\ \\
\begin{samepage}
\textit{Bureaucratic Fallacy}: \textbf{People learn from their mistakes}. \\
\index{bureaucratic fallacy!people learn from their mistakes}
\textit{Why this feels true}: There's an optimistic desire for this to be true. \\
\textit{What this is missing}: 
People repeat mistakes without noticing. People do not naturally reflect on their failings in a constructive way and then apply insights to future situations.
People \textit{can} learn from their mistakes. Doing so requires a low latency feedback loop and incentive to change. In bureaucracies feedback loops are weak so learning may not happen.
\end{samepage}

\ \\
\begin{samepage}
\textit{Bureaucratic Fallacy}: \textbf{Processes are serial}.\\
\index{bureaucratic fallacy!processes are serial}
A conventional approach to process design is a sequence of tasks. As an example, consider approval chains. \\
\textit{Why this feels true}: Serial processes are easier to understand. \\
\textit{What this is missing}: Some tasks that are independent can be carried out concurrently; see the section on \hyperref[sec:reducing-overhead]{reducing overhead}\iftoggle{haspagenumbers}{on page~\pageref{sec:reducing-overhead}.}{.}
\end{samepage}

\ \\
\begin{samepage}
\textit{Bureaucratic Fallacy}: \textbf{Hard work creates results}.\\
\index{bureaucratic fallacy!hard work creates results}
\textit{Why this feels true}: Some results do require hard work. The alternative (results arise serendipitously or with little effort) is not inspiring. \\
\textit{What this is missing}: Hard work can be invested on wasteful effort. Don't confuse being busy with being productive. Sometimes insight is more useful than hard work. 
\end{samepage}

%\ \\

% NOT USEFUL
%\textit{Bureaucratic Fallacy}: \textbf{Motivations for bureaucrats are categorized as individualistic, tribal, organizational, societal, or humanity}.\\



\ \\
\begin{samepage}
\textit{Bureaucratic Fallacy}: 
\textbf{You cannot pay a little and get a lot}; see the \href{https://en.wikipedia.org/wiki/Common_law_of_business_balance}{Common law of business balance}. 
\index{Wikipedia!\href{https://en.wikipedia.org/wiki/Common_law_of_business_balance}{Common law of business balance}}
\marginpar{$>>$ Folk Wisdom} 
\index{folk wisdom!\href{https://en.wikipedia.org/wiki/Common_law_of_business_balance}{Common law of business balance}} \\
\index{bureaucratic fallacy!cannot pay a little and get a lot}
\textit{Why this feels true}: If small investments made a big difference we'd already be making the investment.\\
\textit{What this is missing}: This doesn't account for creative solutions and ignores \href{https://en.wikipedia.org/wiki/Nudge_theory}{nudge theory}
\index{Wikipedia!\href{https://en.wikipedia.org/wiki/Nudge_theory}{nudge theory}}
from behavioral economics. 
\end{samepage}

\ \\

When you are reasoning about bureaucratic systems, there may be a conclusion that is concise and feels explanatory. Then you should try to come up with counter-examples, either logically or from experience.  

Some fallacies are based on an expectation that other people should be more like what you imagine your best self to be. That mindset fails to account for your  shortcomings and the diversity of other bureaucrats. 

 
  \clearpage
  \section{How to Reduce Administrative Overhead\label{sec:reducing-overhead}}

You can read this list from a position of authority to take action, or you can read this list as an advocate for improvement. 

% https://graphthinking.blogspot.com/2018/08/how-to-decrease-bureacracy.html
In isolation, none of these ideas should be surprising.
\begin{itemize}
    \item \textit{Tip}: Automate recurring decision processes. This decreases bureaucracy by removing subjective influence of bureaucrats
\item \textit{Tip}: Where automation is infeasible, make customer advocates with end-to-end authority available. 
This decreases bureaucracy by improving customer's navigation of processes.
\item \textit{Tip}: Make processes transparent to participants. 
This decreases bureaucracy by improving customer's understanding of processes.
\item \textit{Tip}: Make information discoverable (e.g., via search engine) .
\item \textit{Tip}: Make information directly available, rather than mediated by a person.
\item \textit{Tip}: After an interaction is completed, summarize the steps and outcome for the participant. 
\item \textit{Tip}: When a process fails the needs of a participant, investigate the failure and improve the process\footnote{See the Wikipedia enty on \href{https://en.wikipedia.org/wiki/Continual_improvement_process}{Continual improvement process}}. 
\item \textit{Tip}: Make the goals and priorities of the organization clear to all.
\item \textit{Tip}: Define measurable standards of performance, both for individuals and teams.
\item \textit{Tip}: Train bureaucrats how to engage participants effectively; these interactions determine the culture.
\item \textit{Tip}: Train bureaucrats by addressing their immediate problems (e.g., through mentorship).
\item \textit{Tip}: Make employment desirable to people who have desirable characteristics (e.g., educated candidates).
\item \textit{Tip}: Enhance accountability to peers (e.g., peer review of actions and outcomes).
\item \textit{Tip}: Make sure that incentives for organizations and individuals encourage improvement rather than maintenance of status quo.
\item \textit{Tip}: Decision making should be pushed down the hierarchy to the practitioner.
\item \textit{Tip}: When decision making requires cross-organization interaction, form a team of practitioners.
\item \textit{Tip}: Bosses should share workload with their team to gain practical exposure to current challenges.
\item \textit{Tip}: Seek feedback from process participants; then provide status updates on the implementation process.

% https://graphthinking.blogspot.com/2018/08/sequential-approval-chains-are.html
\item \textit{Tip}: Replace sequential approval chains with a concurrent review process. Sequential approval reduces the number of ideas reaching later approvers by filtering good ideas as well as bad ideas. Approvers may perform independent (redundant) due diligence and end up disagreeing. 
\end{itemize}

  \clearpage
  \section{An Effective Bureaucrat\label{sec:effective-bureaucrat}}

You can leverage your improved understanding of bureaucracy by changing your behavior. In this section I provide actionable suggestions. This section provides prescriptions, whereas other sections like the  
\marginpar{See page~\pageref{sec:dilemma-trilemma}.}%
\hyperref[sec:dilemma-trilemma]{dilemmas} 
and 
\hyperref[sec:unavoidable-hazards]{unavoidable hazards} are descriptive. If you only have the prescriptions (and skip the descriptions) you might get confused and frustrated as to why other bureaucrats are not applying best practices. 
%CANTDO\marginpar{See page~\pageref{sec:unavoidable-hazards}.} 


My baseline assumptions are that you are a good person, you have the relevant technical skills for your duties, and you are a good \href{https://en.wikipedia.org/wiki/Project_management}{project manager}.
\index{Wikipedia!project management@\href{https://en.wikipedia.org/wiki/Project_management}{project management}}\iftoggle{WPinmargin}{\marginpar{$>$Wikipedia: project management}}{}
Enacting Process Empathy is in addition to those elements. 

\subsection*{Characteristics of an Effective Bureaucrat}

Bureaucracy is a system of distributed knowledge and distributed decision-making, so coordination is critical. Here are characteristics to strive for. 
\begin{itemize}
    \item You communicate more effectively than anyone around you. This applies to verbal and written communication, \hyperref[sec:effective-presentations]{formal presentations} 
    \marginpar{See page~\pageref{sec:effective-presentations}.}% 
    and \hyperref[sec:walk-arounds]{informal interactions}. \\
%CONFLICTING\marginpar{See page~\pageref{sec:walk-arounds}} \\
    \textit{Why this doesn't happen by default}: people plateau when there is no incentive to improve. What level of quality is good enough depends on the situation-specific risks and costs.
    \item You facilitate meetings so that productivity is noticeably improved. This means creating and sharing an agenda, providing rules on interaction (e.g., raising hands), taking notes, sharing notes, and following on with assigned tasks after the meeting.\\
    \textit{Why this doesn't happen by default}: facilitation is rarely taught explicitly. The work of facilitation is considered secondary to the discussion, even when the productivity of discussions is disastrous and demoralizing. 
    \item You take part in meetings, whether that means actively contributing or intentionally supporting other attendees. You leverage relationships with other attendees.\\ 
    \textit{Why this doesn't happen by default}: active participation requires effort and exposes you to more risk than remaining silent.
    \item You invest effort in written communication (emails, text-based chats, reports). Your writing empathizes with readers, captures relevant context, is concise, and is clearly worded.\\
    \textit{Why this doesn't happen by default}: writing-as-a-skill is rarely the focus of a bureaucrat's work. Writing is seen as an entry-level skill that doesn't require improvement. Finding examples of clear writing in bureaucratic organizations is challenging, so you may struggle to find role models. 
    \item You communicate verbally concisely and precisely. You listen and you teach. You seek shared definitions. \\
    \textit{Why this doesn't happen by default}: Without a feedback mechanism motivating urgency, bureaucrats tend to pontificate. 
    \item If you are in-person with colleagues, you walk around and talk with people one-on-one.  \\
    \textit{Why this doesn't happen by default}: bureaucrats interested in ``just doing the work'' view professional interactions as a distraction that decreases productivity. 
    \item You are 
\marginpar{See page~\pageref{sec:professional-vulnerability}.}%
    \hyperref[sec:professional-vulnerability]{professionally vulnerable}. \\
    \textit{Why this doesn't happen by default}: Sharing stories of your personal experiences in the organization requires coherent narratives that the listener can learn from. That talent is not taught to all employees in most organizations. 
    \item In your role as bureaucrat you leverage project management skills: you have a vision, you make and share plans, all while building consensus with stakeholders. \\
    \textit{Why this doesn't happen by default}: Consensus requires knowing who is relevant to include and then investing in relationships that allow iterative feedback. That use of time is costly. The level of extroversion may not be attractive to every bureaucrat.
    \item You apply your negotiation skills~\cite{1982_Cohen} that improve your interactions and outcomes.\\
    \textit{Why this doesn't happen by default}: bureaucrats are not taught negotiation unless it is an explicitly named responsibility for their role. 
    \item You reply quickly to incoming requests, whether that means answering directly or acknowledging the request and providing a timeline for when you will answer. This allows other people's tasks to either be resolved or have a clear response about when progress can be expected. \\
    \textit{Why this doesn't happen by default}: The coping skills of responding to a large number of inbound requests either grows through experience, is learned by watching role models, or just doesn't happen.
    \item You strive for and demonstrate transparency. You share information with stakeholders. Transparency  enables coordination without interpersonal relationships.\\
    \textit{Why this doesn't happen by default}: Transparency endangers incumbents. Transparency requires an ongoing investment that doesn't contribute to addressing the primary issue. Transparency requires empathy with subjects.
\end{itemize}

\ \\

The superpowers of a bureaucrat that facilitate cooperation and progress in any organization include the following.
\begin{itemize}
\item You seek information from stakeholders without burdening them. When your requests are burdensome, you acknowledge that and seek ways to pay back (or forward) their investment.
\item You apply consistent processes (rather than being reactionary and applying ad hoc responses).
\item You hold others (and yourself) accountable for their actions. Accountability is created by clearly stating objectives and then measuring results.
\item You adapt to the varying incentives and reference experiences of those around you. Flexibility enables interactions with diverse coworkers. 
    \item You effectively multi-task or, more accurately, switch tasks. The switch among tasks is triggered when the current task encounters an externally-generated pause. You can \href{https://en.wikipedia.org/wiki/Pipeline_(computing)\%23Concept_and_motivation}{work on a task while waiting for another task to finish}. 
    \index{Wikipedia!pipeline computing@\href{https://en.wikipedia.org/wiki/Pipeline_(computing)\%23Concept_and_motivation}{pipeline (computing)}}
\iftoggle{WPinmargin}{\marginpar{$>$Wikipedia: pipeline (computing)}}{}
    \item You are prepared with a backlog of ideas (in writing) if someone shows up with resources.
    % https://graphthinking.blogspot.com/2016/12/life-lessons-i-learned-from-experience.html
    \item You know when to change communication channels (from text chats to phone calls to in-person). 
    \item If you are dependent on someone else getting something done to enable your progress, you can demonstrate priority by physically showing up -- \underline{presence creates priority}. 
    \index{mantra!presence creates priority}
    Being physically at a person's desk motivates that person to respond better than calling  or emailing them. Showing up where someone works and talking with them conveys how much priority you place on the actions of the person you're talking with.
    \item You have intellectual empathy -- theory of mind for thinking (whereas empathy refers to emotions). How to grow your intellectual empathy: shadow peers and bosses and coworkers and subordinates.
    \item You have \gls{process empathy}. You recognize the deviations and exceptions that cause processes to come into existence. 
    \item You avoid unintentionally neglecting tasks or requests. This requires capturing incoming requests and then providing a response about prioritization and status. This matters because other participants in the organization should regard you as reliable in a positive sense. 
    \item You focus on value delivery in relationships to a degree that exceeds the scope of your formal role.
\item You have altered your job description to fit the growth you're seeking.
\item You are willing to engage on a personal level and know stakeholders outside their professional role.
\item Each of your tasks has a customer, a deadline, and a deliverable artifact. You iterate towards a result. 
\end{itemize}

You occasionally ponder and discuss with other people introspective questions like
\begin{itemize}
    \item How can I be successful?
    \item What are the ways I could fail?
    \item How would the organization be characterized as successful?
    \item What are the ways the organization can fail?
\end{itemize}
% https://graphthinking.blogspot.com/2017/07/questions-to-ask-mentors.html
You find mentors and ask them questions like
\begin{itemize}
    \item What do you like most about your career? 
    \item Given a chance, what would you do differently?
    \item How do you manage work/life balance?
    \item What's the big challenge for our industry in the next two years?
    \item How would you tackle that challenge?
    \item What advice do you have for a young person starting in this industry? In this organization? 
    \item Are there mistakes to avoid?  
    \item How can I be successful?
    \item What books would you recommend reading?
\end{itemize}

Although you strive to enact Process Empathy, do not expect that of other people. 

\subsection*{Frames used by Effective Bureaucrats}

Being an effective bureaucrat is a mindset. I refer to these perspectives as framing. 

\ \\
\textit{Framing}: A bureaucrat can do more as part of an organization than by working alone. Being a member of an organization means the bureaucrat's identity is subsumed into service for the organization they are part of.\footnote{See Wikipedia entry on \href{https://en.wikipedia.org/wiki/Deindividuation}{Deindividuation} -- the loss of self-awareness in groups.
\index{Wikipedia!deindividuation@\string\href{https://en.wikipedia.org/wiki/Deindividuation}{deindividuation}}
%%%CANTDOINFOOTNOTE\marginpar{$>$Wikipedia: deindividuation}
} At the same time, bureaucracy enables the bureaucrat to amplify their presence by being part of a larger organization.  Sometimes the cost of being part of the organization exceeds the force multiplier of working together. 

\ \\
% https://graphthinking.blogspot.com/2019/05/definition-of-progress.html
\textit{Framing}: Measuring your personal growth in a bureaucracy is difficult due to the lack of \hyperref[sec:feedback-loop-and-ripples]{feedback loops}. One approach is to measure your capabilities for a specific task. If you can complete the task in less time, with fewer resources, and with less effort, that's progress. If you can now complete a task that you previously wanted to but weren't able to, that's progress.


\ \\
\textit{Feelings}: Bureaucracy induces an emotional response in participants because things don't work the way each person wants. This can lead emotionally to frustration and then apathy. Understanding how things work in a bureaucracy can help decrease the anger.

% Wandering the maze of bureaucratic processes as a subject.

Another emotional response to bureaucracy is a sense of powerlessness. 
\begin{quote}
``Some third person decides your fate: this is the whole essence of bureaucracy~\cite{1921_Kollontai}."
\end{quote}
That sense of powerlessness applies both to bureaucrats and to subjects of bureaucracy. 

The sense of powerlessness is somewhat valid, in that you are as a bureaucrat giving up some power compared to your ability to act individually. That is the trade for working with other people.

\ \\
\textit{Feelings}: A process feels bureaucratic when there seems to be dissonance. The subject faces a multi-step task that they can imagine one person doing the same action in fewer steps and taking less time. The illusion of decreased bureaucracy is created by consolidation from the subject's perspective. The barrier to enacting consolidation is the necessary coordination among different stakeholders who receive no benefit from the consolidation. Externalizing the coordination to the subject is what causes the sense of bureaucracy. 

\subsection*{How to be an Effective Bureaucrat}

The following advice is specific to situations that you may not have yet encountered. 

\ \\
\textit{Do}: \textbf{Work on three tasks concurrently.} Don't rely on one person or one idea for your success. On the other end of the spectrum, don't spread yourself too thin.
\marginpar{$>>$ \href{https://en.wikipedia.org/wiki/Goldilocks_principle}{Goldilocks balance}}%
\index{Goldilocks balance!number of tasks}%
Working on a multitude of projects decreases risk of any one outcome failing but also decreases the amount of attention spent thinking about a specific challenge. 

The three tasks should be at off-set stages (early, mid-way, and nearing completion) rather than all being the same level of maturity.

\ \\
\textit{Do}: When you get stuck, use the following changes of perspective:  
\textbf{Look upstream, look downstream. Look to peers. Zoom in (narrower scope) and zoom out (broader scope)}. These are all ways of changing the context. That change may help you identify assumptions that are holding back progress.

\ \\
\textit{Do}: When you are asked to take on more work, \textbf{avoid responding with ``That's not my job."} If the request is misguided and your perception is that another person has the responsibility, ask if the requester is aware of that other person's responsibilities. If you are to take on the work, get guidance on re-prioritizing and ensure the request is documented in writing. 

\ \\
\textit{Do}: If there's something you want to do, \textbf{strive for influence without authority} instead of working to gain control over resources (e.g., through promotion). Avoid the following: ``In this organization X is important to me, but I can't do X right now because I don't have enough power in the organization. So I'll get promoted and then do X."

\ \\
\textit{Do}: \textbf{Learn the perspectives of those around you.}\\
The relevance of knowing the paradoxes including dilemmas and unavoidable hazards, is that you should talk explicitly to your fellow bureaucrats about these specific topics in conversation. Not that the goal is to find consensus or agreement. But to find what the other person is thinking so that you can account for their processes

\ \\
\textit{Do}: \textbf{Account for holistic view.}\\
The specific circumstances of the challenges you face as a bureaucrat depend on the individual people involved, what the purpose of the bureaucracy is, what technology is available for enacting bureaucracy, and the resources (staffing, money, time). 

\ \\
\textit{Do}: \textbf{Learn the history of the challenge.}\\
This goes beyond
\href{https://en.wikipedia.org/wiki/G._K._Chesterton\%23Chesterton's_fence}{Chesterton's fence}\iftoggle{haspagenumbers}{ (see page~\pageref{concept:chestertons_fence}), }{,}%
\index{Wikipedia!Chesterton's fence@\href{https://en.wikipedia.org/wiki/G._K._Chesterton\%23Chesterton\%27s_fence}{Chesterton's fence}}%
\iftoggle{WPinmargin}{\marginpar{$>$Wikipedia: Chesterton's fence}}{}%
which focuses on why the current approach is in place. Learning the history of a problem means what has been tried before and failed. How did the previous iterations evolve into the current situation? Was the cause personalities, insufficient resources, inadequate technology? What's changed that enables this approach to be better? What do you know that prior attempts didn't?

\ \\
\textit{Do}: \textbf{For a given challenge, work on three remedies}.\\
One solution you're working on may fail, and another may be inadequate. 

\ \\
\textit{Do}: \textbf{Minimize \href{https://en.wikipedia.org/wiki/Externality}{imposing costs on other people}}.\\
\index{Wikipedia!externality@\href{https://en.wikipedia.org/wiki/Externality}{externality}}\iftoggle{WPinmargin}{\marginpar{$>$Wikipedia: externality}}{}
A solution that externalizes costs harms the greater organization and creates bureaucratic debt.

\ \\
\textit{Do}: \textbf{Exploit the flexibility of rules for the benefit of all parties.}\\
% https://graphthinking.blogspot.com/2019/07/winning-game.html
change the rules of the game and the objectives of the game such that every participant wins.

\ \\
\textit{Do}: \textbf{Align selfish interests with social interests.}\\
See also \href{https://en.wikipedia.org/wiki/Nudge_theory}{nudges}
\index{Wikipedia!Nudge theory@\href{https://en.wikipedia.org/wiki/Nudge_theory}{Nudge theory}}\iftoggle{WPinmargin}{\marginpar{$>$Wikipedia: Nudge theory}}{}
(page 66 of Schneier's Liars and Outliers~\cite{2012_Schneier}).

\ \\
\textit{Do}: \textbf{Find ways to rephrase negative complaints}.\\
% https://graphthinking.blogspot.com/2019/07/how-to-rephrase-negative-observations.html
Negative observation: ``Logging into my computer takes a long time."\\
Positive statement and explanation of impact: ``If the latency for logging into my computer were lower, I could make more progress on X."


Negative observation: ``The service team I need support from doesn't offer a ticketing queue."\\
Positive statement: ``If the service team I need support from offered a ticketing queue, I would be able to track the work done on my behalf."

\ \\
\textit{Do}: \textbf{Share lessons learned}.\\
This can build your network. People value vulnerability in others; you can make the first move.  
  \clearpage

\chapter{Why Bureaucracy is Hard}
  \section{Subject's View of Bureaucracy\label{sec:subjects-view}}

This section takes on the perspective of the subject of bureaucracy but is meant to be read by bureaucrats who want to improve their \gls{process empathy}. This book doesn't provide advice for subjects of bureaucracy.\footnote{For advice navigating bureaucracy, listen to National Public Radio's \href{https://www.npr.org/2022/03/16/1086915600/get-what-you-want-customer-service}{How to talk with customer service}~\cite{2022_LifeKit}.} 

\ \\

When you have a positive experience engaging with bureaucracy, your positive attribution is to the people involved. Or the ease of a solution makes bureaucracy less visible and the solution seems obvious. 
When you have a negative experience with bureaucracy, complaints are about the incompetence of the people involved or the incomprehensibleness of the system. Don't these bureaucrats know how to do their job? Why isn't the solution obvious? Why does this system not work for me? David Graeber summarized this view:\footnote{\href{https://harpers.org/archive/2015/03/in-regulation-nation/}{In Regulation Nation. Harpers Magazine, 2015.}}
% May have actually come from "The Utopia of Rules"
\begin{quote}
 Amongst working-class Americans, government is generally seen as being made up of two sorts of people: `politicians,' who are blustering crooks and liars but can at least occasionally be voted out of office, and `bureaucrats,' who are condescending elitists almost impossible to uproot.   
\end{quote}


%\subsection{Sources of complexity}

The scale of bureaucracy (the number of people in an organization) and the processes of an organization can seem disproportionate to the complexity of the task. Typically when a person (being a subject of bureaucracy) interacts with the bureaucratic organization the artifacts are simple, like a form to fill out. The simplicity of the artifact does not correlate to the number of decisions made, the tracking of information, or the precautions taken by the organization. All these aspects are invisible to the subject because they are internal to the organization.

As an example, consider when you go to the doctor and they mend your broken arm with a cast. That seems straightforward because all you see is the doctor putting a cast on. You don't get insight into the decisions they had to make. Why did they need 20 years of focused schooling to carry out a procedure that took 15 minutes?

Judging bureaucracy by the artifact visible to the individual subject undersells the complexity of the decision-making necessary to take action. The contingencies that you were not exposed to because everything went well make the amount of investment from the bureaucrat appear wasteful.

As the subject of bureaucracy, you also lack the ability to distinguish how much work is attributable to the bureaucrat generating justifications for their actions (colloquially, \href{https://en.wikipedia.org/wiki/Cover_your_ass}{covering the bureaucrat's ass}). 
\index{Wikipedia!\href{https://en.wikipedia.org/wiki/Cover_your_ass}{cover your ass}}
\iftoggle{WPinmargin}{\marginpar{[Wikipedia] Cover\\your ass}}{}
These justifications are needed both within the organization and potentially for external stakeholders. Each bureaucrat's rationalization may not be reviewed, but it needs to be available for review later.

As the subject of bureaucracy, you typically can't distinguish when work is caused by a bureaucrat's ill-informed decisions. Is the person stupid, mistaken, or is there something you are not taking into account?
Sometimes the work is carried out by insufficiently trained bureaucrats, but you don't get to know whether you're working with an experienced and knowledgeable bureaucrat or a new untrained bureaucrat. 

As the subject of bureaucracy, you don't have visibility on the many nuances of an organization. The internal power struggles and organizational politics that depend on personalities, resources, and competing prioritization are not clear to outsiders.
The inconsistency of an organization's policies may not be felt by bureaucrats in that organization. Each bureaucrat may have a different opinion, resulting in a lack of consistent guidance.
There may be legacy policies in effect.

% summary paragraph
Effective action by a bureaucratic organization is complicated by the need for relevant information. Gathering, analyzing, and sharing that information persistently requires bureaucratic processes. To further complicate the ideal process, sometimes the organization lacks the staffing with relevant skills. 


% Transition paragraph
The next section documents why bureaucracy is hard from the perspective of the bureaucrat. Even without getting into the specifics of a bureaucrat's role or the purpose of an organization, there are generic reasons that bureaucracy is a burdensome responsibility. % section
  \clearpage
  \section{Bureaucrat's experience of Dilemmas\label{sec:dilemma_trilemma}}

% Not included here: 
% https://en.wikipedia.org/wiki/The_Innovator%27s_Dilemma
% because it is the etic view of change


When a decision has two viable options (neither being best), that presents a \href{https://en.wikipedia.org/wiki/Dilemma}{dilemma}. The name for a decision with three viable options is a \href{https://en.wikipedia.org/wiki/Trilemma}{trilemma}. This section describes the bureaucrat's experience of operating within an organization in terms of dilemmas and trilemmas. 
% https://en.wikipedia.org/wiki/Defeasible_reasoning

%Dilemma are not unique to bureaucracy. 
Dilemmas are a \href{https://en.wikipedia.org/wiki/Defeasible_reasoning}{simple way} of discussing decisions but are not the only relevant aspect of bureaucracy.
Dilemmas are merely a useful framing to highlight the following ideas:
\begin{itemize}
    \item You, an individual bureaucrat, face decisions that you may not have recognized. Failing to recognize a choice can lead to suboptimal results. The dilemmas below are generic to any bureaucratic process and are intended to stimulate your ability to identify decisions. 
    \item You face complex trade-offs in your role as a bureaucrat. Dilemmas are intended an entry point to more nuanced reflection that is specific to your situation. The dilemmas below are not exhaustive; this list is merely illustrative. 
    \item You can build \href{https://en.wikipedia.org/wiki/Theory_of_mind}{intellectual empathy} with fellow bureaucrats when you recognize they face the same dilemmas. These dilemmas are generic to the situation and the person facing the decisions. You now have a topic to discuss with them.
    \item You can be curious about the choice other bureaucrats select for a given dilemma. Everyone in the organization faces these decisions, so these dilemmas give you a topic of conversation to better understand each bureaucrat's view.
    \item You can identify and negotiate potential sources of friction. Other bureaucrats may arrive at different selections for a given dilemma, so recognizing this and discussing it can improve the effectiveness of all involved.
\end{itemize}


% claim
Dilemmas explain the inherent complexity of bureaucracy even when bureaucrats are honest and the purpose of the organization is clear.
% relevance
The point isn't that one should \href{https://en.wikipedia.org/wiki/False_dilemma}{select one of the two options}. The point of recognizing a dilemma is that it is a marker that there is a decision to be made.
% consequence
The action for the bureaucrat in response to this assessment is to identify nuances, enumerate alternatives, and talk with fellow bureaucrats about these decisions. Rather than seeking consensus, strive for comprehension of other people's perspective. That way you can navigate processes more effectively.




\subsection{How Dilemmas Arise in Bureaucracy}

Given a task, the simplest response for an individual is to take action and be done. 

If the task is more challenging, it may be necessary to first make a plan, then take action and be done. The source of the challenge may be task complexity, scale, number of people impacted, diversity of stakeholders, amount of time needed, how the current task shapes future decisions, number of collaborators, etc. Regardless of why the task is challenging, a dilemma has already arisen: how much time to spend planning versus doing. 

If the task is even more challenging, it may be useful to gather data for the plan, then make a plan, then take action and be done. A new dilemma arises: how much time to spend gathering data versus planning. The previous dilemma still exists -- planning versus doing. As an alternative to the dilemma framing, how much time should be allocated to three categories of activity?

But wait, how did the person in the first scenario (just do the task) know no plan was necessary? Either it was an explicit choice or they didn't perceive the need to decide about whether to plan. Someone else faced with that same task might choose differently, i.e., to make a plan. The task complexity is relative to the person's skills, experiences, expectations about stakeholders, and potential ramifications.

In this escalating sequence of increasingly challenging tasks, suppose the task involves you and another person. The overhead of coordination inflicts additional dilemmas. The examples of complexity arising from coordination are identified in the list of dilemmas below.

An independent source of friction is when people involved in the task don't agree on how challenging it is. The framing of the difficulty matters because it can lead different participants to distinct conclusions about how much time should be spent planning versus for carrying out the action. The choice of ``how complicated is the task?'' shapes the team dynamics and informs the need for hierarchical roles. 

\subsection{Folk Wisdom on Decision Making}

Bureaucrats, whether hired for their expertise or simply to provide labor, are rarely experts on decision making. There are multiple domains in which \href{https://en.wikipedia.org/wiki/Decision_theory}{decision making} is studied (e.g., \href{https://en.wikipedia.org/wiki/Rational_choice_theory}{economics}, \href{https://en.wikipedia.org/wiki/Game_theory}{mathematics}, \href{https://en.wikipedia.org/wiki/Decision-making}{psychology}), but practicing bureaucrats are more likely familiar with colloquialisms that feel descriptive. A few are provided here to give a sense of both the conciseness and the lack of action embedded in each meme.

\ \\
\href{https://en.wikipedia.org/wiki/Hick\%27s_law}{Hick's law}\marginpar{[Tag] Folk wisdom}: ``Increasing the number of choices will increase the decision time logarithmically.''

\ \\
\href{https://en.wikipedia.org/wiki/Hanlon\%27s_razor}{Hanlon's razor}\marginpar{[Tag] Folk wisdom}: ``Never attribute to malice that which is adequately explained by stupidity.''

\ \\
\href{https://en.wikipedia.org/wiki/Parkinson\%27s_law}{Parkinson's law}\marginpar{[Tag] Folk wisdom}: ``Work expands so as to fill the time available for its completion.''

\ \\
\href{https://en.wikipedia.org/wiki/Murphy\%27s_law}{Murphy's law}\marginpar{[Tag] Folk wisdom}: ``Anything that can go wrong will go wrong.''

\ \\
\href{https://en.wikipedia.org/wiki/Law_of_triviality}{Law of Triviality}\marginpar{[Tag] Folk wisdom}: ``People within an organization commonly or typically give disproportionate weight to trivial issues.''

\ \\

In comparison to the folk expressions above, dilemmas are intended to be accessible to practicing bureaucrats and trigger reflection and discussion. Effective action is enabled by improved understanding of the trade-offs faced within an organization. 

\subsection{Dilemmas are a poor framing}

Although the following concepts are presented as dilemmas and trilemmas, these are necessarily simplifications. For example, two ways to simplify situations into a dilemma are
\begin{itemize}
    \item Start with a single variable, e.g., ``how much data gathering", and force it into a binary choice ``more data gathering" versus ``less data gathering."
    
    % a reduction of a complicated situation to one variable. 
    \item Start with a complex trade-off space with many opportunities and reduce it to a \href{https://en.wikipedia.org/wiki/False_dilemma}{false dichotomy} of \href{https://en.wikipedia.org/wiki/Zero-sum_thinking}{zero-sum options}: ``more data gathering" vs ``more planning." The actual trade-space involves optimization of multiple objectives, like (maximize productivity) and (minimize risk) and (maximize quality) and (maximize employee satisfaction) and (minimize latency). 
    % https://bennorthrop.com/Essays/2022/code-ownership-stewardship-or-free-for-all.php
\end{itemize}
These over-simplifications neglect both the continuous nature of the trade-offs and the alternative creative approaches to a specific situation. 

%\subsection{Dilemmas as a Framing to Ease you into Complexity}

The choices described below in the simplified representation of dilemmas are intended as a starting point for introducing the decision space relevant to individual bureaucrats. Do not accept the dilemmas presented here as the final framing. Thinking in terms of a limited spectrum of opportunities neglects nuances that enable more creative approaches. 
Awareness of these dilemmas and trilemmas are intended to spark creative imagination about the nuances specific to your situation.

By recognizing the deficiencies of dilemmas, you can identify nuances specific to the situation you are in. Instead of responding to the recognition of a decision with ``should I do this or that?'' the better option is to assess the complexities and \href{https://en.wikipedia.org/wiki/Brainstorming}{brainstorm} multiple options. Talk with stakeholders and understand the history of the situation before making a choice.


Ponder theses dilemmas prior to the pressure of real-time decision making.  Recognize dilemmas and trilemmas and then avoid them by adapting to the local conditions and specific people available to help.

\subsection{Dilemmas are not Purely Intellectual}
Decision making is not a purely intellectual task; there is emotional stress induced by the process. Dilemmas create cognitive dissonance for the decision maker. Any selection is going to have downsides, and any compromise will be suboptimal. Those burdens weigh on deciders.

To counter this morale weight, talk with other people about decisions. Even if this does not alleviate the responsibility of deciding, discussion can help you arrive at new insights. 

\subsection{Additional Complications}
Adding to the difficulty and stress, dilemmas presented here are occur concurrently and continuously. The dilemmas are inter-dependent due to both to the common variables and the constrained resources.
Selecting an option for one dilemma alters the options available in other dilemmas.

Oscillation between approaches can be caused by change of management, accumulation of experience (dissatisfaction) with one solution, the desire for promotion (change as progress), or a desire for cost savings (efficiency). The rate of oscillation is an indicator of the half-life of \href{https://en.wikipedia.org/wiki/Institutional_memory}{institutional memory for an organization}.  

\ \\

The following two sections categorize dilemmas as personal policies (\S\ref{sec:personal_policy_dilemmas}) and policies regarding structure of the organization (\S\ref{sec:org_dilemma}). The personal policies apply to each bureaucrat in an organization, while the structural policies for an organization are faced by a subset of bureaucrats in the management role. 

\subsection{Personal Policy Dilemmas \label{sec:personal_policy_dilemmas}}

In practice, the following decisions are unordered and are constantly faced by the bureaucrat. As observed by Lindblom in \cite{1959_Lindblom}, this flurry of decisions contrasts to a regularized process that might be envisioned as optimal.


  
\begin{center}
\begin{table}[H] % ht
\begin{tabular}{ | m{\dilemmatablewidth}| m{\dilemmatablewidth} | } 
  \hline
  \textbf{Intervene before the deployment of a policy or process or product, perhaps lacking relevant context.} &
  \textbf{Wait with feedback until deployment.} \\
  \hline
  \textit{Cons}: Engaging prematurely betrays your awareness; future explorations by that team are made less visible. & 
  \textit{Cons}: The team wasted time and attention on something that wouldn't work or may even be harmful. \\
  \hline
\end{tabular}
\caption{As an outsider to a team responsible for a process/policy/product, suppose you learn of something prior to official deployment (i.e., you learn the internal musings of another team). This Dilemma of Early Intervention and is not unique to bureaucracies. There is folk wisdom on both sides: ``Stay in your own lane'' and ``Speak up when you see something wrong.'' See the related Dilemma of micromanaging, \ref{table:micromanaging}.}
\label{table:early-intervention}
\end{table}
\end{center}


\begin{center}
\begin{table}[H] % ht
\begin{tabular}{ | m{\dilemmatablewidth}| m{\dilemmatablewidth} | } 
  \hline
  \textbf{Review status of the work of other people early and often. Many milestones, check-ins, and updates.} &
  \textbf{Review status infrequently; just do the work.} \\
  \hline
  \textit{Description}: Micromanagement. & 
  \textit{Description}: Hands-off management style. \\
  \hline
  \textit{Cons}: Takes up your time and the people you're reviewing. & 
  \textit{Cons}: Team members are unsure how to proceed and don't know what the goal is. \\
  \hline
\end{tabular}
\caption{The dilemma of micromanaging peers and subordinates is not unique to bureaucratic organizations. The right blend of how much engagement by reviewers depends on the personalities and desires of each person.}
\label{table:micromanaging}
\end{table}
\end{center}


\begin{center}
\begin{table}[H] % ht
\begin{tabular}{ | m{\dilemmatablewidth}| m{\dilemmatablewidth} | } 
  \hline
  \textbf{Bureaucrat expects management to provide solutions -- just tell me members what to do.} & 
  \textbf{Bureaucrat dislikes managers micromanaging by telling people what to do.} \\ 
  \hline
  \textit{Cons}: Your manager may not have insight on what needs to be done. Or they may guide you in a less effective direction. &
  \textit{Cons}: No autonomy, unable to exploit your expertise and creativity. \\  
  \hline
\end{tabular}
\caption{This is the complement of \ref{table:micromanaging}. Nominally the manager helps identify the objectives and provides context and the subordinate figures out how to accomplish the objective, but who is responsible for what is negotiable in each relationship.
}
\label{table:solution_provider}
\end{table}
\end{center}


\begin{center}
\begin{table}[H] % ht
\begin{tabular}{ | m{\dilemmatablewidth}| m{\dilemmatablewidth} | } 
  \hline
  \textbf{Write everything down to \href{https://en.wikipedia.org/wiki/Cover_your_ass}{cover your ass}.} &
  \textbf{Don't record sensitive conversations that could be used against you or others.} \\
  \hline
  \textit{Cons}: Takes a lot of time and effort to accurately capture intent. Recording can be done poorly or be misconstrued.  & 
  \textit{Cons}: No written record to point to when someone changes their behavior. \\
  \hline
\end{tabular}
\caption{I personally write things down and share them with other people (e.g., this book), but there are costs and risks to investing in documentation. There are \href{https://en.wikipedia.org/wiki/Dark_pattern}{dark patterns} for this trade-off, like intentionally misquoting another person to bias the documentation in your favor, or only writing down the aspects of conversation that favor the outcome you are interested in.}
\label{table:notes_or_no_notes}
\end{table}
\end{center}


\begin{center}
\begin{table}[H] % ht
\begin{tabular}{ | m{\dilemmatablewidth}| m{\dilemmatablewidth} | } 
  \hline
  \textbf{Ponder what should or could be done.} &
  \textbf{Figure out how to accomplish the objective.}\\
  \hline
  \textit{Cons}: Less time for action. & 
  \textit{Cons}: Prematurely select an action that is suboptimal. \\
  \hline
\end{tabular}
\caption{Brainstorming is useful, as it considering the holistic situation. At some point that transitions to action, but when? This is a question of how much time to spend admiring the forest versus the trees. 
}
\label{table:forest-vs-trees}
\end{table}
\end{center}




\begin{center}
\begin{table}[H] % ht
\begin{tabular}{ | m{\dilemmatablewidth}| m{\dilemmatablewidth} | } 
  \hline
  \textbf{Allocate time for meetings to facilitate coordination.} &
  \textbf{Allocate time for action.} \\
  \hline
  \textit{Cons}: Less time for participants to implement ideas. & 
  \textit{Cons}: Results in uncoordinated activity which can be wasteful. \\
  \hline
\end{tabular}
\caption{Similar to \ref{table:forest-vs-trees}, but here the question is about coordination versus doing the work. The amount of coordination depends on how many stakeholders there are, how familiar the stakeholders are with the challenge, and whether the action is reversible when found to be incorrect.
}
\label{table:meetings-versus-work}
\end{table}
\end{center}


\begin{center}
\begin{table}[H] % ht
\begin{tabular}{ | m{\dilemmatablewidth}| m{\dilemmatablewidth} | } 
  \hline
  \textbf{Operate at the level you are being paid for.} &
  \textbf{Operate above the level that you are being paid for in order to be a promoted.} \\
  \hline
  \textit{Description}: Meet job requirements but nothing extra. &
  \textit{Description}: Exceed job requirements. \\
  \hline
  \textit{Cons}: Risk not being promoted. & 
  \textit{Cons}: Experience wage loss - Organization is getting free labor. \\
  \hline
\end{tabular}
\caption{Work above your pay grade (provide the organization extra labor and you get reduced pay) or at your pay grade (expected labor and pay)?
}
\label{table:work_extra_or_work_as_expected}
\end{table}
\end{center}


\begin{center}
\begin{table}[H] % ht
\begin{tabular}{ | m{\dilemmatablewidth}| m{\dilemmatablewidth} | } 
  \hline
  \textbf{Speak out/speak up if something is wrong or \href{https://en.wikipedia.org/wiki/Moral_injury}{offends you}.} &
  \textbf{Hold back comments and questions to minimize disruptions.} \\
  \hline
  \textit{Cons}: You could be missing context; you might look stupid. & 
  \textit{Cons}: You missed an opportunity to correct something; you missed an opportunity to get educated about a situation. \\
  \hline
\end{tabular}
\caption{There is conflicting folk wisdom on both sides of this dilemma: ``The squeaky wheel gets the greese" and ``The squeaky wheel gets replaced." How you raise the issue, with whom, and in what context all matter to either correcting the situation or getting better educated.
}
\label{table:speak-up-or-hold-back}
\end{table}
\end{center}


\begin{center}
\begin{table}[H] % ht
\begin{tabular}{ | m{\dilemmatablewidth}| m{\dilemmatablewidth} | } 
  \hline
  \textbf{Send bad news up the chain of command.} &
  \textbf{Minimize bad news up the chain of command.} \\
  \hline
  \textit{Pros}: You are a reliable source of news. &
  \textit{Pros}: You minimize the burden of managers. \\
  \hline
  \textit{Cons}: You are viewed as a source of problems. & 
  \textit{Cons}: Harmful events eventually catch up with the organization.  \\
  \hline
\end{tabular}
\caption{The canonical example is the \href{https://en.wikipedia.org/wiki/Space_Shuttle_Challenger_disaster}{Challenger disaster}.}
\label{table:bad-news-up-the-chain}
\end{table}
\end{center}



\begin{center}
\begin{table}[H] % ht
\begin{tabular}{ | m{\dilemmatablewidth}| m{\dilemmatablewidth} | } 
  \hline
  \textbf{Prepare for disasters and emergencies, invest in mitigation.} &
  \textbf{Wait for the specific problem to arise before responding.} \\
  \hline
  \textit{Pros}: Lessen the impact when bad things happen; decrease the number of problems from occurring in the first place. &
  \textit{Pros}: deal with the specifics of the scenario at that time and thus be better informed. \\
  \hline
  \textit{Cons}: fewer events evolve into emergencies because you're prepared, or the impact of disasters is lessened. Both make you look overly paranoid and wasteful. & 
  \textit{Cons}: Unexpected events result in worse outcomes.  \\
  \hline
\end{tabular}
\caption{The ``Preparation versus Cleanup'' Dilemma: How much to invest in contingency planning and preparedness.}
\label{table:emergencies-vs-ignore}
\end{table}
\end{center}

% https://graphthinking.blogspot.com/2019/08/two-misleading-simplifications-when.html
\begin{center}
\begin{table}[H] % ht
\begin{tabular}{ | m{\dilemmatablewidth}| m{\dilemmatablewidth} | } 
  \hline
  \textbf{Focus on the immediate problem.} &
  \textbf{Ponder the systemic issues and adjacent contexts.} \\
  \hline
  \textit{Description}: Focus on isolated problematic aspects and do worry about the interdependencies and feedback loops and stakeholder incentives. &
  \textit{Description}: Philosophical musings with a holistic view. \\
  \hline
  \textit{Cons}: Misses systemic issues, causes that exist outside the immediate scope, or issues that occur due to interacting processes. & 
  \textit{Cons}: Relies on knowledge of the wider system that you may have less awareness of. Less emphasis on getting things done. May reveal problems that you don't have authority to address. \\
  \hline
\end{tabular}
\caption{Scope of problem solving can be narrow or broad. Rather than limit your investigations to one or the other, flipping between the two on a recurring basis (but not too frequently) can help.
}
\label{table:focus-vs-systemic}
\end{table}
\end{center}



\begin{center}
\begin{table}[H] % ht
\begin{tabular}{ | m{\dilemmatablewidth}| m{\dilemmatablewidth} | } 
  \hline
  \textbf{Only let good ideas through as determined by a detailed review process of clearly specified plan.} &
  \textbf{Give resources to untested ideas.} \\
  \hline
  \textit{Pros}: Less waste of resources and time. Everyone has confidence in the investment. & 
  \textit{Pros}: High-risk/high-reward ideas that are disruptive can be implemented. \\
  \hline
  \textit{Cons}: Burdensome review process. & 
  \textit{Cons}: some ideas will fail. \\
  \hline
\end{tabular}
\caption{How much vetting should novel ideas get before implementation?
}
\label{table:idea-filtering}
\end{table}
\end{center}



% https://graphthinking.blogspot.com/2016/05/claim-innovation-is-either-disruptive.html
\begin{center}
\begin{table}[H] % ht
\begin{tabular}{ | m{\dilemmatablewidth}| m{\dilemmatablewidth} | } 
  \hline
  \textbf{Work on disruptive innovation.} &
  \textbf{Work on iterative (evolutionary) innovation.} \\
  \hline
  \textit{Description}: Start from scratch, aim for revolution, replace the legacy.  &
  \textit{Description}: Starts with change to existing solution, adjust the legacy path.  \\  
  \hline
  \textit{Pros}: High reward. & 
  \textit{Pros}: Low risk, low cost. \\
  \hline
  \textit{Cons}: High risk, high cost. & 
  \textit{Cons}: Low reward. \\
  \hline
\end{tabular}
\caption{Incremental change may not suffice. Disruption can be costly.
}
\label{table:disruptive-or-iterative}
\end{table}
\end{center}

\begin{center}
\begin{table}[H] % ht
\begin{tabular}{ | m{\dilemmatablewidth}| m{\dilemmatablewidth} | } 
  \hline
  \textbf{Innovate in a novel-to-your-team environment.} &
  \textbf{Innovate in your team's standard environment.} \\
  \hline
  \textit{Pros}: More likely to allows people to drop their expectations.  & 
  \textit{Pros}: Easy to operate in. \\
  \hline
  \textit{Cons}: Loses access to connections vital to actually create success. Use extra space, logistics of moving. & 
  \textit{Cons}: Allows conventional processes to take effect. People hold onto their assumptions. \\
  \hline
\end{tabular}
\caption{Where (physically, spatially) innovation takes place matters because the environment sets context for assumptions.
}
\label{table:where-to-innovate}
\end{table}
\end{center}

% https://graphthinking.blogspot.com/2016/06/innovation-in-open-versus-behind-curtain.html
\begin{center}
\begin{table}[H] % ht
\begin{tabular}{ | m{\dilemmatablewidth}| m{\dilemmatablewidth} | } 
  \hline
  \textbf{Work on innovation in the open.} &
  \textbf{Work on innovation in hiding.} \\
  \hline
  \textit{Pros}: More likely to be criticized. Criticism can be positive (as in an incubator setting).& 
  \textit{Pros}: Less drama -- the incumbent won't attack the innovation since they don't know about it. \\
  \hline
  \textit{Cons}: Negative criticism intended to harm -- an incumbent has reason to be defensive. The incumbent attacks the innovation before it is sufficiently developed or has time to build a user base. & 
  \textit{Cons}: Less opportunity for feedback. Project is easier to kill since the value is not advertised.\\
  \hline
\end{tabular}
\caption{How is innovation carried out within the organization?
}
\label{table:innovate-open-obscure}
\end{table}
\end{center}


\begin{center}
\begin{table}[H] % ht
\begin{tabular}{ | m{\dilemmatablewidth}| m{\dilemmatablewidth} | } 
  \hline
  \textbf{Seek recognition for your work.} &
  \textbf{Work in obscurity.} \\
  \hline
  \textit{Pros}: Helps with promotion. & 
  \textit{Pros}: Less distraction. \\
  \hline
  \textit{Cons}: Devalues the contributions of other people. & 
  \textit{Cons}: No one knows the value of your work and you won't get feedback. \\
  \hline
\end{tabular}
\caption{This dilemma is magnified when the task you work on is high risk or is resource intensive. 
}
\label{table:recognition-obscurity}
\end{table}
\end{center}

\begin{center}
\begin{table}[H] % ht
\begin{tabular}{ | m{\dilemmatablewidth}| m{\dilemmatablewidth} | } 
  \hline
  \textbf{Gather lots of data.} &
  \textbf{Gather minimal data.} \\
  \hline
  \textit{Description}: Gather lots of data for a well-informed decision. &
  \textit{Description}: Minimal information because decision maker knows what to do or outcome is irrelevant.  \\  
  \hline
  \textit{Cons}: High cost of gathering data (time, resources). \href{https://en.wikipedia.org/wiki/Opportunity_cost}{Opportunity costs}. & 
  \textit{Cons}: Lack of data results in decisions based on oversimplified assessment. \\
  \hline
\end{tabular}
\caption{How much data to gather for a decision. See Fig.~\ref{fig:data_collection_cost_uncertainty}. See also \S~\ref{sec:bureaucratic_debt} on bureaucratic debt.
%{\tiny Tag: Decision making.}
}
\label{table:gather_data_lots-vs-little}
\end{table}
\end{center}

\begin{figure}[H] % ht
        \centering
        \includegraphics[width=0.8\textwidth]{images/cost_and_uncertainty_for_data_collection}
        \caption{Collecting more data costs money and decrease uncertainty. See Dilemma~\ref{table:gather_data_lots-vs-little}.}
        \label{fig:data_collection_cost_uncertainty}
\end{figure}

The logistics of gathering data can be measured, but there are other subjective aspects to account for as well. Making a decision has an emotional toll on the decider due to the risk of failure. Also, decisions are made in a social context, with decision makers accounting for the ramifications on people they have relationships with. 


Gathering data (Dilemma~\ref{table:gather_data_lots-vs-little}) is distinct from planning (Dilemma~\ref{table:planning}). It is possible to do a lot of planning with only a little information gathered, and it is feasible to have lots of data and do no planning. 

\begin{center}
\begin{table}[H] % ht
\begin{tabular}{ | m{\dilemmatablewidth}| m{\dilemmatablewidth} | } 
  \hline
  \textbf{Extensive planning upfront (proactive)} & 
  \textbf{Iterative improvement of plans (reactive)} \\ 
  \hline
  \textit{Description}: Lots of time spent brainstorming potential scenarios and contingency options prior to taking action. & 
  \textit{Description}: Start taking action and use feedback to shape next actions. \\ 
  \hline
  \textit{Cons}: ``No plan survives contact with the enemy.'' & 
  \textit{Cons}: Less prepared. \\  
  \hline
\end{tabular}
\caption{How much time to invest in planning.
%{\tiny Tag: Decision making.}
}
\label{table:planning}
\end{table}
\end{center}



\begin{center}
\begin{table}[H]
\begin{tabular}{ | m{\dilemmatablewidth}| m{\dilemmatablewidth} | } 
  \hline
  \textbf{Involve people who disagree.} & 
  \textbf{Ignore people who disagree.} \\ 
  \hline
  \textit{Pros}: Get constructive feedback; account for factors you didn't consider; build a robust solution. & 
  \textit{Pros}: Save time by not interacting. \\  
  \hline
  \textit{Cons}: Results in a compromise or partial solution that minimizes aggregate unhappiness. & 
  \textit{Cons}: Miss a vital aspect you didn't consider. \\  
  \hline
\end{tabular}
\caption{Engagement with opposition to process or change.
%{\tiny Tag: Decision making.}
}
\label{table:opposition}
\end{table}
\end{center}

Making a decision imposes a bound on how much time is available for both gathering data and planning. Time is zero sum, so more time gathering data is less time planning. Similarly, the number of people available for data gathering and planning is bounded, and tasking people is a zero sum choice.

\begin{figure}[H] % ht
    \centering
    \includegraphics[width=0.8\textwidth]{images/planning_and_data_gathering.pdf}
    \caption{Planning (Dilemma~\ref{table:planning}) and data gathering (Dilemma~\ref{table:gather_data_lots-vs-little}) trade-off.}
    \label{fig:pareto_frontier}
\end{figure}

In practice, gathering data and planning rarely terminate -- they evolve.




When planning (Dilemma~\ref{table:planning}), aspects to consider include
%\begin{itemize}
%    \item 
the amount of risk seeking or tolerance (Dilemma~\ref{table:risk})
and
%    \item 
the intended scope of impact  (Dilemma~\ref{table:scope_broad-vs-narrow}).
%\end{itemize}

\begin{center}
\begin{table}[H] % ht
\begin{tabular}{ | m{\dilemmatablewidth}| m{\dilemmatablewidth} | } 
  \hline
  \textbf{Take on big risks and big rewards.} & 
  \textbf{Take on small risks and small rewards.} \\ 
  \hline
  \textit{Description}: High risk tolerance. &
  \textit{Description}: Low risk tolerance. \\
  \hline
  \textit{Pros}: Potential for failure and harm is significant. &
  \textit{Pros}: If any one investment fails, you can continue other efforts. \\
  \hline
  \textit{Cons}: Costly investment, longer feedback cycle. & 
  \textit{Cons}: Incremental can be slower. \\
  \hline
\end{tabular}
\caption{\href{https://en.wikipedia.org/wiki/Risk_assessment}{Risk tolerance}. 
%{\tiny Tag: Personal choice.}
}
\label{table:risk}
\end{table}
\end{center}

\ \\

\begin{center}
\begin{table}[H] % ht
\begin{tabular}{ | m{\dilemmatablewidth}| m{\dilemmatablewidth} | } 
  \hline
  \textbf{Broad scope of impact.} &
  \textbf{Narrow scope of impact.} \\
  \hline
  \textit{Description}: The consequence of the work has many stakeholders. &
  \textit{Description}: Small number of stakeholders. \\  
  \hline
  \textit{Pros}: Benefit more people. &
  \textit{Pros}: Niche impact means less dependencies on other people. \\
  \hline
  \textit{Cons}: Harder to get everyone in agreement. & 
  \textit{Cons}: Less visibility to the rest of the organization. \\
  \hline
\end{tabular}
\caption{Scope of impact of your work. 
%{\tiny Tag: Personal choice}
}
\end{table}
\label{table:scope_broad-vs-narrow}
\end{center}


Once data is gathered (Dilemma~\ref{table:gather_data_lots-vs-little}) and a plan is made (Dilemma~\ref{table:planning}), the result is disseminated. The choice on how to disseminate is Dilemma~\ref{table:consistency} and Dilemma~\ref{table:disseminate_one-by-one}.

\begin{center}
\begin{table}[H] % ht
\begin{tabular}{ | m{\dilemmatablewidth}| m{\dilemmatablewidth} | } 
  \hline
  \textbf{Guidance updated frequently; Incremental change.} & 
  \textbf{Consistent application of policy over time. Rules persist; then sudden drastic change.} \\ 
  \hline
  \textit{Pros}: Adapt policy to new information and changing conditions. &
  \textit{Pros}: Stability is easier to predict between regime changes.  \\
  \hline
  \textit{Cons}: More work needed. Accused of lacking stability. & 
  \textit{Cons}: Doesn't adapt as conditions change. Accused of being inflexible to evolving conditions. \\
  \hline
\end{tabular}
\caption{Consistency over time. Stability of rules; how change is implemented. Can also be characterized as when to tell other people: sooner or later (when firmer information is available) See \S~\ref{sec:static-dynamic_processes} for static versus dynamic processes.
%{\tiny Tag: Organization's culture. Tag: Personal choice.}
}
\label{table:consistency}
\end{table}
\end{center}

Deployment of products and deployment of policies face similar dilemmas. \href{https://en.wikipedia.org/wiki/Diffusion_of_innovations}{Diffusion of Innovation}

\begin{center}
\begin{table}[H] % ht
\begin{tabular}{ | m{\dilemmatablewidth}| m{\dilemmatablewidth} | } 
  \hline
  \textbf{Tell people one-by-one.} & 
  \textbf{Tell everyone at once.} \\ 
  \hline
  \textit{Pros}: One-on-one allows a freer response from audience. &
  \textit{Pros}: Save time for the speaker. \\
  \hline
  \textit{Cons}: Order matters for relationships. & 
  \textit{Cons}: Overwhelming feedback all at once. \\  
  \hline
\end{tabular}
\caption{How to disseminate information.
%{\tiny Tag: Personal choice.}
}
\label{table:disseminate_one-by-one}
\end{table}
\end{center}

Once a decision has been made, the decision is executed or enforced. How many rules are there (Dilemma~\ref{table:number_of_rules}) and
how strictly are the rules enforced (Dilemma~\ref{table:rule_strictness})?

\begin{center}
\begin{table}[H] % ht
\begin{tabular}{ | m{\dilemmatablewidth}| m{\dilemmatablewidth} | } 
  \hline
  \textbf{Enforce rules strictly.} & 
  \textbf{Lax rule enforcement.} \\ 
  \hline
  \textit{Pros}: Predictable. &
  \textit{Pros}: Bureaucrats feel empowered. \\
  \hline
  \textit{Cons}: Insensitive to nuance. & 
  \textit{Cons}: Tolerance for changing conditions or exceptional cases.  \\  
  \hline
\end{tabular}
\caption{Strictness of rules.
%{\tiny Tag: Organization's culture.}
}
\label{table:rule_strictness}
\end{table}
\end{center}

\begin{center}
\begin{table}[H] % ht
\begin{tabular}{ | m{\dilemmatablewidth}| m{\dilemmatablewidth} | } 
  \hline
  \textbf{If it's not against the rules, it must be okay.} & 
  \textbf{I can only do what is allowed by the rules and nothing more.} \\ 
  \hline
  \textit{Pros}: Autonomy &
  \textit{Pros}:  \\
  \hline
  \textit{Cons}: . & 
  \textit{Cons}: .  \\  
  \hline
\end{tabular}
\caption{Adherence to rules.
}
\label{table:rule_adherence}
\end{table}
\end{center}





\begin{center}
\begin{table}[H] % ht
\begin{tabular}{ | m{\dilemmatablewidth}| m{\dilemmatablewidth} | } 
  \hline
  \textbf{Control via rules.} & \textbf{Freedom/autonomy/agility.} \\ 
  \hline
  \textit{Description}: High number of rules to cover a variety of situations. & 
  \textit{Description}: Low number of rules to enable flexibility. \\ 
  \hline
  \textit{Cons}: The more rules that exist the more likely it is that someone will find a way to exploit them to their own advantage. & 
  \textit{Cons}: The fewer rules that exist the more likely it is that someone will try to get away with something bad. \\  
  \hline
\end{tabular}
\caption{Number of rules.
%{\tiny Tag: Organization's culture}
}
\label{table:number_of_rules}
\end{table}
\end{center}
Alternative approach: guidance derived from principles that can be adapted to specific situations. That has the problem of requiring good knowledge of the situation and wise judgement.

\ \\

\begin{center}
\begin{table}[H] % ht
\begin{tabular}{ | m{\dilemmatablewidth}| m{\dilemmatablewidth} | } 
  \hline
  \textbf{I can only do what is mandated by the organization.} & 
  \textbf{I can do anything that's not illegal.} \\ 
  \hline
  \textit{Cons}:  &
  \textit{Cons}:  \\  
  \hline
\end{tabular}
\caption{The scope of your actions bound by mandates and legality, but the way you interpret that is subjective. 
}
\label{table:legality}
\end{table}
\end{center}


\begin{center}
\begin{table}[H] % ht
\begin{tabular}{ | m{\dilemmatablewidth}| m{\dilemmatablewidth} | } 
  \hline
  \textbf{Quickly complete tasks.} & 
  \textbf{Methodically complete tasks.} \\ 
  \hline
  \textit{Description}: implementing a solution quickly to address urgent needs. &
  \textit{Description}: Methodical well-planned design and execution yield robust solutions. \\
  \hline
  \textit{Pros}: Rapid solution. &
  \textit{Pros}: More like to get the solution right. \\
  \hline
  \textit{Cons}: Risk of quick task is that the result is ineffective, inefficient, or wrong. &
  \textit{Cons}: \href{https://en.wikipedia.org/wiki/Opportunity_cost}{opportunity cost} \\  
  \hline
\end{tabular}
\caption{Speed versus accuracy of task completion.
}
\label{table:quick-methodical}
\end{table}
\end{center}

\ \\

\begin{center}
\begin{table}[H] % ht
\begin{tabular}{ | m{\dilemmatablewidth}| m{\dilemmatablewidth} | } 
  \hline
  \textbf{Push people to work really hard.} & 
  \textbf{Create a comfortable work environment.} \\ 
  \hline
  \textit{Cons}: Burn out and leave. & 
  \textit{Cons}: Lower instantaneous productivity. \\  
  \hline
\end{tabular}
\caption{Policy enforcement rate
}
\label{table:rate-of-work}
\end{table}
\end{center}

\ \\

% https://bennorthrop.com/Essays/2022/code-ownership-stewardship-or-free-for-all.php
\begin{center}
\begin{table}[H] % ht
\begin{tabular}{ | m{\dilemmatablewidth}| m{\dilemmatablewidth} | } 
  \hline
  \textbf{One person or team owns an area of responsibility.} & 
  \textbf{Anyone take on any task.} \\ 
  \hline
  \textit{Cons}: Staffing capacity may not be as flexible as varying workload. & 
  \textit{Cons}: Not everyone is skilled at everything. \\  
  \hline
\end{tabular}
\caption{Swimlanes, task boundaries.
}
\label{table:swimlanes}
\end{table}
\end{center}


\ \\

\begin{center}
\begin{table}[H] % ht
\begin{tabular}{ | m{\dilemmatablewidth}| m{\dilemmatablewidth} | } 
  \hline
  \textbf{Seek out experienced collaborators.} & 
  \textbf{Work with less experienced people.} \\ 
  \hline
  \textit{Pros}: Quicker to get something done. &
  \textit{Pros}: Less set in their ways and open to more novelty. \\  
  \hline
  \textit{Cons}: Experienced people who are good are probably busy. &
  \textit{Cons}: Slower progress. \\  
  \hline
\end{tabular}
\caption{People with experience are useful but less accessible.
}
\label{table:experience}
\end{table}
\end{center}


\ \\

\begin{center}
\begin{table}[H] % ht
\begin{tabular}{ | m{\dilemmatablewidth}| m{\dilemmatablewidth} | } 
  \hline
  \textbf{Say yes to new opportunities.} & 
  \textbf{Say no to new opportunities.} \\ 
  \hline
  \textit{Pros}: Positive attitude, collaborative. &
  \textit{Pros}: Able to prioritize and focus. \\
  \hline
  \textit{Cons}: Fail to complete tasks. &
  \textit{Cons}: Not a team player. \\  
  \hline
\end{tabular}
\caption{Acceptance or rejection of additional work.
}
\label{table:new-opportunties}
\end{table}
\end{center}

\ \\

\begin{center}
\begin{table}[H] % ht
\begin{tabular}{ | m{\dilemmatablewidth}| m{\dilemmatablewidth} | } 
  \hline
  \textbf{Share less data.} &
  \textbf{Share more data.} \\
%  \hline
%  \textit{Description}:  &
%  \textit{Description}:  \\  
  \hline
  \textit{Pros}: Restricting data access saves money for the data owner and improves flexibility.&
  \textit{Pros}: Sharing data improves transparency and accountability. \\
  \hline
  \textit{Cons}: Other people (inside and outside the organization) are unable to extract maximum value from data & 
  \textit{Cons}: Sharing data cost resources (people, money, time) \\
  \hline
\end{tabular}
\caption{How much data to share.
%{\tiny Tag: Personal choice.}
}
\label{table:data_share-vs-hide}
\end{table}
\end{center}

\ \\

\begin{center}
\begin{table}[H] % ht
\begin{tabular}{ | m{\dilemmatablewidth}| m{\dilemmatablewidth} | } 
  \hline
  \textbf{Compete for resources.} &
  \textbf{Cooperate for productivity.} \\
  \hline
  \textit{Description}: individuals compete for attention and promotion; teams compete for money and staffing resources &
  \textit{Description}: cooperation improves productivity \\  
  \hline
  \textit{Cons}: Fail to synergize skills resources & 
  \textit{Cons}: Not clear who to assign responsibility for success or failure \\
  \hline
\end{tabular}
\caption{Cooperate or Compete -- applies to teams and to individuals. 
%{\tiny Tag: Personal choice.}
}
\label{table:cooperate-vs-compete}
\end{table}
\end{center}

\ \\

\begin{center}
\begin{table}[H] % ht
\begin{tabular}{ | m{\dilemmatablewidth}| m{\dilemmatablewidth} | }
  \hline
  \textbf{Consistent application of policy across cases.} &
  \textbf{Adapt policy to specific cases.} \\
  \hline
  \textit{Description}: Maximize broad applicability; minimize exceptions. &
  \textit{Description}: Demonstrate flexibility for unique scenarios. \\  
  \hline
  \textit{Cons}: Less sensitive to the nuances of a specific situation. & 
  \textit{Cons}: Takes more work. More likely to be accused of bias. \\
  \hline
\end{tabular}
\caption{Case consistency vs adaptability.
%{\tiny Tag: Personal choice.}
}
\label{table:policy_consistency_across_cases}
\end{table}
\end{center}

When change to policies is desired, there are options on how to advocate for change -- Dilemma~\ref{table:how_to_change}.

\begin{center}
\begin{table}[H] % ht
\begin{tabular}{ | m{\dilemmatablewidth}| m{\dilemmatablewidth} | } 
  \hline
% https://graphthinking.blogspot.com/2019/07/not-too-loose-not-too-tight-determining.html
  \textbf{Adhere strictly to the scope of your role.} & 
  \textbf{Stray outside (or outright ignore) the scope of your role.} \\ 
  \hline
  \textit{Description}: Inflexible to novelty. Specialization of tasking. & 
  \textit{Description}: Lack of structure. Generalization. \\ 
  \hline
  \textit{Cons}: efficiencies of cooperation and specialization would not occur. & 
  \textit{Cons}: Deadlock condition arises due to a scheduling constraint -- no one can proceed because everyone is waiting on everyone else. Responsibilities are unclear when no ones scope is clear. \\  
  \hline
\end{tabular}
\caption{Both strict adherence to role scope and ignoring scope can decrease an organization's productivity. 
What happens when a person deviates from their role?
How are people who do not conform identified? Are they confronted?
}
\label{table:scope_of_activity}
\end{table}
\end{center}

\begin{center}
\begin{table}[H] % ht
\begin{tabular}{ | m{\dilemmatablewidth}| m{\dilemmatablewidth} | } 
  \hline
  \textbf{Delegate; share work with other people.} & 
  \textbf{Work alone; don't rely on other people.} \\ 
  \hline
  \textit{Cons}: Your success is dependent on other people. & 
  \textit{Cons}: Can't accomplish as much on your own. \\  
  \hline
\end{tabular}
\caption{Sharing work can improve productivity and build relationships but also incurs risks to reputation and success.
}
\label{table:delegate-or-not}
\end{table}
\end{center}

\begin{center}
\begin{table}[H] % ht
\begin{tabular}{ | m{\dilemmatablewidth}| m{\dilemmatablewidth} | } 
  \hline
  \textbf{Many small tasks or objectives.} & 
  \textbf{Fewer big tasks or objectives.} \\ 
  \hline
  \textit{Cons}: Enables a fail fast approach. & 
  \textit{Cons}: Less overhead to manage. \\  
  \hline
\end{tabular}
\caption{Chunk size.
}
\label{table:chunk_size}
\end{table}
\end{center}

\begin{center}
\begin{table}[H] % ht
\begin{tabular}{ | m{\dilemmatablewidth}| m{\dilemmatablewidth} | } 
  \hline
  \textbf{Task with many external dependencies.} & 
  \textbf{Task with few external dependencies.} \\ 
  \hline
  \textit{Cons}: Risk of failing because of a failed dependency. & 
  \textit{Cons}: Have to develop everything yourself; waste of resources due to redundancy. \\  
  \hline
\end{tabular}
\caption{External dependencies can enable broader scope. See also the Dilemma of Delgation, \ref{table:delegate-or-not}.
}
\label{table:number_of_external dependencies}
\end{table}
\end{center}

\begin{center}
\begin{table}[H] % ht
\begin{tabular}{ | m{\dilemmatablewidth}| m{\dilemmatablewidth} | } 
  \hline
  \textbf{Focused on one role.} & 
  \textbf{Have multiple roles.} \\ 
  \hline
  \textit{Cons}: If a role does not consume 40 hours per week, you'll be idle. & 
  \textit{Cons}: Context switches between roles and delayed responses. \\  
  \hline
\end{tabular}
\caption{The right number of roles for a bureaucrat depends on personality and tasking. 
}
\label{table:number_of_roles}
\end{table}
\end{center}


\begin{center}
\begin{table}[H] % ht
\begin{tabular}{ | m{\dilemmatablewidth}| m{\dilemmatablewidth} | } 
  \hline
  \textbf{Dissent is welcome and discussed freely.} & 
  \textbf{Dissent is suppressed.} \\ 
  \hline
  \textit{Cons}: Can be disruptive to normal operations. Distracts from the task. & 
  \textit{Cons}: Limits novel ideas from spreading. Harms morale. \\  
  \hline
\end{tabular}
\caption{Dissent is caused by dissatisfaction with people or processes. 
}
\label{table:how_dissent_is_responded_to}
\end{table}
\end{center}


\begin{center}
\begin{table}[H] % ht
\begin{tabular}{ | m{\dilemmatablewidth}| m{\dilemmatablewidth} | } 
  \hline
  \textbf{Do share lessons learned.} & 
  \textbf{Don't share lessons learned.} \\ 
  \hline
  \textit{Pros}: Honesty, accountability, self-awareness, and self-reflection. & 
  \textit{Pros}: Look competent, even when making mistakes. \\  
  \hline
  \textit{Cons}: Looks weak, unprofessional. & 
  \textit{Cons}: Limit the growth of bureaucrats in the organization. \\  
  \hline
\end{tabular}
\caption{Sharing lessons learned may seem reasonable unless you want to maintain a pristine reputation. 
}
\label{table:sharing_lessons_learned}
\end{table}
\end{center}

% https://graphthinking.blogspot.com/2019/07/vulnerability-of-organizations-in.html
\begin{center}
\begin{table}[H] % ht
\begin{tabular}{ | m{\dilemmatablewidth}| m{\dilemmatablewidth} | } 
  \hline
  \textbf{Share lessons learned about yourself.} & 
  \textbf{Share lessons learned from observing others.} \\ 
  \hline
  \textit{Cons}: Potentially look stupid. & 
  \textit{Cons}: Potentially hurts their reputation. \\  
  \hline
\end{tabular}
\caption{When sharing lessons learned (option 1 in \ref{table:sharing_lessons_learned}), the lessons do not have to be about you. 
}
\label{table:share_lessons_learned}
\end{table}
\end{center}

\begin{center}
\begin{table}[H] % ht
\begin{tabular}{ | m{\dilemmatablewidth}| m{\dilemmatablewidth} | } 
  \hline
  \textbf{Learn lessons from your own mistakes.} & 
  \textbf{Learn lessons from others (formal training).} \\ 
  \hline
  \textit{Cons}: Potentially look stupid; waste resources discovering what others already know. & 
  \textit{Cons}: Formal training may overemphasize irrelevant or impractical concepts. \\  
  \hline
\end{tabular}
\caption{How much formal training to invest in before learning by doing?
}
\label{table:lessons_learned_source}
\end{table}
\end{center}



\begin{center}
\begin{table}[H] % ht
\begin{tabular}{ | m{\dilemmatablewidth}| m{\dilemmatablewidth} | } 
  \hline
  \textbf{Build a small coalition of interested parties.} & 
  \textbf{Build a large base of support and get everyone on board.} \\ 
  \hline
  \textit{Cons}: May not be representative of all stakeholders. & 
  \textit{Cons}: Takes time away from the work. Many people may disagree or be disinterested. \\  
  \hline
\end{tabular}
\caption{A coalition can provide morale support but takes time to build.
%{\tiny Tag: }
}
\label{table:how_to_change}
\end{table}
\end{center}

\subsection{Dilemmas of Policy for an Organization's Structure\label{sec:org_dilemma}}

The constraints a decision maker faces are informed by the person's environment. Dilemmas \ref{table:people-per-supervisor} through \ref{table:market-vs-monopoly} shape the experience of bureaucrats in an organization.

\begin{center}
\begin{table}[H] % ht
\begin{tabular}{ | m{\dilemmatablewidth}| m{\dilemmatablewidth} | } 
  \hline
  \textbf{Flatter hierarchical organization.} &
  \textbf{More layers of hierarchy.} \\ 
  \hline
  \textit{Description}: More people managed per supervisor. & 
  \textit{Description}: Fewer people managed per supervisor. \\ 
  \hline
  \textit{Cons}: Less feedback/attention per employee. & 
  \textit{Cons}: Fewer people doing work. \\  
  \hline
\end{tabular}
\caption{Shape of hierarchical organization.
%{\tiny Tag: Organization's culture}
}
\label{table:people-per-supervisor}
\end{table}
\end{center}

\ \\

\begin{center}
\begin{table}[H] % ht
\begin{tabular}{ | m{\dilemmatablewidth}| m{\dilemmatablewidth} | } 
  \hline
  \textbf{Staffing: good coverage.} &
  \textbf{Staffing: minimal coverage.} \\
  \hline
  \textit{Description}: Sufficient staff. &
  \textit{Description}: As small of staff as possible. \\  
  \hline
  \textit{Pros}: Cover all edge cases; resilient to changing demands. &
  \textit{Pros}: Less expensive. \\
  \hline
  \textit{Cons}: Slack resources; sometimes inefficient. Increased communication needed. & 
  \textit{Cons}: Fragile when requirements change or workload increases. If one person departs and there's no redundancy, capacity and capability are harmed.  \\
  \hline
\end{tabular}
\caption{Size of team or organization.
%{\tiny Tag: Design of organization.}
}
\label{table:staff_many-vs-few}
\end{table}
\end{center}


\ \\

\begin{center}
\begin{table}[H] % ht
\begin{tabular}{ | m{\dilemmatablewidth}| m{\dilemmatablewidth} | } 
  \hline
  \textbf{In-house services for non-central activities.} &
  \textbf{External dependencies for non-central activities.} \\
%  \hline
%  \textit{Description}:  &
%  \textit{Description}:  \\  
  \hline
  \textit{Pros}: More control. &
  \textit{Pros}: Easier to replace. \\
  \hline
  \textit{Cons}: Expands scope of responsibilities. & 
  \textit{Cons}: Less understanding of problem.  \\
  \hline
\end{tabular}
\caption{Services that are necessary but not central.
%{\tiny Tag: Design of organization.}
}
\label{table:inhouse-vs-external}
\end{table}
\end{center}

Table~\ref{table:inhouse-vs-external} isn't unique to bureaucracy. It applies to businesses in a market and to governments in a global environment. 

\ \\

\begin{center}
\begin{table}[H] % ht
\begin{tabular}{ | m{\dilemmatablewidth}| m{\dilemmatablewidth} | } 
  \hline
  \textbf{Centralized services.} &
  \textbf{Locally distributed services.} \\
  \hline
  \textit{Description}: A single provider of services for the organization; typically top-down mandate. &
  \textit{Description}: Each team has a local service provider; typically bottom-up organic result. \\  
  \hline
  \textit{Pros}: Cheaper. Enables coordination. &
  \textit{Pros}: Quicker response. 
  Enables innovation. 
  Accounts for local deviations from the norm. \\
  \hline
  \textit{Cons}: Less sensitive to local issues. Less responsive. Longer delays. Single point of failure.  & 
  \textit{Cons}: Uneven quality of service. Inconsistent strategies and policies. \\
  \hline
\end{tabular}
\caption{Centralization of services. Oscillation (see Fig.~\ref{fig:central-vs-distributed}) indicates neither solution is optimal.
%{\tiny Tag: Design of organization.}
}
\label{table:central-vs-distributed}
\end{table}
\end{center}

\begin{figure}[H] % ht
    \centering
    \includegraphics[width=0.8\textwidth]{images/dilemma_centralization-vs-distributed.pdf}
    \caption{Dipole oscillation. See Dilemma~\ref{table:central-vs-distributed}. Migrating to the opposite paradigm gives people in charge a chance to show their responsiveness to the needs of participants. The rate of oscillation is a measure of institutional memory half-life.}
    \label{fig:central-vs-distributed}
\end{figure}

Centralization is often carried out for the purposes of cost efficiency. The cost savings are due to de-duplication and having less slack. Both of those ``savings'' are a decrease of redundancy, which has a cost when there are unexpected fluctuations in need. 

Centralization is an intentional monopolization, with the corresponding decrease in choices. 
Centralization (de-localization) can decrease the value assigned to feedback from people using the service because personal relations are de-valued. 

The weaker feedback, lack of redundancy, and decreased emphasis on relationships motivates the creation of local services. 

\ \\

\begin{center}
\begin{table}[H] % ht
\begin{tabular}{ | m{\dilemmatablewidth}| m{\dilemmatablewidth} | } 
  \hline
  \textbf{Decision making lower in a hierarchy} &
  \textbf{Decision making higher in a hierarchy} \\
  \hline
  \textit{Description}: Push decisions down to empower employees. &
  \textit{Description}: Escalate every decision to management. \\  
  \hline
  \textit{Pros}: better information &
  \textit{Pros}: better scope \\
  \hline
  \textit{Cons}: more inconsistency & 
  \textit{Cons}: Decision maker has less skin in the game and may be less well informed. Employees are disempowered. \\
  \hline
\end{tabular}
\caption{Where decisions get made in hierarchical organization.
%{\tiny Tag: Design of organization.}
}
\label{table:decisions_low-vs-high}
\end{table}
\end{center}

\ \\

\begin{center}
\begin{table}[H] % ht
\begin{tabular}{ | m{\dilemmatablewidth}| m{\dilemmatablewidth} | } 
  \hline
  \textbf{Redundant services in a market} &
  \textbf{Monopoly service provider} \\
  \hline
  \textit{Description}: using a market model within the organization &
  \textit{Description}:  \\  
  \hline
  \textit{Pros}: enable customers to choose the best service &
  \textit{Pros}: efficiency of a single service \\
  \hline
  \textit{Cons}: redundancy & 
  \textit{Cons}: might not meet the needs of all customers \\
  \hline
\end{tabular}
\caption{Services within an organization. See also Fig.~\ref{fig:market-vs-monopoly}.
%{\tiny Tag: Design of organization.}
}
\label{table:market-vs-monopoly}
\end{table}
\end{center}


\begin{figure}[H] % ht
    \centering
    \includegraphics[width=0.8\textwidth]{images/dilemma_market_vs_monopoly.pdf}
    \caption{Dipole oscillation: solution A exists but doesn't meet my needs. Rather than tweak A, re-invent solution B which mostly overlaps with A but has independent development and support. See also Dilemma~\ref{table:market-vs-monopoly}.}
    \label{fig:market-vs-monopoly}
\end{figure}


% https://graphthinking.blogspot.com/2021/04/laffer-curve-and-minimum-viable.html
There is a Goldilocks situation for amount of processes in an organization:
\begin{itemize}
    \item Too few processes (all social relationships).
    \item Just the right amount -- a balance process and social, and when to use which is known.
    \item Too much process (not enough leveraging of social relationships)
\end{itemize}
This optimization is similar to the \href{https://en.wikipedia.org/wiki/Laffer_curve}{Laffer curve} in economics.

\ \\

% https://graphthinking.blogspot.com/2021/12/hierarchical-organization-trilemma.html
Trilemma: \textbf{Do you work for your team, your manager, or yourself?}
Being a member of a team that operates within a hierarchy is tough. One reason is the question of who you are working for. The trilemma is whether you work for yourself, work for your supervisor, or work for your team.  Ideally you can find ways to do all three, but that is not always the case. 

Members of the team should work collaboratively, but there is a potential counter-force: accountability to the supervisor. Because each team member is accountable to their supervisor(s), that motivates the action of the individual. The team does not actually have autonomy -- they are accountable to the boss.

In the approach ``team members are directed by their supervisor," the synergy of the team is neglected and the supervisor becomes a bottleneck (for decision making and for creativity and for planning).

The third approach is for a person to ignore their team and their supervisor. This might enable quicker progress, at the risk of going in an unhelpful direction or not leveraging skills of coworkers. 

\ \\

Trilemma:
\textbf{Seek less budget, same, or larger budget.} Less budget is needed if you improved efficiency, but then the proportional power within the organization is decreased. A larger budget may be due to inefficiency or growth. Keeping the same budget indicates no promotions are relevant (although a steady state could result from a combination of growth and improved efficency). 

\ \\

A trilemma applicable to many situations is that options are \textbf{fast, inexpensive, good; choose two}. (This is the \href{https://en.wikipedia.org/wiki/Project_management_triangle}{Project management triangle}.) \\
In other words, the options are
\begin{itemize}
    \item Good and fast is expensive (i.e., requires lots of resources).
    \item Good and inexpensive takes a long time (i.e., a clever solution).
    \item Fast and inexpensive will be low quality.
\end{itemize}

\subsubsection{Consequences of the Dilemma-based Framing}

The dilemmas listed above are numerous but not exhaustive. Even if each bureaucrat considers the same choices (which doesn't necessarily happen), not everyone makes the same selection. One response might be to defer to someone higher in the hierarchy to coordinate. While this would ensure consistency, this would be \href{https://en.wikipedia.org/wiki/Micromanagement}{micromanagement}. The people in the hierarchy above the person facing the choice don't have exposure to the problem. Choices are delegated to leverage expertise. 

\ \\

When you are facing these dilemmas and trilemmas\marginpar{[Tag] Actionable Advice} there are constructive responses that can improve your effectiveness. Improvement benefits your results, your reputation, and the organization. 
\begin{itemize}
    \item Collect quantitative data for each variable. Quantitative arguments can augment qualitative stories. 
    \item Construct the Pareto frontier to identify non-optimal choices that can be eliminated.
    \item Instead of assessing variables in isolation, assess consequences in the context of workflows and impacts on stakeholders.
    \item Discuss subjective decisions with stakeholders so that potential disagreements can be negotiated instead of creating friction.
\end{itemize}
 

  % section
  \clearpage
  \section{Unavoidable Hazards of Bureaucracy\label{sec:unavoidable-hazards}}

In a bureaucracy there are certain challenges that cannot be avoided. The value of recognizing them is to understand that what you're experiencing isn't an anomaly. The problem isn't unique to you, your circumstances, your coworkers, or the organization. The cause is the combination of all those factors.

\ \\

\begin{samepage}
\textbf{Separation of responsibility and accountability}. \\
Each bureaucrat in an organization has responsibilities associated with their role. The ability to complete the tasks associated the role are not wholly within the scope of the bureaucrat's control. Even if there is a desire for action, the action might not be immediately feasible because of a dependence on another person or process. 
\end{samepage}

\ \\

\begin{samepage}
\textbf{Appreciation for being an effective bureaucrat is rare.}\\
If you do your job well, at best no one will notice.

``No one [thanks teachers] for policing cheating. Not the cheaters, not the honest students who feel inconvenienced and mistrusted, and certainly not the school [administrators] who have to process academic dishonesty paperwork.''\footnote{https://dynomight.net/teaching/}
\end{samepage}

The number of thank you cards sent to \href{https://www.fda.gov/}{Food and Drug Administration} meat inspectors, \href{https://www.osha.gov/}{Occupational Safety and Health Administration} regulators, \href{https://www.fcc.gov/}{Federal Communications Commission}, \href{https://www.ftc.gov/}{Federal Trade Commission}, and \href{https://www.sec.gov/}{Securities and Exchange Commission} is likely small. 
% TODO: ask each agency how many thank you letters they receive

There are counter examples in public service bureaucracy. 
Law enforcement is thanked when there is a victim of a crime. The military is held in high regard. 

\ \\

\begin{samepage}
\textbf{Lack of accountability to fellow bureaucrats who rely on you.}\\
You are more likely to feel accountable to your supervisor. Performing for your supervisor and explaining your value matters more than the work done to support coworkers.  This also applies when you depend on other people -- they aren't accountable to you.
\end{samepage}

\ \\

\textbf{Decision-makers are under-informed and inexperienced.}\\
You can make a decision with insufficient information, and there is rarely an expectation of expertise or experience. 
Gathering information takes time and is thus burdensome.
Having experience requires getting experience.

\ \\

\textbf{Gathering data for decisions takes resources and expertise.}\\
When there is a desire to gather data, and there is data available to be gathered, an investment of time is necessary. Well-informed decisions take time, experience, or both.

\ \\

\textbf{Defining success is subjective and dynamic}. \\
Who defines success and for which audience in a bureaucratic organization is subjective because of the lack of feedback mechanisms. Consequences may not be immediately obvious. Worse, how success is defined can be changed at any time -- there's no need for consistency. 

\ \\

\textbf{Change threatens incumbents}. \\
Change of plans, roles, tasks, resources, flatness of org, scope, technology.

\ \\

\textbf{\href{https://en.wikipedia.org/wiki/Diffusion_of_responsibility}{Diffusion of responsibility} in the bureaucracy} \\
A specific task needs to be completed, and action requires involvement of multiple collaborators. 

Once responsibility is assigned, the new risk is that of blame.

\ \\

\textbf{Diffusion of blame in the bureaucracy}. \\
Example: if something goes wrong, who is at fault: the creator of a policy or the person tasked with executing the policy?

\ \\

\textbf{High latency feedback}. \\
See 
%section~\ref{sec:slow_deployment} 
\hyperref[sec:slow-deployment]{slow deployment of processes} and 
%section~\ref{sec:decision-delay} 
 \hyperref[sec:decision-delay]{decision delay}.

\ \\

\textbf{Weak feedback}. \\
Suppose my bureaucratic task depends on a service provided by you, a fellow bureaucrat. If your computer isn't working, your ability to provide the service is blocked by a dependency on a computer repair person. You might feel relieved -- you can relax while you wait on the computer repair and you don't have to do work providing the service to me. You might feel anxious -- you're unable to provide the service until the computer is fixed. Regardless of how you feel about the situation, my task is blocked. If neither you nor I have priority relative to the computer repair person's other tasks, then we wait. This question of priority could be resolved hierarchically (for which there is limited attention bandwidth) or socially (dependent on someone having a relationship with the repair person). 

The delay for my task may be sufficiently small that raising the issue through the hierarchy or socially is not a good use of time.

\ \\

\textbf{The person making the rules that you follow doesn't actually know what they're doing}. \\
Your choices then include
\begin{itemize}
    \item Follow the rules that are not correct. This harms your productivity and morale. 
\item You can violate the rules and be more effective. This puts you at risk for sanctions. 
\item You can work the change the rules. Then you are not doing the work that's needed.
\end{itemize}


\ \\

\textbf{You rarely get to pick who is on the team}. \\
When a task requires collaboration, there is rarely a choice of who you get to work with. 

\ \\

% TODO
\textbf{Rarely able to alter team membership.}

\ \\

\textbf{Fear of the unknown.}\\
Current suffering is tolerable compared to the uncertainty of change, especially when the suffering isn't felt by the decision maker.

By identifying this fear in yourself or those you collaborate with, you can discuss specific concerns and work to address them.
\marginpar{[Tag] Actionable Advice} 

\ \\

\textbf{Fear of change.} \\
Change disrupts the status quo, putting accumulated power at risk and altering relationships. 

An emotional basis for decision making (or avoidance of decision making) may or may not be rational. Discussing the fear with fellow bureaucrats can ease the burden. 

\ \\

\textbf{Fear of conflict.} Refers to conflict of ideas, not physical conflict. Conflict of ideas is not personal, though the consequences may have impacts on careers for stakeholders. \\
Professional disagreement is to be expected in a bureaucratic organization.

Feeling uncomfortable is distinct from feeling unsafe. Communicating while feeling discomfort is a useful skill, rather than avoidance. You can demonstrate courage by talking about your sense of discomfort. 

\ \\

\textbf{Inadequate resources: staffing, time, money.}\\
The resources you have to address a challenge may not match the complexity of the issue.

\ \\

\textbf{Outcomes for the team are ill-defined and constantly shifting.}\\

\ \\

\textbf{The reward for good work is more work.}\\
When a coworker doesn't do their fair share, then the productive employee shoulders more burden. The deficient worker has no incentive to improve

\ \\

\textbf{The organization has a lack of vision; or has vision but no plan; or has vision and a plan but no consensus; or has vision and a plan and consensus but inadequate resources.}

\ \\

% TODO
\textbf{Progress depends on subjective decision making.}\\
Rarely is the optimal path deterministic. 

\ \\

\textbf{Easier to ask for big money than small money}\\
Processes are scale invariant -- Regardless of whether you're asking for \$1000 or \$100,000, the process is the same. 

\ \\

\textbf{Flux of people and processes.} \\
staff turn-over, changing conditions, changing timelines, change of vision, and the need to be promoted. in a bureaucracy, consistency doesn't yield promotion.

\ \\

\textbf{Why are there so many rules?}\\
To address edge cases and malicious or dumb people.

\ \\

\textbf{So much paperwork/forms.}
Why is Red tape endemic to bureaucracy?\\
Paperwork as a form of coordination in processes to facilitate decentralized decision making. 

One of the motives for paperwork is to catch people who are misrepresenting. Documentation makes prosecution more straightforward.

\ \\

\textbf{Everything is slower.}\\
What is meant is "slower than desired" or "slower than I imagined in my simplified model"

One reason progress is slower than expected is because there aren't as many hours available as imagined.


% What's the point of this? Is there a consequence for the reader?

In a bureaucracy the actual time spent working is less than the number of hours you get paid for. Breaks during work, vacation from work, holidays, sick leave. 

% https://rescuetime.wpengine.com/work-life-balance-study-2019/

\begin{figure}[H]
    \centering
    \includegraphics[width=0.8\textwidth]{images/hours_per_activity_per_employed_year}
    \caption{Hours of ``work'' per year when accounting for the rest of life. Assumes 5 weeks of vacation, 2 days of sick leave, and 11 holidays.}
    % footnotes in caption is not recommended; see https://texfaq.org/FAQ-ftncapt
    % however it can be done; see https://stackoverflow.com/questions/67621322/footnote-in-caption-of-figure-on-latex
    \label{fig:hours_per_year}
\end{figure}

Fig~\ref{fig:hours_per_year}
\footnote{\href{https://docs.google.com/spreadsheets/d/1ZaOZZXWkEzX4fFltUdlR4A6ENrAXnkzTW4YrjA4tDO8/edit?usp=sharing}{source for calculations}}

\ \\

\textbf{Bureaucracy inhibits creativity.}\\
See 
%section~\ref{sec:innovation} for 
notes on \hyperref[sec:innovation]{innovation}.

\ \\

% https://graphthinking.blogspot.com/2017/06/measuring-growth-of-bureaucracy.html
\textbf{Bureaucracy grows.}\\
The causes of growth include
scope creep, an increased number of people participating in the process, an increased workload, or the ratio of workers to managers decreases. 

You can measure the growth of bureaucracy using metrics like the latency in addressing a request and 
the number of people involved in addressing a request.  

\textbf{Increasing the size of a bureaucracy is easier than cutting it down.}\\
\ \\

% https://graphthinking.blogspot.com/2017/04/growth-of-bureaucracy.html
\textbf{The larger a hierarchy, the ratio of workers to managers gets worse}. \\

As complexity and workload increase, the number of staff needed increase. To facilitate coordination, managers arise. The amount of work being done in a bureaucracy can be small compared to the number of participants.

In figure~\ref{fig:growth_of_bureaucracy}a, the ratio of workers to participants is 3:3=1. When a manager is added, the ratio is 3:4 = 0.75 (a lower value means less work per employee is being accomplished). Adding a second team with dedicated management puts the ratio at 6:9=0.66. Finally, with a dedicated administrative team (e.g., payroll, hiring, facilities), the ratio is 6:13=0.46.

Adding more people to a hierarchical organization results in diminishing returns for time spent on the central work since bureaucrats invest time maintaining the hierarchy and administrative processes.

    \begin{figure}
        \centering
        \includegraphics[width=1\textwidth]{images/growth-of-bureaucracy.pdf}
        \caption{Growth of an organization as complexity and workload increase. a: three technical workers (green dots); b: administrative functions are delegated to an administrative assistant or manager (blue square); c: as complexity and scale grows, a new technical team (orange dots) is brought in which has their own manager; d: administrative work is delegated to a dedicated team of non-technical staff (purple diamonds).}
        \label{fig:growth_of_bureaucracy}
    \end{figure}

\ \\

\textbf{Each person in a bureaucracy has multiple roles}.\\
To limit the expansion of the bureaucracy. The multiple roles come from multiple relationships necessary to span an organization larger than \href{https://en.wikipedia.org/wiki/Dunbar\%27s_number}{Dunbar's number}. 
\ \\

\textbf{Everything is more complicated}. \\
Bureaucracy administers access to resources, whether tangible items (water, air) or expertise. There's friction for tangible resources because the easiest solution is you get the resource (without an intermediary). There's friction for expertise because the expert understands things you do not. 

\ \\

\textbf{Bureaucracy is inefficient and wasteful.}\\
The inefficiency of an organization is not just due to a breakdown of communication among its many members. Different individuals in different roles faced distinct incentives and have a different view of what problem is most relevant. There is a mixture of competition and cooperation internally for resources.


Can you explain a Pareto frontier for multiple objectives like speed, cost, accuracy?

Because decisions by bureaucrats are subjective, there is significant risk of being wrong or being called out by others as being wrong. Therefore, a motive to cover your ass for decisions made. Unnecessary work is carried out. 

Different participants have different motives, and the aggregation needed for coordination is inefficient regardless of what metric efficiency is measured against.

Another example: Inefficiency of changing the requirements on a project partway through. If an objective quantitative measure were available, the return on investment could be determined. 

Efficiency is typically assessed from the perspective of a serial process -- a single worker could accomplish this task faster, so why involve 10 people and get a slower, more burdensome result? The Mythical Man Month manifests \href{https://en.wikipedia.org/wiki/Amdahl\%27s_law}{Amdahl's law}. Dividing a task among 10 people does not make the task 10x quicker. It's not valid to expect that a task which takes 1 person 1 hour to perfectly scale so 10 people should be able to accomplish 10 results in 1 hour.

Inefficiency of process is similar. I need to file a request to replace the lightbulb in my office. Compared to the serial ``I'll just replace the lightbulb myself by running to the hardware store and purchasing one."

The excess resources characterized as waste are similarly driven by the serial-based assessment of what one person would require to carry out the task. 

\ \\

% https://graphthinking.blogspot.com/2020/07/scope-creep-is-experienced-differently.html
\textbf{Scope creep.} \\
The customer wants more, and the exploration of what's possible is exciting (a positive experience).

The creator hears more work. This implies a few trade-off options, all of which are negative for one or both parties.
\begin{itemize}
    \item sticking with the original terms (telling the customer "no", which is negative for both parties)
    \item re-negotiation for additional compensation (a burden to both parties)
    \item the programmer doing more for the same pay, which means less money per effort (yielding a less happy programmer)
    \item decreasing existing efforts to fit the additional new requirements (yielding a less happy customer)
    \item even if the programmer is getting paid by the hour, additional work means the end product will be delayed to accommodate additional features (yielding a less happy customer)
\end{itemize}

\ \\

% https://graphthinking.blogspot.com/2020/09/why-migrating-from-current-to-new.html
\textbf{Migrating technologies.} \\
The person implementing the transition has to be educated in both the old and new technology. 
the legacy code has to be migrated to the new implementation
convincing stakeholders; may require synchronization
difficulty scales with the number of stakeholders   % section
  \clearpage


%\chapter{Bureaucracies are made of Humans\label{b_made_of_humans}}
% instead of presenting specific lessons, here I'm going to present insider perspectives on the fundamentals described in previous chapters. 

%This chapter provides an insider's perspective of working in a bureaucracy. What are all the aspects to consider?
% unordered essays to be clustered later

  \chapter{Getting started in a Bureaucracy\label{sec:getting-started}}
    This section focuses on a bureaucrat's initial exposure to bureaucracy. 
  
    \subsection{What to Read\label{sec:to_read}}

The content of this book is independent of the specific role you as a bureaucrat will be playing. I assume you are striving to be a good person, and that you have relevant education/experience for the position in the organization.

The reading lists 
\href{https://github.com/LappleApple/awesome-leading-and-managing}{Leading and managing}, 
\href{https://github.com/kdeldycke/awesome-engineering-team-management}{Engineering team management}, 
\href{https://github.com/pdfernhout/High-Performance-Organizations-Reading-List}{High Performance Organizations}, and 
\href{https://github.com/ankitjaininfo/awesome-managers}{Managers}
are outside the scope of a guidebook for bureaucracy. 
%Bureaucracy involves hierarchy, so understanding the mechanics and nuances of hierarchical roles is useful. 


%Books on management might not seem applicable if you are not already in a management role.
Reading management books is helpful regardless of whether you are a manager, want to become a manager, or if you never want to become a manager. In all three cases, gaining awareness of best practices and alternative approaches benefits you as a bureaucrat and the people you coordinate with -- bosses, peer bureaucrats, subordinates, and subjects.

    
%See also 
% "Making of a Manager" has Chapter 7 (pages 161-187) on hiring



\section{Hiring into a Bureaucracy\label{sec:hiring}}


Hiring shapes the culture of an organization and determines what is feasible, both by the skills of those hired and how well new hires integrate with the existing organization. 

Hiring can be an expensive process, both in time and money spent by the organization to recruit and filter applicants, as well as the emotional investment of candidates. %, and the burden of reviewing applicants. 
% TODO: SO WHAT? What's the consequence?
As a candidate, recognize the risk the organization is taking. 
As a reviewer of candidates, hiring is an investment in the future of the organization. %Even if your bureaucratic role does not require an exclusive focus on hiring, there is value in studying this domain. 

%When hiring into an organization 
Organizations are not reflective of the society they operate within. As a candidate you should expect a different set of norms. As a bureaucrat reviewing applicants you will see behaviors not consistent with the norms of your organization. 
Reasons for the dissonance of norms in society versus the norms of the organization include incentives and constraints specific to the organization. 
There is also a \href{https://en.wikipedia.org/wiki/Selection_bias}{selection bias} in
\index{Wikipedia!selection bias@\href{https://en.wikipedia.org/wiki/Selection_bias}{selection bias}}
the people who join a bureaucratic organization, and there is a selection bias to who stays in a bureaucratic environment. 
% TODO: SO WHAT? What's the consequence?
The people in a bureaucratic organization are not representative of the society that the organization serves. Tactics that work inside a bureaucratic organization may not apply in society, and tactics that are useful outside an organization may not work on bureaucrats. 


% https://leadership.garden/onboarding-engineers/
% https://news.ycombinator.com/item?id=30810786

There are specific attributes that make a candidate more likely to be successful in a bureaucracy, regardless of the role they are being hired into. 
In addition to role-specific skills, hire candidates with the ability to reflect on 
thinking (\href{https://en.wikipedia.org/wiki/Metacognition}{metacognition}),
\index{Wikipedia!metacognition@\href{https://en.wikipedia.org/wiki/Metacognition}{metacognition}}\iftoggle{WPinmargin}{\marginpar{$>$Wikipedia: Metacognition}}{}
navigating complex social scenarios with competing interests, and intrinsic motivation are useful . Before becoming a candidate bureaucrat you can develop these traits. 

% https://graphthinking.blogspot.com/2021/04/screening-for-metacognition-in-job.html

% https://graphthinking.blogspot.com/2021/07/screening-for-intellectual-empathy-in.html

% https://graphthinking.blogspot.com/2021/04/questions-to-ask-interviewer-when.html


    \subsection{Tips for Getting Started}

% https://graphthinking.blogspot.com/2019/11/formal-education-for-domain-exposure.html
You may have graduated from school with a degree in a specific subject. Then you go to work and expect the position you're in to apply the domain-specific education you have at work. That work produces an outcome and gives meaning to the struggle of educational process (time, money, effort). When you are not allowed to practice what you went to school for the disconnect may feel dissatisfying.

An alternative attitude is that you went to school to learn how to learn. A degree represents your ability to learn and persist, rather than a domain specialization. Then, when you get to work, you are presented with new opportunities to learn other domains. 

\ \\

Sometimes your role as a bureaucrat may be poorly defined. The lack of direction or objectives may be feel disappointing or frustrating. The positive perspective is that you get to define your tasks and determine what would be useful or interesting. 

Two mentalities that I find useful are enabling others and working myself out of the job. My goal is to facilitate success rather than being integral to the process. Enabling other people to be successful can apply to customers, fellow bureaucrats, or management. Not everyone matters equally, so figuring out the relevance of helping specific people impacts your decision. 

Another tactic I use is to look around, look backwards, and look forward. What are the challenges my team or my organization faces? Then review previous attempts to address the challenges. Why did previous efforts fail to remedy the situation? Once I have a less naive view and more context, what is my vision of the desired scenario? Without this vision there is a danger of doing work that keeps me busy but yields no progress. 

The tactics of identifying challenges and learning the history of a challenge often relies on social interactions. A bureaucratic organization rarely maintains a searchable written record of decisions and failed efforts. If that content is available, it may lack the specificity available from oral narrative by people who were present. When investigating issues in a bureaucracy, I end discovery-oriented conversations with the question, ``Who else should I talk to about this issue?'' Independent of the content or the quality of the content, this question is the entry point to a large second-order social network of expertise that is often separate from the hierarchical chain of command. If multiple separate discussions lead to the same person, then that person is key. 

Documenting your process as you proceed helps both your self-reflection and your ability to measure progress. Written documentation of your status can be used as proactive updates to your management and provides a record when someone asks retroactively why you took certain actions.

Bureaucratic challenges that do not violate the constraints of Nature (such as conservation of momentum, conservation of mass, speed of light, the \href{https://en.wikipedia.org/wiki/Uncertainty_principle}{Heisenberg Uncertainty principle}) are solvable. There are typically political, social, personality, security, and technical aspects that need to be accounted for. Laws can be changed, regulations altered, people replaced, and norms updated. Constraints that appear to be logically inconsistent may be re-framed to be resolved. Narrowing or broadening the scope may reset the context for a challenge.




    \subsection{Setting Your Emotional State}

As a member of a bureaucratic organization, you may feel frustrated by the bureaucracy, the happiness of success, excited by possibilities, or the fear of uncertainty of operating in a complicated environment. Your emotional state, whether set by your work or set by events outside the bureaucracy, impacts your productivity and the well-being of your coworkers. 

You may see emotions as getting in the way of productivity and creating problems, or emotions as unprofessional. This view is not shared by everyone, and there are real benefits to concepts like passion, enthusiasm, excitement, and joy. Yes, joy can be felt by bureaucrats. 
    \section{Learning from Failure\label{sec:learn-from-failure}}

This section on individual failure is separate from
%\ifsectionref
%section~\ref{sec:org-failure-and-success} on 
%\fi
\hyperref[sec:org-failure-and-success]{how an organization characterizes failure}\iftoggle{haspagenumbers}{ (see page~\pageref{sec:org-failure-and-success}).}{.}

\ \\

Making progress in a bureaucracy is not a linear sequence of steps. Ideally there is an ongoing cycle of trying something, not succeeding, and then applying what you learn towards the next try. This idealized ``\href{https://en.wikipedia.org/wiki/Fail-fast\%23Business}{fail fast}'' 
\index{Wikipedia!\href{https://en.wikipedia.org/wiki/Fail-fast\%23Business}{fail fast}}\iftoggle{WPinmargin}{\marginpar{$>$Wikipedia: Fail-fast}}{}
cycle does not occur naturally -- the participants have to either have intrinsic motivation or external incentives to make progress. Failing as a path for learning is justified after having exhausted less expensive alternative education opportunities like reading books, going to school, and talking with experts. 

The recognition of failure depends on a clear measurement. Do the participants know the measure, and can they make the measurement? Having a defined measure of failure and regularly making the measurement relies on having an understanding of the expectations for the situation. Expectations are assumptions about the future.

Assumptions about the future can be categorized as
\begin{itemize}
    \item Opinions. For example, the interpretation of policy. Policy interpretation can be bent or exceptions can be made. 
    \item Guesses; more formally interpolation -- ``Given multiple options, X seems most likely."
    \item Experience; more formally extrapolation -- ``Last time X happened, so next time ...''
\end{itemize}

Failure with respect to expectations implies thinking about the future so that you can measure your progress (or failure). That is a subset of planning. You can fail if you don't plan, and you can fail if you do plan. This does not indicate that planning is irrelevant. Plans can be overly detailed and prescriptive, or plans can be inadequately specified; both are unhelpful.

If your measurements indicate failure, this could be due to an invalid assumption (with a valid objective) or you may have selected a bad objective. Declaring failure means re-evaluating assumptions, resetting objectives, or giving up on the concept and doing something else.

\ \\

There are ways to avoid failing: by not setting an objective, or setting an objective that isn't measurable. You can avoid the appearance of failure by not telling anyone the objective. Extrinsic accountability for failure requires other people to have awareness. If these tactics are being employed, it may be because the person (or
%\ifsectionref
%, in section~\ref{sec:org-failure-and-success}, 
%\fi
\marginpar{See page~\pageref{sec:org-failure-and-success}.}
\hyperref[sec:org-failure-and-success]{the organization}) doesn't want to share how they've failed until they have a success. 

Understanding rationales like the failure avoidance mentality is critical if you want to avoid feeling baffled the first time you encounter it. By learning these bureaucratic anti-patterns, you can inoculate yourself emotionally and cognitively. You now have a chance to think ahead before encountering the explanation.

A counterargument when you hear that someone is avoiding failure is to explain that noticing failure is essential to improving. Detecting indicators of failure, addressing the cause, and then changing is the process of improvement. When someone wants to avoid failure, I don't understand how they expect to improve. 


    \section{Success (and failure) in Your Organization\label{sec:org-failure-and-success}}

This section on how a team or organization conceptualizes failure is separate from %section~\ref{sec:learn-from-failure}
\hyperref[sec:learn-from-failure]{learning from failure} 
\marginpar{See page~\pageref{sec:learn-from-failure}.}  
as an individual bureaucrat. 

Understanding the ways your organization succeeds or fails matters to you because it affects how you are evaluated and how you feel about your work. Success and failure are measured against a quantitative objective. 


Bureaucratic organizations have varying definitions of success and failure because organizations are made up of individual bureaucrats with a multitude of motives and weak feedback mechanisms. There are a few macro-level patterns for organizations. Each pattern has an affect on the individual bureaucrats operating within the organization. 

\ \\

A bureaucratic organization that provides infrastructure service (such as a public water utility, an electrical power company, or an organization's internal computer support) has a specific failure mode. When infrastructure is unavailable or degraded, that's a failure. Success means satisfactory operation at a minimal cost. An organization operating in this context is motivated to avoid failure. Success is merely a lack of failure. An effort to improve processes or otherwise innovate will face the inertia of a culture organized around avoiding failure. 

\ \\

A bureaucratic organization that claims ``wins'' but no failures (because either more or different work is needed) is difficult to be productive in when the successful outcome has ill-defined value. The success can be counted (number of wins), but the relative importance of the success is unclear.

Examples of win-focused organizations include law enforcement, research teams, and teams with a responsibility for innovation. 

For an organization focused on ``wins,'' funding depends on how convincing these success narratives are. This biases the organization to focus on investments most likely to get more funding. When the funding for an organization does not come from consumers of products or policies created by the organization, there is a disconnect. The ``wins'' may not be helpful to the consumer.


% NOT INCLUDED -- need an example to make this more concrete
%A team that works to improve the productivity of other teams may have neither successes nor failures. Money and time are invested with the hope that results are useful to someone at some future time. If the team didn't exist no one would be significantly affected. 
    \section{Promotion\label{sec:promotion}}

Bureaucrats get promoted for a variety of reasons: due to their technical competence, a role needs to be filled, to motivate retention, or for their managerial ability.  Not being promoted can feel like a slight against any of those reasons. 

In this section, folk wisdom illustrates the frustration bureaucrats have with promotion. The purpose of you familiarizing yourself with these sentiments is to understand that the failures of promotion are not specific to the team you are on or the organization you are a member of. The deficiencies are generic to hierarchy. With this insight, you can choose a more emotionally-healthy response to the situation. Process Empathy is intended to produce emotional resilience based on an improved understanding of the situation you are in.

\subsection*{Folk Wisdom about Promotion}

The 
\href{https://en.wikipedia.org/wiki/Dilbert_principle}{Dilbert principle}
\index{Wikipedia!Dilbert principle@\href{https://en.wikipedia.org/wiki/Dilbert_principle}{Dilbert principle}}
\marginpar{$>>$ Folk Wisdom} 
\index{folk wisdom!Dilbert principle@\href{https://en.wikipedia.org/wiki/Dilbert_principle}{Dilbert principle}}
says organizations
``Systematically promote incompetent employees to management to get them out of the workflow.''
A similar-in-spirit observation is 
\href{https://en.wikipedia.org/wiki/Putt\%27s_Law_and_the_Successful_Technocrat}{Putt's Law}
\index{Wikipedia!Putt's Law@\href{https://en.wikipedia.org/wiki/Putt\%27s_Law_and_the_Successful_Technocrat}{Putt's Law}}
\marginpar{$>>$ Folk Wisdom}
\index{folk wisdom!Putt's Law@\href{https://en.wikipedia.org/wiki/Putt\%27s_Law_and_the_Successful_Technocrat}{Putt's Law}}
which says,
``\sout{Technology}[Bureaucracy] is dominated by two types of people, those who understand what they do not manage and those who manage what they do not understand."
Those two observations are post-hoc rationalizations for managers lacking management skills and lacking technical skills. 

For bureaucrats who don't seem effective in their role, try to determine whether there is a lack of training, lack of awareness, or lack of interest. Training can take the form of reading books or attending classes. Lack of awareness can be addressed with measurements of turnover rate and surveying team members. 




\href{https://en.wikipedia.org/wiki/Peter_principle}{Peter principle}~\cite{1970_Peter}:
\index{Wikipedia!Peter principle@\href{https://en.wikipedia.org/wiki/Peter_principle}{Peter principle}}
\marginpar{$>>$ Folk Wisdom} 
\index{folk wisdom!Peter principle@\href{https://en.wikipedia.org/wiki/Peter_principle}{Peter principle}}
``People in a hierarchy tend to rise to `a level of respective incompetence': employees are promoted based on their success in previous jobs until they reach a level at which they are no longer competent, as skills in one job do not necessarily translate to another.''
% https://graphthinking.blogspot.com/2021/03/a-tension-between-working-at-discount.html
This incompetence might manifest as recurring failure for a person in a role they are inadequate for. More likely, your incompetence will blind you to your failures.

To avoid falling victim to the Peter principle, do not seek promotion until you have a track record of performing at the next level. The downside is that you are working at a discounted rate since you are providing value beyond what is expected of your grade. (That's Dilemma~\ref{table:dilemma-personal-work-extra-or-work-as-expected}\iftoggle{haspagenumbers}{ on page~\pageref{table:dilemma-personal-work-extra-or-work-as-expected}.)}{.)} 


% https://graphthinking.blogspot.com/2021/04/notes-from-peter-principle-by-peter-and.html
The Peter principle has an element of truth, especially when a person starts a new role. It takes time to find your footing in a new job. However, you can learn their role and become competent. %Also, the Peter principle relies on the simplification of a single dimension of intelligence.  \href{https://en.wikipedia.org/wiki/Theory_of_multiple_intelligences}{Gardner's theory of multiple intelligences}.
%\index{Wikipedia!Theory of multiple intelligences@\href{https://en.wikipedia.org/wiki/Theory_of_multiple_intelligences}{Theory of multiple intelligences}}
%\iftoggle{WPinmargin}{\marginpar{$>$Wikipedia: Theory of multiple intelligences}}{}


\subsection*{Designing Incentives Determines Values}
Promotion is critical when designing incentives for behavior. Promotion is challenging to engineer because people a driven by various motives like money, status, and authority. 


Promotion typically focuses on the past success of a person. This de-emphasizes teamwork and lessons learned from failures. 
Promotion of a person (rather than the team) results in hero culture.


%Lowering the value of learning from failures increases risk aversion. 
The risk tolerance of a bureaucratic organization is driven in part by tensions in the promotion processes: an interest in success (more success is better) and avoidance of failure (organizations don't promote for failing). This imbalance can lead participants to seek ownership of efforts that are likely to succeed, to seek easier successes, and to steal credit for outcomes created by other people. 

%\subsection*{Promotion and Risk}

The promotion process within an organization is coupled to the idea of innovation, even if your promotion is at odds with taking risks. 


    %As a bureaucrat interested in improving your organization, exit interviews can help identify weaknesses. 


% https://news.ycombinator.com/item?id=30910387
    \subsubsection*{Organization chart orientation
\label{sec:org-chart-orientation}}

A common method of describing relations within the bureaucracy is the organization chart (commonly the ``\gls{org chart}"). \iftoggle{glossaryinmargin}{\marginpar{[Glossary]}}{}%
Normally the Chief Executive Officer (CEO) is at the top of the chart, middle management is in the middle, and managed employees are at the bottom. See Figure~\ref{fig:org_chart_orientation_ceo-at-top}\iftoggle{haspagenumbers}{ on page~\pageref{fig:org_chart_orientation_ceo-at-top}.}{.}

Artifacts like org charts subtly convey an organization's culture. 
% What's the point of this section? Is there a consequence, or is this just an observation?
There are emotional connotations to alternative layouts. You can alter expected relations (culture and norms) by playing with the orientation of the org chart.
Org chart orientation can be overanalyzed, so the exploration in this section is limited.

The point of thinking about org chart orientation is to frame how you perceive your supervisors, peers, and the bureaucrats you manage. Notice that the framing is embedded in the words -- prefixes super (over) and sub (under). 
These concepts inform what you expect from relations.
Do I seek support or direction and guidance from my boss? What do I expect from my boss? My peers? The people I oversee? What do I expect to provide them?

%\begin{itemize}
%\item 
%\end{itemize}

The relative orientation of the \href{https://en.wikipedia.org/wiki/Chief_executive_officer}{CEO} 
\index{Wikipedia!Chief Executive Officer@\href{https://en.wikipedia.org/wiki/Chief_executive_officer}{Chief Executive Officer}}\iftoggle{WPinmargin}{\marginpar{$>$Wikipedia: CEO}}{}
to the workers sets expectations for relations. 
Options for orientation are the conventional CEO at the top
(Figure~\ref{fig:org_chart_orientation_ceo-at-top}), 
CEO at the bottom (Figure~\ref{fig:org_chart_orientation_ceo-at-bottom}),
CEO on the right (Figure~\ref{fig:org_chart_orientation_ceo-leads}),
CEO on the left (Figure~\ref{fig:org_chart_orientation_ceo-follows}),
CEO as the center of a star 
(for example, the diagram for the \href{https://en.wikipedia.org/wiki/File:League_of_Nations_Organization.png}{League of Nations} in 1930.)
\index{Wikipedia!League of Nations diagram@\href{https://en.wikipedia.org/wiki/File:League_of_Nations_Organization.png}{League of Nations diagram}}

\begin{figure}
\begin{center}
\includegraphics[width=1\textwidth]{images/org-chart-orientation-ceo-at-top.pdf}
\end{center}
\caption{Standard orientation. The role with the most responsibility and authority is at the top. Left-right ordering is intended to be irrelevant in this view, though left-to-right reading order emphasizes importance.}
\label{fig:org_chart_orientation_ceo-at-top}
\end{figure}

\begin{figure}
\begin{center}
\includegraphics[width=1\textwidth]{images/org-chart-orientation-ceo-at-bottom.pdf}
\end{center}
\caption{Flipping the orientation presents a more realistic view of the CEO's responsibility. The crushing burden of servant leadership is clear. Left-right ordering is intended to be irrelevant in this view.}
\label{fig:org_chart_orientation_ceo-at-bottom}
\end{figure}

\begin{figure}
\begin{center}
\includegraphics[width=0.7\textwidth]{images/org-chart-orientation-ceo-leads.pdf}
\end{center}
\caption{Conventionally time flows from left (old) to right (new), so in this graph the CEO leads the charge into the unknown. Is the CEO dragging workers forward, or are the workers pushing the CEO? The top-to-bottom order can be read as importance. }
\label{fig:org_chart_orientation_ceo-leads}
\end{figure}

\begin{figure}
\begin{center}
\includegraphics[width=0.7\textwidth]{images/org-chart-orientation-workers-lead.pdf}
\end{center}
\caption{The ``chariot view'' with the CEO in the chariot and the workers out front. Workers are in the future; the CEO is in the past operating on old information. As with Figure~\ref{fig:org_chart_orientation_ceo-leads}, top-to-bottom ordering can be read as importance. }
\label{fig:org_chart_orientation_ceo-follows}
\end{figure}

\begin{figure}
\begin{center}
\includegraphics[width=0.8\textwidth]{images/org_chart_wedding_cake_dependencies_-_manufacturing.pdf}
\end{center}
\caption{An internal-customer-oriented view rather than a reporting-oriented view. The center of the bullseye is the team that generates the value that is the focus of the business or the organization.
The outer rings support teams that exist in the inner rings. The diagram is specific to an organization's domain. This visualization identifies which teams are the customers of which other teams in an organization.}
\label{fig:org_chart_wedding_cake_manufacturing}
\end{figure}



%extension of 
% \href{https://en.wikipedia.org/wiki/Conway\%27s_law}{Conway's law}: seating chart reflects org chart
    \subsection{Professional Training, aka Continuing Education}

There is often a disconnect between the results formal education and the specific needs of the organization. Professional training is intended to fill that gap.  
  \clearpage
  \chapter{You are an Individual in an Organization\label{sec:individual-in-org}}
  
    This section focuses on an individual bureaucrat operating within an organization. What actions can you take?  This is a prelude to section~\ref{sec:working-with-other-bureaucrats} on working with others.
    
    Understanding your options as a bureaucrat for action matters because trying to withdraw from bureaucracy isn't harmless. There is no neutral member of a bureaucracy. Either you are contributing to the organization, or you are sapping resources from the organization. This is because the organization operates in a zero-sum game for resource allocation.
    \subsection{Ideas for Innovation within a Bureaucracy\label{sec:innovation}}


The innovation lifecycle in a bureaucracy is
\begin{enumerate}
    \item You observe problems and challenges in your environment. This manifests as complaints from both the people directly harmed and observers who see inefficiency.
    \item You share ideas for innovation and gets feedback. Build a coalition of people willing to fight for the idea on your behalf
    \marginpar{[Tag] Actionable Advice}
    so that when you're not present, the idea is still proceeding towards implementation.  That puts the threshold at ``so important other people are willing to pause whatever they were working on and take up your cause.''
    \item Implementing these suggestions require either a change to existing processes or new processes or an investment of work. These changes may not succeed -- there's risk. Your idea could fail because it's not a good idea. It could also fail because someone doesn't like you, or the idea doesn't account for some dependency you weren't aware of, or it might conflict with other changes in progress.
    \item If you do decide to invest effort, the activity takes you away from your current work. Implementing the change might involve skills you don't have; learning those skills takes time. Carrying out the activity with new skills increases the likelihood of novice mistakes.
    \item If someone else implements the idea they get the credit for having done the work.
    \item If the idea saves money or time, there's typically no monetary reward. Recognition and benefit to your reputation isn't required as part of the change process. 
\end{enumerate}

There are many barriers in that lifecycle. The problem has to be observable to someone willing to invest effort in change. That person has to build a coalition of stakeholders. If the person isn't negatively harmed, that person may also lack clear benefit from resolving inefficiency. 

In addition to the work of implementing the change, there is an administrative overhead of documenting the reason for the change. 
Subjective decisions mean choices have to be defensible. 
The need for defensible justifications result in conservative decisions and risk aversion and a decrease of motive for innovation. 



These barriers lead to external observers to conclude something like the following simplification:
\begin{quote}
Bureaucracy destroys initiative. There is little that bureaucrats hate more than innovation, especially innovation that produces better results than the old routines.
\marginpar{[Tag] Folk wisdom}
Improvements always make those at the top of the heap look inept. Who enjoys appearing inept?\footnote{Frank Herbert (1987). ``Heretics of Dune'', page 201, Penguin}
% https://www.azquotes.com/quote/453163
\end{quote}


    \subsection*{Changing your Environment}

% https://quoteinvestigator.com/2017/05/23/culture-eats/
%\begin{quote}
%    Culture Eats Strategy for breakfast
%\end{quote}
%\href{https://en.wikipedia.org/wiki/Peter_Drucker}{Peter Drucker}

Changing the bureaucratic organization you are a member of is difficult. That applies regardless of your motive for change and the thing to be changed. Regardless of motive, there are tactics applicable to any change. 
The two major categories of change are 
\marginpar{See page~\pageref{sec:change-a-process}.}%
\hyperref[sec:change-a-process]{processes} (e.g., technical solutions, changing policy)
%\ifsectionref
%described in section~\ref{sec:change-a-process} 
%\fi
and changing the \gls{culture} of the members of an organization.
%\iftoggle{glossaryinmargin}{\marginpar{[Glossary]}}{}%

\ \\

\iftoggle{glossarysubstitutionworks}{\Gls{culture}}{Culture} is the set of norms that determine acceptable behavior for members of the organization. Norms are the expectations of your fellow bureaucrats. Each member of the organization can model desired behavior and inflict punishment for deviation from norms. The culture of the organization is practiced by members. 

% https://graphthinking.blogspot.com/2021/01/why-active-shaping-of-culture-is.html
The expectations of an organization should be stated explicitly so that members cannot use the defense that they were not aware. If the stated norms conflict with what is experienced by the organization's members and left unaddressed, the hypocrisy decimates trust. 

Besides modeling desired norms in the organization, there are a few more tactics you can take. 
\marginpar{$>>$ Actionable Advice}
\index{actionable advice}
You can reinforce examples of desired behavior with public praise and meaningful rewards. You can counsel people who go against the culture and punish instances of undesirable norms. 

Altering the culture of an organization means altering member behavior and expectation. This is accomplished by adjusting incentives to reinforce the desired behavior, punishing undesirable behavior, and providing resources to make the intended outcomes easier. 

Although aspects of culture may be problematic, thinking of cultural change as problem-solving can be a misleading framing. Rather than solve problems think of change as addressing challenges. %The word ``problem" implies a solution. 
Altering the culture often introduces novel challenges.
%The problem-solution framing is ina since . 
Even though you don't solve a challenge, you can still get a sense of emotional reward from having improved the situation. 

\ \\

Before investing effort in change, reflect on the emotional and career impact of your effort. Will the work bring you joy, or is your emotional reward tied to the outcome? If the outcome is the only source of reward, you'll either be disappointed in failure or you will be more conservative (intending to increase the likelihood of success). The work of creating positive change can be joyful: you can form relationships and focus on learning. 

One of the challenges of change in a bureaucratic organization is figuring out what is worth changing. There are many opportunities for improvement. Talk to stakeholders and enumerate their challenges.
\marginpar{$>>$ Actionable Advice}%
\index{actionable advice}%
From the myriad stories, identify recurring themes and root causes. Typically bureaucratic challenges intertwine technical, budget, staffing, historical context, personalities of participants, and organizational politics. 

Select an aspect that you have some competence in and would facilitate personal growth. If the challenge is big (many people involved, takes a long time, costs a lot), can it be broken down into incremental challenges? Distinguish wide-scale changes you envision from local changes over which you have more influence.

Having surveyed challenges worth addressing and having picked an issue with a high return on investment, the guidance of
\href{https://en.wikipedia.org/wiki/G._K._Chesterton\%23Chesterton\%27s_fence}{Chesterton's fence}
\index{Wikipedia!Chesterton's fence@\href{https://en.wikipedia.org/wiki/G._K._Chesterton\%23Chesterton\%27s_fence}{Chesterton's fence}}%
\marginpar{$>>$ Folk Wisdom}%
\index{folk wisdom!Chesterton's fence@\href{https://en.wikipedia.org/wiki/G._K._Chesterton\%23Chesterton\%27s_fence}{Chesterton's fence}}%
\iftoggle{haspagenumbers}{(page~\pageref{concept:chestertons_fence})}{section~\ref{concept:chestertons_fence}} 
is that you shouldn't change something until you understand why things are the way they are. Figuring out the history of an issue requires work when there is a lack of historical records.

Talking to stakeholders in the survey phase may yield tidbits of history. Why have previous efforts been rejected or not taken effect? Your review of written documentation of the history should be augmented by seeking out oral folklore. 

If your idea has a novel justification for success, identify all possible methods of progress and then rank the tasks. If you can't do all the work, who can you collaborate with? Can you delegate to people with skills you lack, or will you have to gain new skills? 
If no one has the relevant skills, becoming the widely-acknowledged expert improves your influence. 



    % management isn't leadership
    \subsection*{Creating Change in the Organization\label{sec:creating-change}}

If the organization you are in has no problems or challenges, this section can be skipped. For bureaucrats in organizations that do have issues, this part of the guidebook provides points to ponder independent of the specific problem.

As a bureaucrat, you have unique insight on the challenges the organization faces, and you have unique leverage to alter the situation.  While you could proceed haphazardly, an effective bureaucrat has vision. The vision is a story of what-could-be. This vision can be broken out as goals: specific measurable outcomes. Each goal requires contingency plans for achieving the goal. Plans are marked milestones which indicate whether the plan is proceeding successfully. 

Perspectives to consider when assessing change include what the situation is, what the situation could be, and what the situation looks like for different stakeholders. Each person has different information, a different way of responding to issues, and shifting intents.

People who you depend on who have conflicting visions or no vision, confounding your ability to improve the organization. There are different views on whether something is actually a problem, distinct characterizations of an issue, and competing priorities.

A trade-off to consider is that having niche impact is easier than broad change (the \hyperref[table:scope-broad-vs-narrow]{Dilemma of Scope of Effect}). There's also a trade-off of the quick fix versus more robust solutions (the \hyperref[table:quick-methodical]{Dilemma of Speed and Accuracy}).


%Determine social/political/technical impediments. 

\ \\
% https://graphthinking.blogspot.com/2018/08/potential-paths-when-faced-with.html
\textit{Tip}: Understand and empathize with the fear of change. \\
When faced with a challenge, the options are to take action or to not take action. If you take action, the options are failure, success, or iteration. 
The choice of no action is attractive if you are not suffering. Delay minimizes work, enables beneficial moves since the player can wait out other players, and there are other things to focus on.

Action incurs risk due to uncertainty and costs work to create change.

Fear of failure is justified if and only if the cost of the failure is greater than the value of lessons learned.

Fear of iteration is a fear that the process might be stopped prematurely -- before perfection.

\ \\
\textit{Tip}: Before starting a new effort, check to see whether this has been tackled before.
\marginpar{[Tag] Actionable Advice}
Learn the history of the problem. Why hasn't this been solved?

\ \\
\textit{Tip}: Query your first and second order social network.
\marginpar{[Tag] Actionable Advice}

\ \\
\textit{Tip}: Get feedback early before polishing.
\marginpar{[Tag] Actionable Advice} 

\ \\
\textit{Tip}: Advertise the result.

\ \\
\textit{Tip}: Hear criticism and respond.

\ \\
\textit{Tip}: Leverage both social networks and bureaucratic processes. 

\ \\
\textit{Tip}: Identify sources of power (hierarchical positions and titles, social influence, reputation and credibility, buzzphrases or popular paradigms) and leverage them.

\ \\
\textit{Tip}: Professional respect (for what the other person knows) and professional curiosity (for what you don't know) \\
Example: Getting approval from multiple overseers in different hierarchies is hard. Often different objectives and incentives.


\ \\
% https://graphthinking.blogspot.com/2016/01/methodology-for-people-acting-as.html
\textit{Tip}: Use social recommendations to name relevant individuals.\\
Leverage the trust already in an existing social network by starting with ``Person A recommended I talk to you about X."


\ \\
% https://graphthinking.blogspot.com/2016/01/methodology-for-people-acting-as.html
\textit{Tip}: Sit in on meetings, listen to topics, see who is talking, who is attending. After the meeting talk to individuals about the meeting. Set up one-on-one informal discussions. Keep the first conversation  brief - 10 or 15 minutes. Your body language should indicate engagement and interest. ``Who else would you recommend talking to?" is the last question in the first conversation.

\ \\
\textit{Tip}: Consensus doesn't mean everyone agrees on the problem, the remedy, or the approach, or who's taking the action. Consensus in a bureaucracy means people aren't going to resist the change.
    \section{Measuring Bureaucratic Maturity of Bureaucrats}
% https://graphthinking.blogspot.com/2021/07/three-measures-of-bureaucratic-maturity.html

Bureaucratic maturity of bureaucrats in an organization can be broken into three behaviors: 

The first and most wide spread behavior is to see a problem and then complain about the situation. This indicates awareness of the environment but no creativity regarding what to do about the problems.

The second, less common behavior, is to see a problem and recognize the situation as an opportunity. This requires creativity and reframing. 

The third behavior is to see a problem and then nudge the situation towards a vision. The nudges could be through either direct action or by influencing others. 



An individual bureaucrat can show one or more of these behaviors. That is, being capable of the third behavior does not mean the person lacks the ability to complain on other topics. 

There is not a specific amount of experience within the organization needed to arrive at any of these three paradigms. A holistic understanding of the system certainly helps.

The ``vision" of the third behavior could be in the form of a long-term (temporally distant) outcome, or the vision could be of immediate multi-party cooperation. 
  \clearpage
  % person-to-person relations
  \section{Working with other Bureaucrats\label{sec:working-with-other-bureaucrats}}

Delegation of tasks
Asking for help
Seeking input
Offering to help
Offering input
Understand what someone else is priorities are and why those are priorities

These all directly impact your reputation; see \S~\ref{sec:reputation}.


enumerate tropes to figure how to respond

% https://graphthinking.blogspot.com/2021/10/why-i-dont-like-being-in-management-role.html
Solo work may be more emotionally rewarding due to fewer external constraints, but the cost is complexity and scope being limited to the skills of the individual. 

Working with others allows you to occasionally accomplish complex results beyond your own skills or your own bandwidth in spite of collaborators not being under your control. How? Through persuasion. 

The challenge of collaboration is to multiply productivity rather than merely sum the output of a set of individuals. 

inside an organization, cooperation/coordination is not held together by internal contracts or even service level agreements. What holds the organization together? Force of will of participants. 

two networks: formal an informal % section
    \subsection{Motivation of Bureaucrats\label{sec:motivations}}

Appreciating the diverse motives of the bureaucrats you work with is instructive. If you expect everyone to have the same motives as you then you will be surprised by the friction created by diverse motives. 

Motivations of participants are rarely ``how can I make the organization more successful" or even ``how can I sell/produce more product"? Usually motivation is based on personal success in various manifestations, which leads to emergent phenomena which appears confounding to observers outside the bureaucracy. 


% see https://en.wikipedia.org/wiki/Social_influence

Each bureaucrat has a motive, even the bureaucrats who do nothing. 
% https://graphthinking.blogspot.com/2020/02/there-is-no-idle-status-for-paid.html
In an organization where you are a paid bureaucrat, you are either actively working for improvement of the organization, or your existence is parasitic to the organization. There is no ``idle" status for paid employees in an organization with limited resources.

As an example motive for a bureaucrat, I want to avoid being too efficient such that I eliminate the need for my job, and not so inefficient that the organization fails and I lose my job. Increasing the efficiency of bureaucracy is good for the organization and the outcomes, but can be harmful to the bureaucrat's career.

Career stability within an organization is a benefit, and it can be leveraged to take more risk. However, it typically manifests as inaction by an employee. There's no harm to the employee in not taking action. If an employee doesn't do anything, nothing bad will happen to that employee. Career stability decreases extrinsic motivation.


Example motivations for bureaucrats: 
stability (aka job security, the comfort of a routine),
money (current pay or future earnings), 
travel, 
problem solving, 
status, 
exerting power or control, 
credibility of being associated with the organization (if the organization has a positive reputation), 
logistical convenience (``the office is near where I lived''), 
service to people the organization serves.



The consequence of diverse motives is that expecting bureaucratic organizations to be logical, fair, consistent, and efficient is unreasonable even when every participant wants those features. Each bureaucrat thinks, ``I am logical, fair, consistent, and efficient.'' Therefore each bureaucrat expects other bureaucrats to meet those same (unrealistic) standards. Next, anthropomorphize the team or organization and expect the group to meet those standards. 

Even if each bureaucrat were logical, fair, consistent, and efficient (they are not, and neither are you), each person has a different motivation. Each person wants to accomplish something different using their unique skills and referencing their own experiences. Compounding the confusion, each bureaucrat has to coordinate using communication that has latency and limited bandwidth and isn't precise.

An expectation that bureaucracy feels illogical, unfair, inconsistent, and inefficient is a useful baseline. Working against bureaucratic entropy yields improvements even though perfection is inaccessible.
    \section{Bureaucratic Tropes\label{sec:tropes}}

Any sufficiently large organization will have a diverse group of bureaucrats. From that diversity certain tropes emerge\footnote{J.~Duffy, personal communication, 2017}. The value of recognizing these recurring patterns is to leverage their strengths so that you and that person can be more effective. 

% TODO: a commenter at work provided this. Can I use it?
\begin{itemize}
    \item Excited intern or new employee. \\
    \textit{What you can do}: Their enthusiasm can be channeled to take on challenges other bureaucrats are scared of taking on for political reasons. Their naivete should be remedied by pairing with a mentor who can explain the organization's history and specific challenges.
    
    \item Retired in place. \\
    \textit{What you can do}: Recognizing this attitude allows you to temper your expectations about their productivity. 
    
    \item Font of institutional memory. Likes to explain how the current situation relates to his experience ten years ago. \\
    \textit{What you can do}: You can benefit from talking to this person about historical processes and what has been tried before. 
    
    \item ``Let them eat cake'' managers who are removed from challenges and deny the problem is a problem. \\
    \textit{What you can do}: It is worth figuring out whether these managers actually don't care, or if they are busy fighting battles you are not exposed to.
    
    \item ``Don't rock the boat'' people who tried changing the system, didn't succeed, and are emotionally beaten down such that they no longer try. \\
    \textit{What you can do}: These people are often creative and can be consulted before you tackle a challenge they have experience with. 
    
    \item The Patriot is willing to sacrifice for the organization. \\
    \textit{What you can do}: Educate them that there are alternative ways to be effective in a bureaucracy.
    
    \item Mismatched -- this person's skills and interests don't match their role, so they are despondent and frustrated. \\
    \textit{What you can do}: You can provide value to this person by networking with them and finding better opportunities.
    
    \item Friendly grouch: friendly and helpful but unhappy in their job due to bureaucracy and other hurdles. \\
    \textit{What you can do}: Provide a copy of this book to increase their effectiveness and improve their spirits. They will be grateful to you and the organization will benefit.
    
    \item Detail-oriented: ask them what time it is and they'll tell you how to build a clock. \\
    \textit{What you can do}: You can serve as their ambassador to other members of the team or organization.  
\end{itemize}


% see related but not relevant
% https://blog.codinghorror.com/which-online-discussion-archetype-are-you/
The above personality tropes apply to individuals. If your team works on multiple projects, you may identify project tropes:
\begin{itemize}
    \item We've always done it this way.\\
    \textit{Assessment}: The team is scared of change. The change may be known (and undesirable) or unknown. Change disrupts power.\\
    \textit{What you can do}: You can brainstorm with the team ways to improve incrementally. 

    \item Individual solutions to systemic challenges.\\
    \textit{Assessment}: Demonstrating small wins helps build credibility. \\
    \textit{What you can do}: Plan to ramp up  complexity (towards the systemic challenges) and \hyperref[sec:reputation]{spend the earned reputation}.
\marginpar{See page~\pageref{sec:reputation}.}
\end{itemize}
Team tropes can emerge from either the top-down directives made by supervisors or by the collective (in)actions of team members.


\ \\

% TRANSITION to solving_problems

Tropes as labeled patterns of behavior can decrease surprise when you first encounter them and can help you develop constructive responses. 

The last set of tropes is on the topic of bureaucratic problems. 
    \subsection*{How to Solve Problems}

%This list is somewhat in jest, somewhat as a source of bureaucratic patterns.

Not every bureaucrat responds to challenges the same way. Understanding the diversity of potential reactions (or inaction) before you encounter them decreases the surprise. 

% TODO: someone else at work wrote this? Seek permission
\begin{itemize}
    \item Solve the problem. This is the engineering approach.
    \item Solve a simpler problem. This is the Physics approach.
    \item Solve a more abstract problem. This is the Mathematical approach.
    \item Solve a different problem.
    \item Recognize and ignore the problem. 
    \item What problem?
    \item Assign someone else the problem.
    \item Redefine the problem so that it is already solved. ``That's not a problem, that's how it's supposed to work.''
    \item Apply for resources (staff, training, money) to solve the problem.
    \item Point to the problem as a justification for more resources.
    \item Give a talk about the problem.
    \item Write a paper about the problem.
    \item Teach a course on the problem.
    \item Talk to your therapist about the problem.
    \item Your manager declares the problem solved and gets a bonus.
    \item Reframe the problem as an opportunity. 
    \item Study the problem in the context of the broader system.
    \item Study the problem from other people's perspective.
    \item Explore the impact of the problem on the broader system.
    \item Develop a set of activities for role playing parts of the problem. 
    \item Solve part of the problem. Then another part.
    \item Research whether anyone else worked on the problem.
    \item Find other instances of the problem that were already solved.
    \item Find similar problems and identify the differences.
    \item Contract out the solution and select a bid.
    \item Explain why the problem does not need to be solved.
    \item Explain why the problem is beneficial. 
    \item Form a team and set a mission statement addressing the problem.
    \item Slow-roll the problem until you don't need to worry about the problem. 
\end{itemize}


Did you recognize any of those responses from your own past experiences? 
\marginpar{[Tag] Reflect and Discuss}
Were any unexpected?
    \section{Not my Job: Task Scope}

% https://graphthinking.blogspot.com/2021/05/not-my-job-task-scope-and-collaboration.html

If I ask someone for help on a task that benefits our team, that person might respond that the task is ``not my job" and time spent on tasks like that ``gets in the way" of their progress. The interesting (emotionally rewarding) work is not that task.

`\href{https://www.urbandictionary.com/define.php?term=scut}{Scut}' is medical slang for the non-clinical yet essential tasks that do not require a doctor's degree or expertise.
This is different from administrivia (administrative tasks).
%as ``scut'' would include taking out trash. 

% TODO: expand on
``Once I realized someone else has the same problem, I stopped worrying about it.'' Same as ``not my job."

The potential reasons for this reluctance to help include
\begin{itemize}
    \item Working on the task will not get them promoted.
    \item Their understanding of the scope of their job and their expertise does not include the task.
    \item The person don't know how to do the task and they don't want to learn.
\end{itemize}
In any case, progress for the individual is not aligned with the success of the team.
    \subsection{Building, Managing, and Spending Reputation\label{sec:reputation}}

Your personal reputation within the organization dramatically impacts your effectiveness.

Organizations have reputations externally. 
Internal-to-the-org there is cultural norms. 
% https://graphthinking.blogspot.com/2021/01/why-active-shaping-of-culture-is.html


% https://graphthinking.blogspot.com/2018/05/my-evolving-view-on-role-of-my.html


Building reputation through multiple small wins or larger risk on bigger bets


Relation between Reputation and Brand and Political Capital? Same thing?

% the following article is useless
% https://www.indeed.com/career-advice/career-development/build-a-reputation
% since it reduces to "be a good person"

How does an individual create and accumulate political capital? What does political capital mean with respect to teams?

% https://graphthinking.blogspot.com/2021/09/notes-from-class-on-being-politically.html

Reputation matters for influence. How other people perceive you impacts what you can accomplish and when people seek out your help or input.

Your reputation is actively changing based on your activities and communication -- both your communication and the stories others tell about you.

Neglecting to manage your reputation means you lose input to the stories others tell about you. Active management of your reputation requires engaging people and generating evidence. 

Reputation is set whenever and where ever you are observed, or artifacts are associated with you. What you wear, when you show up, how your emails appear, body posture in meetings. 


Reputation is perception of the person the bureaucrat is engaging with.  What does that person think of you?

Ideally that would be a function of their technical skill, ability to collaborate with other people, the strength of their network, creativity. None of this matters if the person you're engaging with doesn't know those things. 

Based on your reputation, what trust does that person have? 

To spend reputation is to bend the rules. 

Spending reputation means taking risk that involve other people

Build reputation by doing useful things that are visible to other people

    % https://graphthinking.blogspot.com/2021/07/patterns-anti-patterns-in-bureaucracy.html

% intra- and inter- team dynamics

\subsection{Teams are Subdivisions of an Organization}

\cite{2015_Katzenbach}

In the context of altering teams, there are a few major levers available: create new team, merge teams, dissolve a team. 


For a given set of teams, the lateral interactions are competitive or cooperative. Coordination is required (or conflict will occur) for money, staffing, and resources. Examples of resources include access to or control of data, computer equipment, hardware, floor space, prestige, products (output).

\subsubsection{Roles of Management versus Leadership}

Teams include managers and leaders. Those roles are not necessarily the same person. 

A manager's role involves time management, task tracking, employee evaluations, promotion, pay, requesting resources for team members, and hiring. 

A leader's role focuses on coordinating vision and principles. The vision and principles do not have to originate from the leader -- other team members can contribute ideas. The leader's responsibility is primarily social consensus. 
    \section{Intellectual Empathy\label{sec:intellectual-empathy}}

As a bureaucrat, you will work with people who are smarter on a specific topic than you are. That superior knowledge might come from experience or formal education. 

The person with superior knowledge may be able to convey their knowledge to you, or they may struggle to. That struggle could be due to an inability to break down complex insights, or they don't recognize the mismatch between what they know and what you understand. 
This same scenario can be reversed -- you might be more knowledgeable on a given topic than some other bureaucrat you work with. In either situation, both parties have responsibilities for \gls{intellectual empathy}. The more knowledgeable person needs to understand what they can build on when talking with the less knowledgeable person, and the less knowledgeable person has to be clear about what makes sense and what does not. 

You can start conversations with people who are new to you by asking about their professional background -- their education, past work experience, and other aspects that would allow you to build on their existing knowledge. 
\marginpar{$>>$ Actionable Advice}
\index{actionable advice}
This allows you to tailor the conversation to the gap of knowledge that exists when discussing a challenge. 

During a conversation you can monitor body language for reactions like a squint, a head tilt, an eyebrow raise or an eyebrow furrowed. 
\marginpar{$>>$ Actionable Advice}
\index{actionable advice}
Those are common reactions when information doesn't conform to what the listener expected. 

You can ask about the other person's knowledge of the topic. The phrasing ``Do you know about (name of topic)?'' is less effective than ``What do you know about (name of topic)?'' 
\marginpar{$>>$ Actionable Advice}
\index{actionable advice}
Respondents might claim experience or knowledge of a topic, but probing for the difference between ``I've heard those words before'' versus ``I was part of a team that implemented (topic)'' versus ``I came up with the original idea for (topic)'' is helpful. 

Sometimes the person with less knowledge on a topic feels uncomfortable with indicating their confusion, or sometimes the person doesn't realize they are missing the main point. The more knowledgeable person can check in with the listener to confirm. ``What did you get from what I just talked about?'' is a better confirmation than ``You understood what I meant, right?'' The person with less knowledge can proactively check their knowledge by providing a read-back: ``Here's what I understood you to mean:~...'' These check-ins, regardless of who initiates, should be relatively frequent when exploring complex topics. 

Admitting ignorance can be intimidating and may feel like you risk losing the other person's respect. A gentler way of saying ``I'm lost'' is ``I was following up until the point where you described~...'' Intellectual empathy is the responsibility of both parties in a conversation. 


    \subsection{Leveraging Expertise}

Experts are cautious to include caveats and limitations

Translation process is necessary from "here's what's known in general" + "situation specific facts" to "plan of (in)action" and what's realistic with what confidence. 

Experts straying outside their expertise are dangerous because the boundary may not be clear to the audience.

non-experts think and speak less precisely


    \section{Does Anyone Want to Volunteer?}

% https://graphthinking.blogspot.com/2020/06/what-to-ask-instead-of-does-anyone-want.html

The facilitator of a meeting asking, ``Does anyone want to volunteer for this task?" to a group will most often be met with silence. 
The lack of engagement can be addressed by not asking for volunteers. The outcome can be higher engagement with the following tactics:
\begin{enumerate}
    \item Ask each person whether they are attending to observe or to volunteer.
    \item Ask each person how much time they are willing to volunteer. Responses could be zero house, one hour (non-recurring), one hour per month, or something else.
    \item Ask each person what their goals in the interaction are. It is usually generic (``I want the group to succeed.") but it can be narrow, in which case that gives you something to focus on.
    \item Ask each person what they are good at and what skills they have. Are they good with personal interaction? Writing? Computers? Coordinating? Logistics? Fundraising? Making phone calls?
    \item From the inventory of tasks, are there any that fit both the skills and time? Can the task be scoped to fit the time? If there are multiple candidate tasks for a volunteer, let them pick the task. (If no existing task aligns with their skills, do not create work to be assigned.)
    \item If there are other volunteers working on the same tasks, put them in contact with the other participants.
    \item Give the volunteer a deadline for the task. Your deadline should not be arbitrary -- it should be based on the task dependencies. 
    \item Confirm that the volunteer is willing to commit the time to complete the tasks by the deadline.
    \item Schedule a check-in with the volunteer before the task deadline to review progress.
    \item To create accountability among multiple volunteers, hold a group review to explain how each participant's work contributes to the goals and how work done by one person enables the next task done by someone else. (Enumerate the dependency graph; include deadlines.)
\end{enumerate}

  \clearpage
  \chapter{Communication within a Bureaucracy\label{sec:communication-within-bureaucracy}}
{\footnotesize Back to the \hyperref[sec:toc]{Main Table of Contents}}
\minitoc


\section{Communication is Critical for Bureaucracy}
\Gls{bureaucracy} 
\marginpar{[Tag] Glossary}
is a system of distributed knowledge and distributed decision making. Distributed decision making relies on effective communication. Bureaucrats communicate by writing and speaking, both of convey incomplete information and are imprecise. The inability to be comprehensive and precise in communication means iteration is key when establishing a shared mental model of the situation and plan.


The dependence on iterative interactions means relationships matter for communication. 
\href{https://en.wikipedia.org/wiki/Metcalfe\%27s_law}{Metcalfe's law} 
\index{Wikipedia!\href{https://en.wikipedia.org/wiki/Metcalfe\%27s_law}{Metcalfe's law}}
says
\marginpar{[Tag] Folk Wisdom}
\index{folk wisdom!\href{https://en.wikipedia.org/wiki/Metcalfe\%27s_law}{Metcalfe's law}}
the value of an organization is proportional to the square of the number of people interacting in the organization. A team of 5 people is not just the aggregated skills of five individuals -- there's the synergy arising from ten bilateral relationships. Broadening your network of collaborators broadens your potential effectiveness and that of your fellow bureaucrats. 

% https://graphthinking.blogspot.com/2018/08/confusion-leads-to-confusion.html
Suppose that you don't invest in relationships with your fellow bureaucrats. Then you'll be less likely to know what is happening in your organization, and you won't understand why certain decisions were made. Because of your confusion you might focus on your own self interests since that is what you have some control over and insight on. That doesn't leverage the potential synergy of collaboration, but at least you look busy. 

The analysis above applies to each person in the bureaucratic organization.
% A Nash equilibrium?
The consequence is that the organization's goals are not achieved when bureaucrats act independently. 
There is a way out of this: broadcast your intent to other people, be transparent with decisions, and share the results of your activities. That decreases the confusion other people may have about your actions and can improve effectiveness. You can make yourself and the organization more effective even when your investment in communication is not reciprocated.

\ \\

Information within a bureaucracy shapes the relationships among bureaucrats. You have a choice: you can decide to not share information, you can be transparent, and you can be vulnerable.  
Not sharing information can harm or protect relationships. 
Being transparent conveys ``here's what is happening." Being vulnerable expands that to ``here's why that's happening or what might happen."
Each of those options informs how your coworkers, peers, management, and team members interpret their perception of your intent. Your communication shapes your reputation.


Another way communication alters bureaucratic organizations is when actions informed by locally-available information produce \hyperref[sec:failure-to-comm]{suboptimal results for participants}, as illustrated by the
\href{https://en.wikipedia.org/wiki/Prisoner\%27s\_dilemma}{Prisoner's dilemma}.
\index{Wikipedia!\href{https://en.wikipedia.org/wiki/Prisoner\%27s\_dilemma}{Prisoner's dilemma}}
The concept of ``locally-available information'' refers to communication among bureaucrats that facilitates delegation of work, allocation of resources, relationship creation/maintenance, and carrying out processes. Each bureaucrat benefits from gathering information that shapes their actions, but communication has a cost: time spent building and verifying consensus delays action. There is an opportunity cost when time is spent talking with people.

% https://graphthinking.blogspot.com/2021/09/why-is-everything-so-hard-in-large.html

 % section

    \subsection{Why of Communication within a Bureaucracy Matters}

Taking action informed by locally-available information can produce suboptimal results for participants, as illustrated by the \href{https://en.wikipedia.org/wiki/Prisoner\%27s\_dilemma}{Prisoner's dilemma} (see \S\ref{sec:failure-to-comm}). Communication among bureaucrats facilitates delegation of work, allocation of resources, relationship creation/maintenance, and carrying out processes. Communication has a cost for participants: time spent building and verifying consensus delays action. There is an opportunity cost when time is spent talking with people.

% https://graphthinking.blogspot.com/2021/09/why-is-everything-so-hard-in-large.html



\subsubsection{Communication Reflects and Shapes Your Thinking}

Can you turn complaints into impact statements? \marginpar{[Tag] Actionable Advice}

\textit{Negative observation}: ``Logging into my computer takes a long time.''\\
Positive statement with explanation of impact: ``If I were able to log into my computer more quickly, then I could accomplish more tasks.''

\textit{Negative observation}: ``The team I need support from doesn't provide a ticket tracking system.''\\
\textit{Positive statement with explanation of impact}: ``If the service team I need support from provided a ticket tracking system, then I would be able to know the status of my request.''

\textit{Negative observation}: ``When I submit a ticket for a support request to the team, I don't have visibility on the status of the request.''\\
\textit{Positive statement with explanation of impact}: ``If the service team I need support from provided visibility in the ticket tracking, then I could proceed with other tasks without worrying about the status of my request.''


Here's a technique to transform generalizations to action: \marginpar{[Tag] Actionable Advice}
\begin{enumerate}
    \item This organization does not like blueberry pie.
    \item No one in this organization likes blueberry pie.
    \item I don't know of anyone in the organization who likes blueberry pie.
    \item I like blueberry pie, how would I find someone else in the org who likes blueberry pie.
\end{enumerate}
In the first statement we personify an organization. Organizations cannot like or dislike tangible items, so this is a meaningless statement. The second statement is slightly more precise (individual people can have likes and dislikes) though still a generalization and therefore likely inaccurate. The source of the generalization might be your lack of relationships with every member of the organization. That leads to the fourth statement which is correct, precise, and curious.  % subsection
    \section{Failure to Communicate\label{sec:failure-to-comm}}

Communication is critical in bureaucracy because bureaucracy is a system of distributed knowledge and distributed decision making. When less effective communication occurs, individual bureaucrats are less able to rely on the knowledge of other experts, and they have to make decisions with less consensus. 

\subsection*{Slowing Communication\label{sec:slowing-communication}}

In an ideal scenario there would be no delay associated with communication -- you would get the information you need when you need it. In practice, there are various causes for why your progress is blocked when you depend on other people. 

Tactics that delay communication within a bureaucratic organization are stonewalling, slow-rolling, bikeshedding, and red herrings. By learning these concepts you will be better able to identify and then respond to their use.

\index{responsiveness!stonewalling}
\iftoggle{glossaryinmargin}{\marginpar{[Glossary]}}{}
\iftoggle{glossarysubstitutionworks}{\Gls{stonewalling}}{Stonewalling} 
is when the recipient of a request or question  doesn't reply. There may be \hyperref[sec:email-responsiveness]{legitimate reasons for the lack of response}.
\iftoggle{haspagenumbers}{\marginpar{See page~\pageref{sec:email-responsiveness}.}}{}%
The person may be busy and didn't see your message, or they did see your message but didn't have a chance to reply yet because a response to you is lower priority than other tasks they have. You cannot differentiate those reasonable causes from when the recipient doesn't want to enable you to proceed. They may disagree with your objective and see silence as \href{https://en.wikipedia.org/wiki/Passive-aggressive_behavior}{less confrontational}
\index{Wikipedia!passive-aggressive behavior@\href{https://en.wikipedia.org/wiki/Passive-aggressive_behavior}{passive-aggressive behavior}}\iftoggle{WPinmargin}{\marginpar{$>$Wikipedia: Passive-aggressive behavior}}{}
than explicit rejection. 

One way of circumventing stonewalling is to ask if the respondent is opposed to your idea. 
\marginpar{$>>$ Actionable Advice}%
\index{actionable advice}%
Then a lack of response indicates no opposition. This tactic applies if you are confident the recipient will read or hear the message.

An unintentional source of stonewalling is when you ask on the wrong channel. Sending an email may result in what appears to be stonewalling if the person relies on chat messages or the phone. The solution for this
\marginpar{$>>$ Actionable Advice}%
\index{actionable advice}%
is feasible, but action is required by the person who doesn't respond. The person who only uses certain channels should explicitly indicate that. An automatic out-of-office email that says, ``Contact me by phone'' tells the sender the \hyperref[sec:communication-preferences]{preferred channel}. 
\marginpar{See page~\pageref{sec:communication-preferences}.}

From the view of the person doing the stonewalling, if you need time to think or gather information before responding, 
\marginpar{$>>$ Actionable Advice}%
\index{actionable advice}%
tell the person who sent a request that you acknowledge their message and will follow up in more detail later (with a specific timeline). While better than no response, this leads to the next challenge.

\index{responsiveness!slow-rolling}
\iftoggle{glossaryinmargin}{\marginpar{[Glossary]}}{}
\iftoggle{glossarysubstitutionworks}{\Gls{slow-rolling}}{Slow-rolling} 
is when you get a response to your request or question, but the response isn't helpful. Progress is delayed because you have to iterate to get an answer. There are valid reasons for a slow-roll and there are uncool reasons for a slow-roll response. Perhaps the person wants to acknowledge your request but doesn't currently have time to provide a full explanation. The person may need to gather more information for the complete response. Or the person is \href{https://en.wikipedia.org/wiki/Passive-aggressive_behavior}{passive-aggressive}
\index{Wikipedia!passive-aggressive behavior@\href{https://en.wikipedia.org/wiki/Passive-aggressive_behavior}{passive-aggressive behavior}}
\iftoggle{WPinmargin}{\marginpar{$>$Wikipedia: Passive-aggressive behavior}}{}%
and may understand your question but does not want to enable your progress. 

The reason for a slow-roll should be made explicit by the respondent, 
\marginpar{$>>$ Actionable Advice}%
\index{actionable advice}%
and a timeline for a complete response is helpful. 


\label{concept:bikeshedding}
\index{responsiveness!bikeshedding}
\iftoggle{glossaryinmargin}{\marginpar{[Glossary]}}{}
\iftoggle{glossarysubstitutionworks}{\Gls{bikeshedding}}{Bikeshedding} 
is when the recipient of a question or request 
\href{https://en.wikipedia.org/wiki/Law_of_triviality}{focuses on unimportant details relative to the primary topic}. 
\index{Wikipedia!Law of triviality@\href{https://en.wikipedia.org/wiki/Law_of_triviality}{Law of triviality}}\iftoggle{WPinmargin}{\marginpar{$>$Wikipedia: Law of Triviality]}}{}%
Whether this behavior is intentional or not, the best response is to refocus the conversation on the core issue. The amount of time allocated for various topics should be proportional to their consequence. 

\index{responsiveness!red herring}
A \gls{red herring}\iftoggle{glossaryinmargin}{\marginpar{[Glossary]}}{}
response is misleading, whether intentional or not. The respondent provides what looks like a reasonable answer but results in unproductive work. Occasionally there is a coincidental benefit of discovering something unexpected, but that wasn't the respondent's intent. 


Even though bikeshedding and a red herring can be a \href{https://en.wikipedia.org/wiki/Dark_pattern}{dark pattern},% 
\index{Wikipedia!dark pattern@\href{https://en.wikipedia.org/wiki/Dark_pattern}{dark pattern}}
that may not be the intent of the speaker or author. Perception matters more than intent.
\marginpar{$>>$ Mantra}%
\index{mantra!perception matters more than intent}%
Process Empathy applies both when you are sharing information (how could the information be perceived?) and when you are receiving information (what was the author's intent?).

\subsection*{Decreased Effectiveness in Communication and Some Remedies}

\textit{Challenge}: The \href{https://en.wikipedia.org/wiki/Allen_curve}{Allen curve} 
\index{Wikipedia!\href{https://en.wikipedia.org/wiki/Allen_curve}{Allen curve}}
is 
\marginpar{[Tag] Folk Wisdom}
\index{folk wisdom!\href{https://en.wikipedia.org/wiki/Allen_curve}{Allen curve}}
an ``exponential drop in frequency of communication between engineers as the distance between them increases.'

\index{mantra!presence creates priority}
Presence creates priority - go to their desk. I once needed some data from a coworker. After trying email and phone calls, I ended up flying across the country. Once I arrived the person was able to provide the data in a few hours.

Merely sitting next to a coworker, even with no official purpose of interaction, results in spontaneous informal discussions. See the discussion of 
\hyperref[sec:prisoner-exchange]{Prisoner exchange} on 
page~\pageref{sec:prisoner-exchange}.

Take advantage of the \href{https://en.wikipedia.org/wiki/Allen_curve}{Allen curve} 
\index{Wikipedia!\href{https://en.wikipedia.org/wiki/Allen_curve}{Allen curve}}
by implementing the Inverse Conway Maneuver. If you know what structure is needed for a product, then design the placement of your team to reflect that.

\ \\
\textit{Challenge}: \href{https://en.wikipedia.org/wiki/Wiio\%27s_laws}{Wiio's law}: 
\index{Wikipedia!\href{https://en.wikipedia.org/wiki/Wiio\%27s_laws}{Wiio's law}}
\marginpar{[Tag] Folk Wisdom}
\index{folk wisdom!\href{https://en.wikipedia.org/wiki/Wiio\%27s_laws}{Wiio's laws}}
``Communication usually fails, except by accident.''\\
This pessimistic take is similar to \href{https://en.wikipedia.org/wiki/Murphy\%27s_law}{Murphy's law}
\index{Wikipedia!\href{https://en.wikipedia.org/wiki/Murphy\%27s_law}{Murphy's law}}
and is indicative of the level of ongoing investment needed for effective communication. 

\ \\
\textit{Challenge}: Periodic status reports sent up the chain of command get sanitized so that only good news is shared. This impedes risk analysis. \\
If your reports are getting sanitized, as for a copy of the sanitized version. If you are the person consolidating and aggregating reports, aim for conciseness rather than good news. 

\ \\
\textit{Challenge}: Decisions by bureaucrats high in the \href{https://en.wikipedia.org/wiki/Command_hierarchy}{chain of command}
\index{Wikipedia!\href{https://en.wikipedia.org/wiki/Command_hierarchy}{command hierarchy}}
are not pushed down the chain. \\
You can request management provide a summary of their activities.

\ \\
% Role of assumptions 
\textit{Challenge}: To assume makes an ass out of you and me, 
\marginpar{[Tag] Folk Wisdom}
\index{folk wisdom!To assume makes an ass out of you and me}
yet assumptions are necessary to making progress in communication.\\ 
This dissonance can be addressed by looking for sources of difference and then talking about them. For example, when I talk with someone for the first time I ask what their educational background it. If they have a different degree than mine, I can tune my language to their academic training. I can make my story more relatable. 

Another technique for detecting differences is to ask about the person's previous experience. What did they work on previously in this organization? What were their jobs before joining this organization? This backstory can provide context for decisions that need to be made in the current context. 

\subsection*{Leveling up Your Communication}

There are levels of enlightenment for bureaucrats. If you know the progression, you can step up to the next level more easily.
\begin{enumerate}
    \item I feel bad.
    \item (complaint) I can't do what I want.
    \item (complaint) I can't do what I want in the way I want.
    \item (complaint) There is a problem.
    \item I have a solution.
    \item I have an implementation.
    \item I tried but my solution didn't work.
    \item (burnout) Life sucks but I get a pay check.
    \item I quit (in hopes of being more successful somewhere else).
\end{enumerate}


\subsection*{Theory of What Could be Done and Why the Theories are Impractical}

% https://graphthinking.blogspot.com/2021/09/why-is-everything-so-hard-in-large.html 

To break the 
\href{https://en.wikipedia.org/wiki/Prisoner\%27s\_dilemma}{Prisoner's dilemma}, 
\index{Wikipedia!\href{https://en.wikipedia.org/wiki/Prisoner\%27s\_dilemma}{Prisoner's dilemma}}
options are 
\begin{itemize}
    \item Expose all participants to the consequence of outcomes. In practice this feels unfair to each participant because the outcome is partially attributable to other people involved in the process. Dividing responsibility limits exposure to consequences.
    \item Have all participants communicate. In practice communication takes time and skill. Not everyone is willing to invest since communication is not seen as ``doing the work.'' Accounting for the \href{https://en.wikipedia.org/wiki/Allen\_curve}{Allen curve}
    \index{Wikipedia!\href{https://en.wikipedia.org/wiki/Allen\_curve}{Allen curve}}
    takes effort. The time needed to arrive through consensus at an optimal approach for a given situation may exceed time available for solving the problem.
    \item Limit everything to what can be accomplished by one person. This hero-based approach is limited to the attention-bandwidth of the person and their skills. As the complexity increases the necessary skills increase and the number of candidate heros decreases. Large organizations accomplish complicated tasks by leveraging diverse skillsets of teams of bureaucrats.

\end{itemize}

%  How does communication among individuals fail?

 
  % subsection
    \subsection{Communication Preferences}
% why this matters: source of friction among bureaucrats
% benefit to reader: understanding why people are different

There are typically multiple communication channels available between bureaucrats in an organization. Phone, email, chat, in-person, voicemail, websites, video calls, calendars, message boards. 

Do you know which channel is preferred for what purpose? Can you list your own preferences, and the preferences of other stakeholders?

Do you keep your calendar up-to-date? Or is your calendar irrelevant? Is your calendar visible to your coworkers?

Text chat is useful for asynchronous interrupts. Text chats are better than stopping by in person or calling on the phone since in-person drive-bys and phone calls interrupt whatever I'm thinking about or discussing. The purpose of the text chat is either a reminder or seeking a convenient time to talk. 

Email is useful for notifications or questions or setting up logistics. If I don't respond to a question or request for action within 2 days, please send a reminder. 

Video calls are my preferred method for group meetings. Group meetings via video call scheduled in advance on my calendar are best. Impromptu video calls are acceptable but as interruptive as a phone call. 

In-person discussions are my preferred method for one-on-one discussions. Stopping by my desk  without an appointment interrupts whatever I'm thinking about. If a calendar invitation is inconvenient, a text chat confirming my availability prior to the in-person discussion is helpful. 

Phone is useful when the caller is unable to be present in-person. Phone, being voice-only, is inferior to video calls. Unscheduled phone calls interrupt whatever I'm thinking about. I've intentionally disabled voicemail. 

Status updates in electronic issue trackers or wikis or forums are best (rather than email-only). Notifying in text chat or email that a wiki page or issue is updated is helpful but not mandatory. 

Given all those caveats, I prefer communication in any form at any time over surprises. 

If you're not sure whether you should communicate, default to communicating. If you're not sure whether communication would interrupt something, check my calendar and then communicate. If it's urgent or blocking progress, ignore my calendar.  % subsection
    \section{Communicating Professional Sense of Vulnerability\label{sec:professional-vulnerability}}

Emotionally vulnerable communication is personal and is a form of intimacy. You are exposing personal issues. This personal emotional vulnerability deepens relationships. The risk is that information could be used to manipulate or harm. Whether you choose to form social friendships with your coworkers is up to you, but isn't a requirement for an effective bureaucrat. There is a more relevant practice of vulnerability.

Professional vulnerability is about going beyond transparency regarding bureaucratic issues. Transparency is about what is happening, while professional vulnerability is about why something is happening. For example, professional topics of conversation include processes and incentives. Bureaucratic processes may be intended to be impersonal, but the consequences are felt by participants.


Discussions of internal intrigues of an organization are a form of gossip among professionals. \href{https://en.wikipedia.org/wiki/Gossip}{Gossip} 
\index{Wikipedia!\href{https://en.wikipedia.org/wiki/Gossip}{Gossip}}\iftoggle{WPinmargin}{\marginpar{[Wikipedia] Gossip}}{}
can be constructive (finding aspects to remedy),  lead to insights, and shapes cultural expectations within the bureaucracy. As with personal gossip, professional gossip can risk harm if used against other organization members. 

As with personal vulnerability, professional vulnerability involves learning who to share what information with and when. In interacting with a person who is new to you, you can experiment by being professionally vulnerable and see whether they reciprocate or at least explore the topic with you. Being open, direct, and curious can help the person you're talking with feel more comfortable. Shared introspection is the objective. 

As an example, you can say something like, ``Perhaps the reason behind (observation) is (reason 1) or (reason 2).'' That gives the other person a chance to brainstorm with me without committing to a position. If the other person is unwilling to explore in-depth, reverting to safer topics is easy.

Being vulnerable does not mean the other person will reciprocate. That is a risk on your part.  % subsection
    % https://graphthinking.blogspot.com/2020/05/invisible-bureaucracy.html

\subsection*{Social and Bureaucratic interactions\label{sec:socializing}}

Change in a bureaucracy can apply to processes and people, but a more amorphous concept is changing the culture of a team or organization. What is meant by ``culture'' usually refers to norms -- the expectations of behavior that individuals hold to. That definition of culture is generic; what is meant within a bureaucratic context requires jargon for specific expectations.

To evaluate expectations, we start by introducing categories of interactions. 
Interactions among members of an organization are either a social interaction or a bureaucratic interaction. 

As examples of each of these,
\begin{itemize}
\item \textit{Social interaction example}: ``Did you see the game on TV last night? Our team did fantastic, right? I wanted to get tickets for the game, but they were sold out."
\item \textit{Bureaucratic interaction example}: ``You'll need to get approval from Sue before presenting your idea to the board for their review. Then talk with Russ and get his thoughts about how to proceed."
\end{itemize}
Both social and bureaucratic interactions are vital to cohesion in an organization. 


Bureaucratic interaction can be broken into two subcategories: 
\gls{visible bureaucracy} \iftoggle{glossaryinmargin}{\marginpar{[Glossary]}}{}%
(procedures and processes are written down and can be discovered by stakeholders) and 
\gls{invisible bureaucracy} \iftoggle{glossaryinmargin}{\marginpar{[Glossary]}}{}%
(procedures and processes are known to some stakeholders and are conveyed verbally to some of the other stakeholders).

Invisible bureaucracy is akin to related topics outside the professional environment: invisible domestic work\footnote{Cleaning your living space, raising children, caring for pets; see~\cite{1987_Daniels}.} and invisible relationship work.\footnote{Consistent need to delegate, being curious without reciprocation.} The work associated with emotional cohesion, logistics, planning, scheduling, and communicating is hard to quantify so it does not get counted.


The relevance of this jargon is to break down the components of an organization's ``culture" experienced by participants.
When someone in the organization advocates for changing the culture, which expectations are they specifically referring to? Invisible bureaucracy is the hardest to alter because it is undocumented and not counted.
%The ratio of social relationship to visible bureaucracy to invisible bureaucracy is a characterization of the culture. There are norms associated with each of these three categories.

Processes default to invisible bureaucracy because creating and maintaining documentation requires work. Making the documentation discoverable requires work.
%, and because some processes are embarrassingly inefficient. 
To make invisible bureaucracy visible, document the work and enable other people to find the documentation.
 % subsection
    \section{Communication Tips}

The role of verbal communication is critical for bureaucrats. 
There is a lot of advice on effective communication (enunciate, speak loud enough to be heard, be humble, be curious), so the advice below is highlighted because of prevalence in bureaucratic organizations. 
The following is generic to interactions outside of bureaucracy. 

For general writing tips, see Strunk's and White's Elements of Style and other resources \footnote{\href{https://www.youtube.com/watch?v=vtIzMaLkCaM}{LEADERSHIP LAB: The Craft of Writing Effectively} and \href{https://www.youtube.com/watch?v=aFwVf5a3pZM}{LEADERSHIP LAB: Writing Beyond the Academy}}.

\subsection*{Tip: Not all interaction challenges are communication problems.}
Sometimes an inability to discuss ideas is not a communication problem but a psychological deficit of personality. Distinguishing ``I'm an ineffective communicator'' from ``the person I'm talking with doesn't communicate effectively'' from ``that person has a diagnosed psychological reason they are unable to communicate'' is tough for those of us who are not psychologists or psychiatrists. 

%how to measure effectiveness: The waste or inefficiency in a bureaucracy is a measure of the lack of coordination or inconsistent decision making among the members

\subsection*{Tip: Avoid relying on stereotypes.}
Within an organization different teams may build up reputations for certain behaviors, or there may be significant events that the team is associated with. 
When interacting with members of a team that has a reputation, avoid relying on that stereotype or event as an opening for discussion. 
You're speaking to an individual, so address that person's behavior.



\subsection*{Tip: Avoid questions that have a binary response\label{sec:yes_no_questions}.}

Responding to a request with ``no'' is advantageous for the person replying to the question. There is less work required, less risk of failure, and better continuity. As an example of a poorly framed question, I could ask, ``Can I have a copy of the data you're using?'' The person I'm asking is less disrupted if they refuse to share. 

A more constructive phrasing is ``I need information on X to work on Y, and I think you have information about X. How can you help me get information on X?'' By clarifying my intent, I allow the person with the data to provide options I may not have considered.

Similarly, when I'm being asked for information, I try to learn the person's intent motivating the question. Sometimes the requester doesn't actually know what to ask for. Instead of ``no'' I try to figure out how to enable the person to be successful. 

\subsection*{Tip: Leverage the other person's experience while focusing on your own\label{sec:advice}}

Advice without context is less effective.\\
\textit{Bad}: ``Here is what I think you should do in that situation.''\\
\textit{Better}: ``Here is what I did in that situation.''

People usually find talking about themselves an easy topic if you are curious about their experiences. 
If you can learn the other person's background and history and motivations, you can weave that into the advice you provide. 
Tailoring your message increases the likelihood of implementation. 

\subsection*{Tip: Avoid Platitudes\label{sec:platitudes}.}
% https://graphthinking.blogspot.com/2017/10/why-platitudes-are-used.html
\href{https://en.wikipedia.org/wiki/Platitude}{Platitudes} are \gls{thought terminating}; the statement feels true and is resistant to debate. Platitudes capture a feeling with sufficient accuracy, but with imprecise language. As a result, there's no specific action.

Because platitudes result in a conclusion, the conversation participants may feel more bonded. However, that bond is shallow.

Example platitudes to avoid:
% https://graphthinking.blogspot.com/2017/02/phrases-which-serve-as-thinking-stoppers.html
% https://graphthinking.blogspot.com/2017/10/a-list-of-platitudes.html
\begin{itemize}
    \item pick your battles
    \item Some things you can't explain
    \item Your time will come.
    \item You can be anything that you want to be
    \item I just want to get through this day
    \item It is what it is
    \item I'm just one person
    \item That's that
    \item Life's not perfect
    \item Life's not fair
    \item There's only so much you can do about it
    \item What is meant to be will be
    \item It is God's will
    \item It's part of God's plan
\end{itemize}

If your goal is to understand a concept or issue deeply, you need to use precise language.

\subsection*{Tip: Strive to use Precise Language}

Imprecise language causes miscommunication. Intent is unclear, as is expected consequence.

If you have a specific definition for a word central to the topic of interest, ask your new collaborator for their definition. Do not expect others to share your definition even when there are established norms for the topic. 

Instead of asking a collaborator, ``Are you taking action on this topic?'' ask ``What actions are you taking on this topic?''

If someone claims, ``We plan to get to that action,'' ask for a timeline. A deadline can be for an artifact or a re-evaluation of the topic.

In the short-term imprecise language takes less work to create and can take less time to convey. 

The importance of precise language is proportional to the potential consequences of action/inaction/wrong action.
Precision also should be proportional to the complexity of the topic being discussed. 

\subsection*{Tip: Word is bond\label{sec:word-is-bond}}

Your communication (verbal, written) is your reputation. People rely on what you tell them even if there isn't legal recourse. Reliance on your word is why precision matters. 

Frustration and disappointment follow when you don't uphold your word, or others misinterpret your imprecise language, or you are misunderstood.

Communication implies responsibility for the content.  There is a corresponding accountability in the relationship between speaker and listener (or writer and reader).

\subsection*{Tip: Take care near the boundaries of knowledge}

Trying to find someone else's extent of knowledge is tricky -- they don't want to appear stupid. They may interpret the exploration at a trap. ``I don't know'' can be an embarrassing statement to make, even if you don't share their embarrassment. 

Knowing the limitations of your own knowledge and disclosing those boundaries to others is critical. Distinguish what you know from speculation about things you don't. 

\subsection*{Tip: Listen all the way to the last word of the speaker}

Formulating a response to the first part of an idea or a sentence is tempting. Waiting for the speaker to finish before thinking of how to response is courteous. Waiting creates a pause which makes you seem more thoughtful. 

This is complicated by the speakers who include long pauses for contemplation and then resume. 

\subsection*{Tip: Eliminate speaking over other people.\label{sec:crosstalk}}
% https://graphthinking.blogspot.com/2017/10/crosstalk.html

Crosstalk occurs two people who are communicating verbally experience interference from another audible conversation. That can occur because a third person is talking at one of the original two participants, or when four or more people are holding two separate conversations concurrently. 

Crosstalk in a bureaucracy is motivated by
limited time available to communicate. A meeting participant may feel inspired by something someone else said and want to interject. 
%Typically manifests as popcorn style stories based on experience. 
%Intended as wisdom for self-validation by others in our community. 
Crosstalk can indicate engagement and enthusiasm, or it can be due to the speaker wanting to dominate the topic through interruption. The likelihood of crosstalk is dependent on the level of aggressiveness of participants.
In either case (enthusiasm or power-seeking), the original speakers are disrespected. The original speaker may feel annoyed at being interrupted.



%Crosstalk has four roles and a minimum of two people participating: the discussion facilitator, the original speaker, the interrupter, and other meeting participants. 
%During crosstalk, the discussion facilitator loses control of the interaction to the interrupter.  
The audience is frustrated by the lack of clarity of where to focus. This distraction causes participants to lose of focus and productivity of the interaction decreases.

Bystander intervention for out-of-control meeting: raise your hand. \marginpar{[Tag] Actionable Advice} This non-verbally reverts focus back to the discussion facilitator. 

\subsection*{Tip: Account for Warnock's dilemma}
% https://graphthinking.blogspot.com/2018/09/dealing-with-warnocks-dilemma-in.html
\href{https://en.wikipedia.org/wiki/Warnock\%27s_dilemma}{Warnock's dilemma}
is the common experience of figuring out how to interpret not getting feedback. This is especially vital in meetings where the speaker or facilitator needs to gauge participant comprehension of delivered content. Simply asking ``Does anyone not understand what I just described?" is likely to get no response from attendees because individuals want to avoid looking stupid.

\ \\
\textit{Technique}: Pick an individual to provide a recap.\\
\textit{Technique}: Survey the audience using multiple choice to gauge understanding.\marginpar{[Tag] Actionable Advice}

\subsection*{Tip: Seek action with a deadline}

% https://graphthinking.blogspot.com/2017/11/collected-wisdom.html
When asking someone for help or input, specify a deadline for their response. \marginpar{[Tag] Actionable Advice}This helps the person prioritize their tasks.

\subsection*{Tip: Identify the cause of miscommunication}

% https://graphthinking.blogspot.com/2019/06/miscommunication-versus-inability-to.html
\begin{itemize}
    \item miscommunication the cause is often due to definitions of words used or differences in context. In this situation, additional time spent communicating and taking different approaches is sufficient to remedy the issue.
\item A speaker may be inarticulate. If the speaker is unable to coherently convey their internal experience to a listener, then the communication failure is of a distinct category. No amount of additional communication will lead to improved understanding on the part of the listener.
\item A speaker simply has nothing to say about a subject. Regardless of whether they are capable of articulating a concept, they may be unable to relate to the topic. Often a person in this situation still wants to participate, but they are unable to meaningfully contribute. 
\end{itemize}
% https://graphthinking.blogspot.com/2019/05/identifying-empty-talk.html
Empty talk is the use of words that are ill-defined, emotionally resonant, inactionable, and impersonal.

\subsection*{Tip: ask-tell-ask}

Collaborating with fellow bureaucrats who have expertise in areas you do not requires extra work. There may be differences in the words used to describe certain situations, more precision in wording that you're used to, or thinking about situations in ways you are not familiar with. In that context, to bridge the differences you can ask, tell, ask\footnote{\href{https://cepc.ucsf.edu/sites/cepc.ucsf.edu/files/Curriculum_sample_14-0602.pdf}{``The 10 Building Blocks of Primary Care: `Ask Tell Ask' Sample Curriculum''} and \href{https://www.the-hospitalist.org/hospitalist/article/125126/qi-initiatives/ask-tell-ask-simple-technique-can-help-hospitalists}{Ask-Tell-Ask: Simple Technique}}. 

The first step is asking what the other person's perspective is on the topic. This helps establish the appropriate level of nuance is and can tell you how that person frames the issue. The second step is to tell the person what you want to say. The ``tell'' step should leverage what you learned from the first ``ask'' step. Use phrasing that is consistent with what you just learned from the other person. The third step is to ask the person what they heard from you. If they are unable to tell you, you may need to refine your delivery. To improve the likelihood of success keep the content in the second step short. 

The ask-tell-ask technique can be used iteratively in the same conversation, especially in discussion complex topics with a new collaborator. \marginpar{[Tag] Actionable Advice}Using ask-tell-ask takes long than just telling but increases the effectiveness of the communication. You also get to learn more about the other person's perspective. 


\subsection*{Tip: Initial responsiveness and status updates}
In a bureaucracy requiring approval, or soliciting input, sometimes waiting can provide value to the person doing the waiting. The request may be overcome by events, or the person asking may remind which indicates priority.

\subsection*{Tip: Make deadlines explicit}

Typically requests have two deadlines. The first deadline is when a response is sought. The second deadline is when a response is no longer useful.  

As an example, suppose I am inviting people to a meeting. I send the invitation 5 days prior to the meeting and I want to know who is able to attend by 3 days before the meeting. The second deadline is the time of the meeting. Replies after that second deadline do not help me understand who is going to attend. 

\subsection*{Tip: Read each email/memo/report to determine the purpose }
% https://graphthinking.blogspot.com/2021/03/read-each-email-to-determine-purpose.html

\textit{Problematic behavior}: scan the text of a message, see if there is immediate action or response needed. If no action or response is needed, go to the next email. \\
 That does not work for emails that contain logistics associated with future events. 

Instead, consider possible intentions of the person writing the email. 

\textbf{Decision needed}. Typically includes context. \\
\textit{Action}: if the team maintains a decision log, update that.
Response is selection of a choice.

Tip: Instead of asking for a decision, ask for if the person is opposed.

Tip: instead of asking for a decision, ask for the go-ahead. This framing biases the respondent towards action (specifically approval) rather than thinking. 

\textbf{Situational awareness}.\\
\textit{Action}: Expected default is no action, but interject if there's an issue.


\textbf{Action or Tasking}.\\
\textit{Action}: Do something within some deadline

\textbf{Approval sought}.\\
\textit{Action}: Confirm or deny

\textbf{Feedback sought}.\\
\textit{Action}: Assessment of proposal


\textbf{Meeting logistics}. Can be an announcement (widely available), registration (limited attendance), or invitation (specific to you). Attendance is optional or require. \\
\textit{Action}: Create or update a calendar event
Response should restate the logistics (time/date/location/purpose) to confirm. 

\textbf{Brainstorming}\\
May provoke a response for building on an idea.
``For your situational awareness, no action needed." Notification of activity by someone else. Or change in plans. 
If needed, a correction to the described direction might trigger a response or even a meeting.

\textbf{Reference} e.g. describing a process or business workflow. Or a citation.\\
\textit{Action}: Copy process documentation to wiki. Copy citation to bibliography.
Acknowledgement response thanking the sender for the update/clarification.

\textbf{Setting a formal policy or issuing an informal edict}\\
\textit{Action}: move the policy/edict documentation to Confluence or Wiki
Acknowledgement response needed only if the edit is aimed at just me or the group I am leading

\textbf{Question}\\
If this is a recurring question, move to a ``Frequently Asked Questions" page on Confluence or Wiki.
Response needed that provides answer or seeks clarification.


Here I'm using ``action" to refer to activities outside the email channel. 

If I read email to figure out the purpose of the email, that will help me determine what action and response are relevant. 

Whether I am the only recipient or on of many receivers can change the intent of the email, and whether I'm in the ``to" or ``cc" field matters. Unfortunately, ``to" versus ``cc" are not reliable indicators since email senders do not reliably conform to the expected use. 



%This categorization of text within emails is a useful natural language processing challenge for machine learning. Currently a few email providers already do some of this with identifying meeting logistics, providing reminders to follow-up, and providing reply snippets. A browser plug-in that differentiates the various purposes of text could help readers determine relevant actions and responses. 

An email sent to multiple recipients may have different purposes for different readers. The reader's role or knowledge may factor into how they interpret the content. The inclusion or exclusion of recipients alters how the content is understood. 

\subsection*{Tip: Don't seek attribution for contributions; credit others\label{sec:credit-others}}

Give credit to others for good ideas and beneficial actions. Either they accept credit and you are seen as a contributor to their success, or they push back and you look generous. Credit is not a \href{https://en.wikipedia.org/wiki/Zero-sum_game}{zero sum game}.

\subsection*{Tip: Offer to take blame\label{sec:take-blame}}

Before an action commences, tell collaborators that you are willing to accept blame if something goes wrong. This alleviates their fear of risks.

\subsection*{Tip: Survey stakeholders}
% https://graphthinking.blogspot.com/2016/01/how-to-solve-and-not-solve-problems.html

Suppose you are a \href{http://www.peacecorps.gov/}{Peace Corps} worker in Africa. You show up and the village doesn't have easy access to clean water. Villagers walk a long ways in dangerous areas for dirty, unsafe water. This is a very obvious problem and all the villagers agree that they don't have good water and that this problem should be fixed.

Implementing the solution would take about a week - get the equipment to the village, drill a well, build a pump.

You could take additional time and involve the villagers in this project. They could participate in getting the equipment, which should lead to a sense of ownership.
But then when the equipment shows up, they don't take action to drill the well. If the well is drilled, it soon falls into disrepair and the villagers are back to doing things they way they used to. What happened?

The villagers don't see access to clean water as the most significant issue. You came in and imposed your view of what the problem is and how to fix it. When you impose your view of what the problem is, the solution won't be adopted by villagers because they don't prioritize it.
It is better to survey the community to see how they operate. What do they think the problems are?
Both leadership and the community members need to provide priorities.

This issue is exacerbated if you come to the village as a representative of a company providing wells. You are biased when you ask, ``Do you have any problems?"

Of course the villagers have water problems which could be fixed with better wells. However, when you get into the details of placing or improving a well, they lose interest. What the community really wants is free installation, zero maintenance, easy to use, and no operational costs. That would improve their life.

When you say there's cost (both initial investment of capital and then operations/maintenance) and a learning curve associated with the solution, then the user's interest wanes -- you are presenting another cost/benefit ratio for them to evaluate. Then they ask ``Can we get by without the well?" Yes, they don't need the well -- they've survived without it.

Novel solutions (in this example, drilling a well and installing a pump) have have barriers to adoption. Two barriers are the current priorities of the community and the incumbent solution/processes.

If there are problems with higher priority, the community will delay implementing your solution. That's fine if the higher-ranked priorities are bounded, but they are often not. An example of this is the following:
Suppose a person has three tasks, and you introduce a solution which is a fourth task.
If the first task is ``go from point A to point B," then that task will eventually be eliminated and there will be three remaining.
If the second task is ``secure your village," that is an unbounded task. The person won't get to or won't prioritize your low-ranked task.

How will your solution impact their higher-ranked priorities?

\ \\

\subsection*{Email Structure\label{sec:email-structure}}

% from https://graphthinking.blogspot.com/2021/10/structuring-email-content-for.html

If implementing all these tips sounds like a lot of work, that's because it is. Effective written communication requires intentional effort because of the lack of augmenting channels (compared to voice or video or in-person). 



Use consistent design and structure for your emails. Emails are part of your professional reputation.

Emails start with a greeting: Hi, Hello, Good morning, Good afternoon, Good evening. 
Email greetings include the name of the targeted recipient(s). 

Emails terminate with a professional closing, e.g., ``Kindly", ``Regards", etc

Emails contain a signature block with contact information -- phone number, normal hours of response, which timezone you're in if your team spans timezones, how long to wait for a response before asking again, which communication channel I prefer, etc.

\begin{figure}
\includegraphics[width=1\textwidth]{images/email_template.pdf}
\caption{Template for new email messages. Greeting has a space after the comma -- that is where the recipient's name will go. Signature block uses smaller after the name.}
\label{fig:email_template}
\end{figure}

Email signature blocks do not include unnecessary images, as that uses more storage for recipients. 
Email threads focused on a specific instance of a recurring event include the date (YYYY-MM-DD) in the subject line. 

Based on the purpose of the email, example key phrases for subject lines include: ``meeting notes" versus ``agenda" versus ``question about".

Revising an existing subject line can disrupt the ability of email software to thread conversations. However, sometimes the revision is worth breaking threading.

When replying to an ongoing thread, retain the original message as part of the thread to provide readers historical context.

When replying to threads with sensitive messages, sanitize the included content as appropriate by removing name or identifying details.

If an email contains multiple requests or questions, at the top of the email (after the greeting) explicitly say how many of each type. Then, in the body of the message, number them.

\begin{figure}
\includegraphics[width=1\textwidth]{images/email_two_questions.pdf}
\caption{Distinct items the recipient should address in a reply.}
\label{fig:email_two_questions}
\end{figure}

If an item corresponds to a requested action, separately highlight the action and indicate who is supposed to take the action and what the deadline for response is.

\begin{figure}
\includegraphics[width=1\textwidth]{images/email_meeting_notes.pdf}
\caption{Who has what action due when?}
\label{fig:email_meeting_notes}
\end{figure}

Computer commands should be distinct separate fixed-width font. This distinguishes the text from the rest of the narrative. 


\begin{figure}
\includegraphics[width=1\textwidth]{images/email_computer_font.pdf}
\caption{The computer commands use fixed width font. The author distinguished input from output through the use of bold vs non-bold respectively. The author highlighted the error message using red. Inline text like ``cat'' in the last line is also fixed width.}
\label{fig:email_computer_font}
\end{figure}

References to documents include a direct full path.

If referring to a previous separate thread, include the subject and the date+time that email was sent

For bullet points, explicitly specify that items are joined by one of the following: OR, XOR, AND

If you have an unordered list, explicitly state that order is irrelevant.

If you have a sequence of steps, number them appropriately and indicate which steps are required versus optional

Use visual sketches to illustrate concepts rather than always relying on text. Don't use pictures all the time, and don't have too many pictures in an email. 

Know how to both embed pictures inline and how to attach files and when to use which. 
Email replies should preferentially be at the top of the thread. 
If replying to multiple points in the previous email, embed replies inline, mark the distinction, and highlight the authorship. 

\begin{figure}
\includegraphics[width=1\textwidth]{images/email_reply.pdf}
\caption{Bob's reply to Sue's questions. The third question is not shown in this illustration.}
\label{fig:email_reply}
\end{figure}

If replying inline, explicitly state that at the top of the thread.

The email trilemma is to balance the amount of detail against providing sufficient context and being concise. 

If the email is longer than a paragraph, provide a \href{https://en.wikipedia.org/wiki/BLUF_(communication)}{B.L.U.F} or \href{https://en.wikipedia.org/wiki/Wikipedia:Too_long;_didn\%27t_read}{tl;dr} or summary. In general emails should be short. Longer discussions should be held on the phone or in person, with a summary report after the discussion. Reliance on a BLUF or tl;dr risks resulting in the reader skipping the content. 

Emails convey both emotional tone and facts. Your intent is practically irrelevant; the reader's perception is paramount. 

Every email should have a purpose. What are you asking the recipient to do? How do you want them to feel? How should they respond?

When replying, starting your email with an expression of gratitude for the work the recipient has done so far sets a positive tone by acknowledging their investment.

\ \\

\subsection*{Summary of what action should be carried out} 

As the outsider, you should help the community enumerate and document all of the problems they identify. Then you can help enumerate and document how the problems are related (dependencies). Only then can you help the community identify and document the root causes.

If the solution you, the outsider, identified really is the root cause, then the community will arrive at that independently. If that is the case, then you can enable them to implement a solution which addresses the root causes. The community will then have a sense of ownership.
 % subsection 
    \subsection*{Friction between teams within an organization}

Ideally there is a clear division of responsibilities among teams. Even in that context there is necessarily some interaction among teams -- one team may depend on the output from another team. Coordination among the teams regarding transfer of data or products or projects or knowledge is critical to the smooth operation of the organization. 

An organization has finite staffing, money, time. Therefore, teams within the organization face a \href{https://en.wikipedia.org/wiki/Zero-sum_game}{zero-sum}
\index{Wikipedia!\href{https://en.wikipedia.org/wiki/Zero-sum_game}{zero-sum game}}
distribution of resources.

Difference of decision making perspective based on local conditions, stupidity, or different incentives, different definitions of success

When attempting to resolve friction between teams, there is an authority common to the teams, but that person lacks the nuanced insight, doesn't have time to get involved in every challenge, and doesn't want to micromanage multiple teams. % subsection

    %There are many books on how to do meetings well. 
    % How is this section (in a book about bureaucracy) distinct?
    % how does understanding bureaucracy matter in the context of a meeting? what's the change in behavior for the effective bureaucrat?
    Stakeholders care about the outcome of a decision. 
    In some sense, every participant in an organization is invested the outcome of every decision made because there is consequence to how resources are allocated and the direction of the organization. Every org member has an opinion, even lacking experience or expertise. 
    
    \subsection{Financial models of communication}

Late to a meeting is theft, and poorly written email and email to the wrong people has financial cost.

\subsubsection{Meeting time compared to theft}
% https://graphthinking.blogspot.com/2021/02/organizations-value-things-more-than.html

In large organizations, there can be significant bureaucracy associated with even small purchases. A multi-step review process may be incurred for a \$2000 acquisition.

Another measurement of value is that if an employee were to steal even \$200 worth of materials, the organization would likely punish that employee.


In the book High Output Management \cite{1995_Grove}, Grove points out that those metrics apply to tangible goods, but not to people's time. Consider a meeting of 10 people and each person's cost is \$200 per hour. A wasted meeting is not unusual and certainly would not incur bureaucratic review processes. The cost to the organization is fiscally the same -- \$2000. Similarly, consider an employee who is late and causes a loss of productivity. Merely depriving the organization of \$200 worth of time is not punished in the same way theft is.

In fact, organizations default to meetings (even recurring meetings) rather than not meet. And being late to a meeting is accepted. 

We can debate the differences between theft of materials and theft of time. The financial argument is clear. 


\subsubsection{Email is not free}

Total cost of email:
* time spent writing
* time spent reading
* infrastructure maintenance

Cost per email = 
(hourly rate of writer)*(time spent writing) + (hourly rate of reader)*(time spent reading)*(number of readers)+(annual salary of maintainer)/(number of emails per year) + (email server cost)/(number of emails per year), etc

50*(5/60) + 50*(5/60)*4 + 100000/1000000 + 10000/1000000 = \$21 for one email. 

\begin{figure}
    \centering
    \includegraphics{images/email_costs_5minutes_4people.pdf}
    \caption{If email is 2 hours a day, and staffing is 50\% of the organizations budget, then 
    % (2/8)*0.5
    12.5\% of the organization's budget is spent on email. Same math for meetings, so that's 25\% of the organization's budget on coordination.}
    \label{fig:my_label}
\end{figure}

% https://docs.google.com/spreadsheets/d/1ysV5PA3cEcneKv5BViUhAfgNq26cOlNRj6l5fXuHM3Q/edit?usp=sharing

Paradigm for reading emails
* Unread
* Does this message have any bearing on my current task? 
* Integrate with my understanding 
    
    \subsection{Well-run meeting\label{well-run_meeting}}

Identify essential attendees. If someone does not need to be present, notify them in advance that you will share the meeting notes afterwards. 


Bad: no meeting agenda\\
Good: agenda\\
Better: agenda share with other participants
For formal meetings, share agenda in writing prior to meeting. 

TODO: Why an agenda matters in a bureaucracy: 

TODO: forces conspiring against agendas


For formal in-person meetings, Verify meeting venue has sufficient space, seating, working IT equipment

For formal virtual meetings, ensure participants are familiar with virtual meeting controls

TODO: why logistics/infrastructure matter in a bureaucracy:

TODO: forces conspiring against logistics/infrastructure



\subsubsection{Meeting facilitator skills}

To run a smooth and productive meeting, I explicitly state two ground rules to the attendees:
\begin{itemize}
    \item If you want to talk, raise your hand and I will call on you. If there are multiple people wanting to talk, I'll track the order of speakers.
    \item If you talk too long, I'll cut you off. 
\end{itemize}

As a facilitator, my focus is on structure (distinct phases of the discussion) and ensuring participation. I remove myself from taking part in the discussion. % subsection
    \section{Characterizing Meetings}
% relevance of this section:
Characterizing meetings is critical to distinguishing which norms are applicable, and what people expect from the different formats. 



Types of meetings: internal meetings, customer meetings, conferences, scheduled one-on-ones, impromptu walk-around.   


%What is the purpose of a meeting?
% https://graphthinking.blogspot.com/2019/12/what-is-purpose-of-this-meeting.html
There are many potential purposes of a meeting, including
\begin{itemize}
    \item To gather input from attendees.
    \item To make a pronouncement to attendees.
    \item To educate attendees.
    \item To educate one person.
    \item To signal interaction.
    \item To brainstorm ideas.
    \item To make progress towards an goal.
\end{itemize}
When the purpose is not explicitly stated, confusion arises. 
When multiple purposes occur in one meeting and the transition is not explicitly stated, confusion arises.
The reason for this confusion is that the assumptions and expectations and norms of each purpose are different. When the attendees don't know the purpose or the purpose shifts, the expected behaviors and roles are unclear. 

An attendee can ask what the purpose of the meeting is during the meeting but that is generally considered rude. An attendee can try to deduce the purpose of a meeting, but this takes time and attention and can result in the wrong conclusion. An attendee can try to set the purpose of the meeting during the meeting, but this can conflict with the intent of other attendees. 

A meeting's purpose can shift during a meeting. If done intentionally, the changes should be stated explicitly. Otherwise an attendee may continue to work under the previous set of expectations rather than the current norms. 



% https://graphthinking.blogspot.com/2014/12/how-to-understand-meetings-at-work.html
Level of formality, start time (early or on time or late), 
end time (early or on time or late), utility, 
duration, number of attendees, number of speakers, number of participants.



Meetings involve people, either known or strangers.
Meetings involve information, either relevant or irrelevant. Relevant information is either new or related to previous work.
Meetings either have a leader or no leader (brainstorming). If there's a leader, the leader may be disseminating info to participants, or gathering information from attendees.


\subsection*{Phases of a meeting}

\begin{enumerate}
    \item Establish understanding of what the topic is. Without a shared focus and a common goal a meeting is unlikely to be productive. 
    \item Set aside topics that are not the focus. Either discuss outside the current meeting or defer to a later meeting. Managing scope is critical to accomplishing the goal. 
    \item Establish a shared language specific to the topic and to the participants.
    \item Establish each participant's viewpoint on the topic.
    \item Brainstorm options.
    \item Build consensus or nominate a decider.
\end{enumerate} % subsection
%    \subsection{Make Effective Presentations\label{sec:presentations}}

% https://graphthinking.blogspot.com/2011/10/presentation-notes.html

By paying attention to a bad presentation, you can find problems that you do not want to repeat.

Speaking is vital at decision points in your career progress - at interviews, competitions, gaining new collaborators at conf, impressing peers/boss/students. Thus, it's best not to make these mistakes in those situations.

Aspects to a presentation
\begin{itemize}
    \item speaker: appearance, verbal accent, enunciation, smell
    \item slides: layout, content
    \item audience: background
    \item venue: size, technology available, lighting
    \item lead time: how much time do you have to prepare?
\end{itemize}
The following list is meant to be reviewed as a checklist.

Planning your presentation

In introducing your research, there are a few ways to open the talk
\begin{itemize}
    \item compliment the audience
    \item humor: relevant joke
    \item how does this work make me feel
    \item explain context and relevance in terms of money, number of people involved, size of system
\end{itemize}
Estimate the experience/education background of your audience prior to creating the presentation. The question being answered in the presentation is independent of audience, but the level of delivery depends on audience background.


Speaking to an audience outside your field:
\begin{itemize}
    \item Use jargon your audience is familiar with. For example, a physicist uses "quantization of energy" whereas an audience of mathematicians may be more comfortable with "discrete energy levels."
    \item Look for commonality. For example, both physicists and mathematicians use assumptions and build models.
\end{itemize}
Visual presentation:
\begin{itemize}
    \item Switching between dark and light slides in a dark room stresses the eyes (need time to adjust to varying light levels)
    \item In figures, use both color contrast and distinct symbols, in order to deal with colorblind audiences. See also http://jfly.iam.u-tokyo.ac.jp/color/
    \item making slides appear "professional" means adding non-informational content. This added content should be consistent, not distracting
\end{itemize}
Software:
Common choices include Latex (beamer), Microsoft Powerpoint. Less well-known are Prezi (and its open source equivalent, InkScape+Sozi add-on).

If you use Microsoft Word, a document PDF, notepad, or any other non-presentation software to make a presentation, then you are sending a few messages to your audience:
The audience isn't worth your time needed to develop or learn a proper presentation
You are not technically savvy

Presenting multiple topics within one presentation:
Disparate topics require a segue to show why you are transitioning
Inter-relate the multiple topics 
Dealing with technical failure
If you plan to give a live demo, have screenshots of the process in the presentation. That way, if the demo fails you can show what was supposed to happen in the slides. If the demo works, skip the slides.
If possible, use the setup as close to reality as possible for practice sessions. Project onto the screen using the projector. This will show if the color contrast is sufficient.

Prior to giving a talk at a remote (or unfamiliar) environment

Questions to ask your host:

Will I have a projector and screen for the presentation?

If Power point available, what version is in use?

Will I be allowed to bring a USB drive and/or laptop into the location (security restrictions), or should I email the pptx file to you?

Day of talk

Appearance (suit and tie or jeans and t-shirt?):
\begin{itemize}
    \item under-dressed = I don't respect the audience
    \item overdressed = I'm better than you
    \item similar level of dress = I'm a peer
\end{itemize}
Some of these suggestions may seem glaringly obvious. The reason they are here is because I have personally seen them in "official" presentations given by a "professional"
\begin{itemize}
    \item Do not curse (profanity while speaking, or in the slide presentation, or even the name of the file)
    \item Practice the presentation (out loud in real time) at least once
    \item Run spell check before presenting
    \item Translate your slides from your native language to the language of your audience
    \item Tell a coherent story with a unified theme. Each slide should be logically connected to the following slide. Don't just put a bunch of slides with data together. You risk disorienting your audience.
\end{itemize}
Before the talk begins
\begin{itemize}
    \item Make sure your facility has power, a screen, a projector, pointer, any other necessary equipment. [Your host may not think of these things for you.]
    \item Make sure all equipment works and functions together
    \item Before the presentation begins, use a slide to make announcements and reminders:
\begin{itemize}
    \item list agenda if there are multiple speakers
    \item time talk begins, how long it will last
    \item Reminder: Turn OFF Cell Phones
\end{itemize}
    \item Announce whether to ask questions during talk (interrupt) or to hold questions until done
\end{itemize}
During the talk
\begin{itemize}
    \item Opening: Thank hosts/inviters/organizers. Establish a connection between you (the speaker) and the audience.
    \item Do not pace
    \item Do not stand frozen
    \item If you are the only person laughing at your joke, it isn't funny
\end{itemize}
 
After the talk
\begin{itemize}
    \item Let the person asking the question finish
    \item Restate the question to ensure the audience heard it and that you understood it
    \item Give the minimal answer. Save the long answer for followup
\end{itemize}

    % how to engage your audience
    % stages of dysfunction
%    % https://graphthinking.blogspot.com/2021/02/organizations-value-things-more-than.html
\section{meeting time compared to theft}

In large organizations, there can be significant bureaucracy associated with even small purchases. A multi-step review process may be incurred for a \$2000 acquisition.

Another measurement of value is that if an employee were to steal even \$200 worth of materials, the organization would likely punish that employee.


Those metrics apply to tangible goods, but not to people's time. Consider a meeting of 10 people and each person's cost is \$200 per hour. A wasted meeting is not unusual and certainly would not incur bureaucratic review processes. The cost to the organization is fiscally the same -- \$2000. Similarly, consider an employee who is late and causes a loss of productivity. Merely depriving the organization of \$200 worth of time is not punished in the same way theft is.

In fact, organizations default to meetings (even recurring meetings) rather than not meet. And being late to a meeting is accepted. 

We can debate the differences between theft of materials and theft of time. The financial argument is clear. 

Source: Andy Grove in "High Output Management" % no value to the reader
    
    \subsection{Walk-arounds}

Value of daily walk-around: 
\begin{itemize}
    \item social engagement; see \S\ref{sec:socializing}
\item \href{https://en.wikipedia.org/wiki/Allen_curve}{Allen curve}; 
\item presence creates priority
\end{itemize}

% https://graphthinking.blogspot.com/2016/04/role-models-for-great-leaders.html
\href{https://en.wikipedia.org/wiki/Friedrich_Wilhelm_von_Steuben}{Baron von Steuben} and George Washington both did walk arounds to understand needs of their army.  % subsection
    \section{One-on-one check-in meetings}

% TODO: I'm avoiding tips on "how to manage." 
% Should I also avoid tips on "how to coach"?
% I have a lot of potential content on 
% https://graphthinking.blogspot.com/2021/09/notes-from-half-day-course-on-coaching.html

Formal one-on-one meetings of a supervisor and a team member. Opportunity for the supervisor to coach the team member and for the supervisor to learn what the problems are on the team. From the team member's perspective, a one-on-one is an opportunity to ensure alignment with the team's direction and to provide insight on how to improve the team. Education in both directions. 

% https://graphthinking.blogspot.com/2021/09/notes-from-half-day-course-on-coaching.html
One-on-one discussion take lot of time, so the mindset ``no news is good news" or ``everything is going well, no need to engage" is tempting. As a consequence only negative feedback causes interaction. 

% https://graphthinking.blogspot.com/2021/05/the-agenda-for-one-on-one-meeting.html

One-on-one meeting questions for helping the supervisor understand the team member's status. 
The one-on-one check-in should be tailored to the phase of the employee's progression. Frequency of check-ins depends on the newness of the team member, the complexity of the work, or how quickly the conditions are changing.
\begin{enumerate}
    \item \textit{Name of phase}: \underline{Welcome to the team!}\\
    \textit{Scenario: New team member, either new to the team or new to the organization. }\\
    Here the focus of the one-on-one is to ensure a smooth on-boarding process. Get them up-to-speed on the technical challenges, professional norms, and integrated with other team members. Resolve administrative blockers: Does the employee have the necessary computer log-in accounts? Do they have an email account? Are they on the mailing list? \\
    Questions:
    \begin{itemize}
        \item What are the objectives for the team?
        \item What items on the onboarding checklist are not yet completed?
        \item Who have you met on the team? What is your understand of their role on the team?
    \end{itemize}
\textit{The duration of this phase could last between a day and two weeks.}
    \item \textit{Name of phase}: \underline{Initial contributions}\\
    \textit{Scenario: Team member is responsible for small tasks. }\\
    The one-one-one is for discussions on training and planning and task reviews. Characterized by the team member being dependent on others for their success. In this phase the employee collaborates on tasks.\\
    Questions:
    \begin{itemize}
        \item What are the objectives for the team?
        \item What are your task objectives?
        \item What are you expecting to deliver to the team? When? 
        \item What dependencies does that deliverable have (external to the team or internal to the team)?
    \end{itemize}
\textit{The duration of this phase could last a few months to years.}
    \item \textit{Name of phase}: \underline{Experienced contributor}\\
    \textit{Scenario: Team member is responsible for large tasks (which get broken into subtasks). }\\
    The one-on-one is to help the team member define their success. Activities include planning, resource allocation, assessment. Characterized by the need to coordinate with others on the team or other teams. Team member understands task scope and intent and relevant processes. Team member decomposes task into subtasks.\\
    Questions:
    \begin{itemize}
        \item How do the artifacts you're working on support your plan for the team's progress?
        \item What dependencies does that deliverable have (external to the team or internal to the team)?
        \item What insights do you have about the team or organization?
        \item What insights do you have about the relevance of the task relative to the intent of the organization?
        \item What should management be doing to enable the team's success?
    \end{itemize}
\textit{The duration of this phase could be the rest of a career.}
    \item \textit{Name of phase}: \underline{Facilitator}\\
    \textit{Scenario: Facilitating the productivity of others.}\\
    Rather than being task-oriented, this team member supports coworkers. \\
    Questions:
    \begin{itemize}
        \item What observations from mentoring team members do you have?
        \item What collaborations should we be fostering?
    \end{itemize}
    \item \textit{Name of phase}: \underline{Peer}\\
    \textit{Scenario: Peer check-in.}\\ 
    This one-on-one is a form of mentorship. The value of the exchange is to get a different perspective and to hold each other accountable.
\end{enumerate}

%TODO: How does the team member and the supervisor know when the next phase is appropriate?

%TODO: What are the thresholds for change?

In addition to phase-specific questions, there are questions that fit all phases. A supervisor should ask the team member
\begin{itemize}
    \item What have you been successful with since we last met?
    \item What is blocking our team's progress?
    \item What are your plans?
    \item How are you collaborating with the rest of the team?
    \item If there was just one thing you could change about our organization, what would it be and why?
    \item How do you plan to train your coworkers on topics you understand and they don't?
    \item What have you learned in the past month?
    \item What are the biggest risks for the team?
    \item What's limiting your productivity?
\end{itemize}
Responding to these questions takes time (an hour) and willingness to be open. 

If the supervisor doesn't ask these questions, the team member has the ability to include these in the agenda and to bring them up. 


\ \\

A one-on-one meeting requires preparation. Before the meeting the team member should document
\begin{itemize}
    \item What was discussed previously
    \item What progress has been made since the previous meeting
    \item What is blocking progress
\end{itemize}


\ \\

% lots of comments, not much useful info
% https://news.ycombinator.com/item?id=30152268

\ \\

\subsection*{One-on-one meeting questions to spur discussion}

Not sure what to discuss? An extensive list of questions on topics like career development, conversation starters, and job satisfaction are available on \href{https://github.com/VGraupera/1on1-questions}{https://github.com/VGraupera/1on1-questions}\footnote{\href{https://news.ycombinator.com/item?id=22341138}{comments}}. % subsection
    \subsection*{Listening to Bad Presentations\label{sec:bad-presentations}}

Occasionally you attend a meeting with a bad presentation. The slides may appear slick but the content is poorly thought out. Or the presenter does not understand the topic well. Or the presenter has good content but does not convey it well. Or the presenter is wrong about the topic. Regardless of the cause, you should assume the presenter is making their best effort. 

You could remain silent, complain, criticize, ask leading questions, or offer constructive feedback. Your silence may result in other attendees and the presenter leaving with incomplete or wrong information. If you speak up you'll prolong the meeting or limit the presenter's time to convey their material. 

Your assessment of the presentation may be wrong. You may lack relevant information. A reliable technique for interjection is to assume a state of confusion instead of confidently asserting that the presenter is wrong. 
If you believe the presenter is wrong, asking about the source of their information is a good entry point.

You should ask for clarification when the information is correct but presented poorly (or above the level you understand). 

\ \\

% TRANSITION

By paying attention to a bad presentation, you can  identify issues that you do not want to repeat. The next section summarizes a few insights so you don't need to learn from bad presentations. % subsection
%    \subsection{How to be a Successful Conference Attendee}
% https://graphthinking.blogspot.com/2012/02/how-to-succeed-as-attendee-at.html
%    \subsection{How to be a Successful Conference Organizer}

% https://graphthinking.blogspot.com/2011/11/how-to-organize-conference.html
  % multi-person interactions = process
  \chapter{Bureaucratic Processes\label{sec:process}}
    % Claim: you already have habits; these are similar to processes use by bureaucrats in organizations

Because bureaucratic policies and processes seem convoluted, let's start from a more relatable point: your personal habits.
You might eat at the same time every day, or you might go to bed at the same time. Those are both examples of personal policies. Sticking to a routine decreases the need to think about options. In the same way, bureaucrats in organizations seek routines to decrease uncertainty. 

You have experience with creating and using personal policies (your habits).  Bureaucratic processes for organizations are similar.
Both are routines which ease the burden of decision-making. A habit can be intentional (e.g., my policy is to brush my teeth before going to bed), just as a process can be designed. Habits can be unconscious (e.g., arriving at work without recalling driving there), just as processes can arise without an intentional design. 

Consider your intentional habits -- they are likely motivated by a personal policy. Similarly, bureaucrats acting on behalf of an organization apply processes that support policies of the organization.

\ \\

The sections of this chapter address the what (page~\pageref{sec:definition-of-process}), why (page~\pageref{sec:why-processes-exist}), and how (page~\pageref{sec:process-chaos}) of processes. 
The last sections of this chapter (design on page~\pageref{sec:design-of-processes} and change on page~\pageref{sec:change-a-process}) assume a level of control and autonomy that you may not think of you have. Regardless of where your role is in the hierarchy, your purpose is to provide value to the organization. That value can be negotiation with fellow bureaucrats on improvements to the design of task workflows. 

The foundation of your \gls{process empathy} is based on the sections in this chapter, whereas previous chapters have been focused on your skills as a bureaucrat. 
    \subsection{Definition of Process}

% definition
Policy says what action is allowed or required
Process is a sequence of tasks, with each task invoking policy enforcement 


Forms are meant to ensure consistent application of policy, and to catch people who shouldn't get the resource (both in cases of accidental request and malicious request). For malicious requests, a more burdensome process merely filters the low-cost grifters. 



A \gls{process} is a task broken into a specified set of subtask dependencies, typically with subtasks in a sequence. 
Each task is associated with the application of policies enforced by bureaucrats. 
% Also known as a procedure. 
Two distinguishing features in the context of bureaucracy are authorization and justification.  
A process has inputs and outputs. 
A process can be decomposed into other processes. 
Processes operate on both information and tangible objects. 
Processes require \href{https://en.wikipedia.org/wiki/Work_(physics)}{work} and time. 
Processes are carried out by people or machines.

% process and roles
processes depend on roles that have responsibility/authority
When there aren't sufficient staffing, then individuals wear multiple hats
Commonly leads to conflict of interest among roles, which is defeated by one person in multiple roles
eg "review the implementation" and "implement"

\href{https://en.wikipedia.org/wiki/Responsibility_assignment_matrix}{Responsibility Assignment Matrix}

% how does understanding process help me?
The point to understanding process and policy is so that you know when to work within versus work around, when to accept versus change, when to ignore, how to leverage, and how to design.



% dangers of process
The simplification and the neglect of specific circumstances results in process friction. Process friction manifest in practical waste, temporal and efficiency, emotional frustration, and social distrust of institutions.




\begin{figure}
    \centering
    \includegraphics{images/process_loop_harm-oversight-improvement}
    \caption{loop: harm-oversight-improvement}
    \label{fig:harm-oversight-improvement}
\end{figure}



Bureaucracy is defined by processes.
As a counter force, social relationships.
Actually, processes are social relationships formalized

Processes are not consistent across people because the relationships vary


processes are about a defensible story for the bureaucrats involved -- not the actual equitable distribution



% https://graphthinking.blogspot.com/2021/04/laffer-curve-and-minimum-viable.html
social processes can work well for small tasks that are infrequent
social processes can be effective in anomalous situations
bureaucratic processes are useful under typical conditions

When no process exists, only people with relationships succeed. Processes enable novices to an organization to contribute value and gain experience. 





Does a process exist?

% https://graphthinking.blogspot.com/2015/07/notes-on-bureaucracy-and-social-network.html
Do participants know about existing processes?

Processes can be undocumented. Then oral folklore is the mechanism. 

% process design trade-offs
Processes with fewer people and fewer steps can be quicker and use fewer resources, but they are more fragile and less representative. Having more people involved helps with edge cases, but slows down the process. 



    \section{Why do Processes Exist?\label{sec:why-processes-exist}}

Bureaucratic organizations manage shared resources (tangible or expertise). 
A subjective decision about a resource that is  administered by bureaucrats is a \gls{policy}.
Processes inform the decision making of bureaucrats and result in access to or denial of shared resources. 

The following sections provide more explanation for why organizations rely on processes. The reasoning matters to you as a bureaucrat when you design a process or try to revise an existing process. Similarly, as a person going through a process, the reasoning below explains why you are experiencing \gls{process friction}. 

\subsection*{Processes Ensure Consistent Application of Policy}

Organizations are designed to be insensitive to individual participants; that applies to both bureaucrats and subjects. Processes address risk to the allocated resource and risk to the bureaucrats carrying out the process.

The risk addressed by the use of processes is both current bad behavior and the potential of things going bad in the future. Process is the guardrail limiting harm to bureaucrats and the shared resource. 

One motive for bureaucracy is the \href{https://en.wikipedia.org/wiki/Tragedy_of_the_commons}{tragedy of the Commons}, which says that when there is a shared resource someone will try to get away with behavior that is harmful to the community of users. Limiting harmful behavior can take the form of oversight processes. If oversight processes are not in place malicious actors dominate.

Another example guardrail is keeping stakeholders informed (justifications) so that intervention can be taken if needed. 


In isolation each process for oversight/review/approval may be justifiable, but the aggregate can feel unreasonably burdensome to both subjects and bureaucrats.


% https://advancedbiofuelsusa.info/rodney-hailey-sentenced-to-more-than-12-years-in-prison-for-selling-9-million-in-fraudulent-renewable-fuel-credits/



\subsection*{Processes exist to Simplify}
Compared to an ad hoc approach, processes are intended to simplify what could otherwise be complicated moral decisions, coordination challenges, or financial assessments. %The simplification benefits the bureaucrat's workload; rarely are benefits to the subjects of the bureaucracy.  



A positive view is that specialization allows narrow focus and thus deeper understanding and skill. An alternative perspective is that specialization allows for dumbing down the role, thus enabling a cheaper workforce. In that view process allows dumb individuals to, as a group, do complicated things.  As a specific example, consider the process of designing and building a car. That complexity is feasible to undertake for a skilled and knowledgeable individual, but the cheaper approach is to hire individuals capable of installing the passenger-side doors in an assembly line.

The relation between complexity of a task and skill of the bureaucratic workforce is another motive for the creation of processes. See figure~\ref{fig:complexity-and-size} on page~\pageref{fig:complexity-and-size} for a visualization of the trade-off.

\subsection*{Processes Enable Increased Throughput}

Processes that leverage specialization can improve scalability. Adam Smith's \href{https://en.wikipedia.org/wiki/Business_process#Adam_Smith}{pin factory example} cites a productivity gain of 240x. Similarly, the introduction of moving assembly lines in Ford's car factory produced an 8x improvement of throughput \footnote{Wikipedia entry describing \hyperlink{https://en.wikipedia.org/wiki/Assembly_line\%2320th_century}{Ford's assembly line}}. 

While social processes may work well for low-volume ad-hoc requests, formalized bureaucratic processes become necessary for recurring high-volume activities.

\subsection*{Processes Decrease Reliance on Social Relations}

The relation between bureaucratic processes and social bonds was noted by Selznick in his 1943 paper~\cite{1943_Selznick}. The needs of individual members of an organization may be at odds with (or at least not aligned with) the purpose of the organization. This dissonance is observed as the interplay between personal relationships and bureaucratic processes.

Generalizing claims about social bonds in an organization is unfounded. 
For a given complexity and given scale, there are people who are more social and less social and therefore desire more or less process.
There are bureaucrats who want documented processes and are confused as to how things are operating when those aren't present. 

The invisibility of social relationships limits bureaucrats not attuned to the importance of these bonds. 
Social influence among bureaucrats lacks transparency-- discoverability and documentation. 
In contrast, processes are easier to understand and to track. The counter to social influence is the hierarchical \gls{org chart} described on 
page~\pageref{sec:org-chart-as-guide-and-lie}.
Social influence is not antithetical to processes. There's always a mixture of the two.


\subsection*{Processes Help New Bureaucrats and Subjects}

Bureaucrats new to a team are likely to ask, ``What are the processes and where is the documentation?'' Bureaucrats who have been on a team for a long time say, ``Don't burden me with processes, I just need to get things done using the relationships with the people I have.'' One motive for processes is to help with onboarding team members who don't have relationships available.

Besides not having formed a social network, there's a second reason new people seek processes. New team members who recently graduated from school are used to the existence of formal processes at the high school or university level. Therefore when they join a job, they expect a similar set of conditions for processes to exist and follow.

Once this motive for processes is recognized, the relevance when onboarding new hires is apparent. New hires will need to discover the existing processes while they form social bonds. Discovering processes comes through oral folklore or written documentation. To augment both process discovery and creation of relationship, one technique is to have new hires sit with experienced team members.
\marginpar{[Tag] Actionable Advice}
This is described in the section on  
\hyperref[sec:prisoner-exchange]{Prisoner exchange} on 
page~\pageref{sec:prisoner-exchange}.
Where you sit matters to your effectiveness as a bureaucrat.

\subsection*{Processes exist as Artifacts for Promotion}

Processes can arise organically (bottom up) or be created top-down. In either context, creating new processes counts as bureaucratic innovation. Process creation is confused as promotion-worth because of the distinction between novelty and innovation. Real innovation incurs risk, whereas new processes are merely novel. Both cause change, but processes are typically designed to decrease risk. 


    \section{Riding the Chaos\label{sec:process-chaos}}

Bureaucratic organizations are imperfect because they are composed of humans. %Somethings work, others do not. 
A bureaucratic organization is a complex and chaotic environment.\footnote{Chaotic in the \href{https://en.wikipedia.org/wiki/Chaos_theory}{mathematical sense} 
\index{Wikipedia!\href{https://en.wikipedia.org/wiki/Chaos_theory}{chaos theory}}
of non-linear dependence on initial conditions, not just colloquial sense of disorder} An understandable response is to impose processes and hierarchy and roles. Imposing structure can feel intellectually fulfilling and emotionally satisfying. Imposing structure looks like progress and may even help with your promotion in the hierarchy. The challenge (or continual source of employment) is that chaos is dynamic. A process created in response to chaos create new challenges and is disrupted by changes in number of staff, changes to who is part of the process, changes to amount of work. The half-life of structure depends on the rate of change of tasking and the rate of 
\marginpar{See page~\pageref{sec:turnover}.}
\hyperref[sec:turnover]{personnel turn-over}\iftoggle{haspagenumbers}{; see page~\pageref{sec:turnover}.}{.}

A bureaucrat in the organization could strive for perfection, run away from the chaos (to something less chaotic), exploit the chaos, or adopt an attitude of ``I do what I can'' or ``I'll wait out the problem.'' 
Another option, and the focus of this section, is to build skills for navigating the chaos of a bureaucratic organization.

You can build and maintain your network of professionally connections with fellow bureaucrats. 
\marginpar{$>>$ Actionable Advice} 
\index{actionable advice}
This augments the chain of command of a hierarchy. There is a turn-over in both hierarchical roles and in your professional network, so you need to continually invest effort in creating new bonds and maintain existing relations. 

You can see change as an opportunity for improvement rather than a destruction of your previous investments. The \href{https://en.wikipedia.org/wiki/Sunk_cost}{sunk cost fallacy}
\index{Wikipedia!\href{https://en.wikipedia.org/wiki/Sunk_cost}{sunk cost fallacy}}
applies both to your emotional state and the resources you invested. Fear of change is understandable -- it disrupts the comfort of stability and known processes. Anticipating change and preparing for it eases the stress of change.

You can leverage personalities and unique traits rather than expecting everyone to be interchangeable and treating differences as flaws. 
\marginpar{$>>$ Actionable Advice}
\index{actionable advice}
By knowing the strengths, interests, and weaknesses of coworkers you can facilitate the process of change. 
    \section[Processes Involve a New Person]{Processes Involve a Person New to the Process\label{sec:new-person}\iftoggle{shortsectiontitle}{\sectionmark{Processes Involve a New Person}}{}}
\iftoggle{shortsectiontitle}{\sectionmark{Processes Involve a New Person}}{}

% LONG1 shows up in the TOC
% LONG2 is the section title

% https://graphthinking.blogspot.com/2022/07/bureaucratic-processes-typically.html

% The IPYNB notebook is in
% https://drive.google.com/drive/u/0/folders/1awZYk4EisFIR3mY1iF2X7Hb6CwcNDavp

Bureaucratic processes require working with other people. One source of friction is that participants may lack familiarity with the process. 

This observation can be quantified with a small number of assumptions. The distribution of team membership tenure might fit a \href{https://en.wikipedia.org/wiki/Power_law}{power law distribution}%
\index{Wikipedia!power law@\href{https://en.wikipedia.org/wiki/Power_law}{power law}}
%\footnote{See \href{https://numpy.org/doc/stable/reference/random/generated/numpy.random.power.html}{numpy.random.power} documented on https://numpy.org/doc.}
-- there are more inexperienced people than experienced people. What matters in this context is how long the team member has been in their role (rather than how long they've been a member of the organization). Let's assume a max tenure of ten years~\cite{2022_BLS_tenure}. If the process involves five people, the least experienced member will have a median of 100 calendar days of experience.

% image float options:
% https://tex.stackexchange.com/a/32605/235813
\begin{figure}[!htb] %[H]
    \centering % 0.8 is too small
    \includegraphics[width=1\textwidth]{images/tenure_power_distribution_a5_with_max_tenure10_and_5_participants.pdf}
    \caption{The power law distribution of tenure using the \href{https://en.wikipedia.org/wiki/Probability_density_function}{probability density function}
    \index{Wikipedia!probability density function@\string\href{https://en.wikipedia.org/wiki/Probability_density_function}{probability density function}}
    described by $a\ x^{(a-1)}$ where $a=5$. The median tenure of the youngest participant  in a process with five people is 100 days. Population size has been normalized to 100 people.}
    \label{fig:tenure-powerlaw-5-participants-tenure10}
\end{figure}


Three calendar months (or 70 business days) may be inadequate for complex processes or processes that are infrequent (quarterly or annual), especially if there was no training. When the number of participants is ten people, then the median tenure of the youngest participant is 51 calendar days (37 business days).

\ \\

\begin{samepage}
To recap, the assumptions made were:
\begin{itemize}
    \item Random, independent sampling of organization members. 
    \item Tenure in your role matters, not the tenure in the organization. (This assumption excludes transfer learning among roles.)
    \item Tenure in role fits a power law distribution.
    \item The power law distribution is characterized by (a=5, max tenure=10 years). 
    \item There are five people involved in the process.
\end{itemize}
\end{samepage}
If that last parameter is varied, the estimate of three months as the median  is reasonable for processes with five or more participants.

% image float options:
% https://tex.stackexchange.com/a/32605/235813
\begin{figure}[!htb]  %[H]
    \centering % 0.8 is too small
    \includegraphics[width=1\textwidth]{images/tenure_power_distribution_a5_with_max_tenure10.pdf}
    \caption{Median tenure of the youngest participant (in calendar days) as a function of the number of process participants with tenure following a 
    \href{https://en.wikipedia.org/wiki/Probability_density_function}{probability density function},
    \index{Wikipedia!probability density function@\string\href{https://en.wikipedia.org/wiki/Probability_density_function}{probability density function}}
    $a=5$, and max tenure of ten years.}
    \label{fig:tenure-powerlaw-5-participants}
\end{figure}


Even when there's a single participant, the median tenure is less than two years because of the power law distribution of tenure.

% Uniform Distribution of Tenure instead of a Power Law Distribution
What if the power law distribution of tenure were replaced with a uniform distribution?
Surprisingly the shape of the distribution of the youngest participant's tenure is not uniform, see Figure~\ref{fig:tenure-uniform-5-participants}

% image float options:
% https://tex.stackexchange.com/a/32605/235813
\begin{figure}[!htb]  %[H]
    \centering
    \includegraphics[width=1\textwidth]{images/tenure_uniform_distribution_with_max_tenure10_and_5_participants_median472.pdf}
    \caption{The median tenure of the youngest participant
 in a process with five people is 472 days when the distribution of tenure is uniform. The population size has been normalized to 100 people.}
    \label{fig:tenure-uniform-5-participants}
\end{figure}

The median tenure of the youngest member of a process with five participants is 15 months. 

% https://latex.org/forum/viewtopic.php?t=17111
\FloatBarrier
    \section{Exceptions to a Process\label{sec:exceptions-to-process}}


Process designers are usually motivated to address a specific issue associated with access to shared resources. When bureaucrats execute a process to support a policy, there may be cases that violate the designer's assumptions. Now the person with the exceptional case and the bureaucrat carrying out the policy has extra work to deal with the broken process. This is a source of \hyperref[sec:bureaucratic-debt]{bureaucratic debt}. 
\marginpar{See page~\pageref{sec:bureaucratic-debt}.}
%\ifsectionref
%(see section~\ref{sec:bureaucratic-debt} for details). 
%\fi

When the routine process is inadequate, there are routes for exceptions:
\begin{itemize}
    \item Change the process to address the exceptional case. This takes extra time and work for the subject and the bureaucrat.
    \item Ignore the process. Use relationships instead.
    \item Ignore the task. If the result isn't crucial to your efforts and the work is too burdensome compared to alternatives, stop the process.
    \item Seek an exception to the process. The exception can be recurring or one-time. The justification for the exception can be based on social capital of participants (title, reputation) or be a data-driven argument or some mixture of both.
\end{itemize}

The designer of a process may not be able to anticipate future circumstances, but expecting the capacity to be responsive to change is reasonable. The process designer should use sunset provisions and escape hatches.  
\marginpar{$>>$ Actionable Advice}
\index{actionable advice}
A sunset provision can be time-based (this policy expires after three years) or based on a threshold (this policy should be evaluated for renewal after 1000 cases). An escape hatch specifies the conditions under which the participants are excepted from the policy. A good escape hatch for a policy provides directions about what follow-on actions should be taken.

Another reason processes may be inadequate is due to urgency. A result that is needed quickly may not fit the timeline of sequential tasks for a process.

\subsection*{Time Sensitivity Exceptions}
Exceptions to a process may be required when a result is needed quickly. In the medical setting, there are categories of urgency; a task may be routine, priority, urgent, stat. 
Each of those categories is associated with a time-bound, though the specific duration varies by medical domain and by location~\cite{2022_priority_definitions, 2018_order_priority}.


Common failure modes for processes include abuse of prioritization (everything is important) or abuse of urgency (everything is needed quickly). Either of these abuses need to be corrected to have a healthy bureaucracy, but typically the cost of labeling an effort as ``priority'' or ``urgent'' is low. 

Interaction between the bureaucrats tasked with work and the people characterizing relative value is needed to limit abuse of processes. This interaction is a negotiation of what work gets done when and in what order. Labeling every task as ``urgent'' is a failure to communicate the relative value of a set of tasks.


Process friction can arise from the separation of roles needed to support a process. Those separate roles can happen within a team of bureaucrats, or the process can span distinct teams. The difficulty imposed across teams within an organization is moderated by the existence of a common arbiter in the chain of command.  
% TRANSITION to next section: processes_two
Troubleshooting processes is even more complicated when processes span organizations.  
    \section{Process Mistakes\label{sec:process-mistakes}}

Any process involving humans incurs mistakes  
\footnote{\href{https://en.wikipedia.org/wiki/Murphy\%27s_law}{Muphy's law}}. 
The purpose of pondering bureaucratic mistakes is to come up with limitations on the ability of mistakes to cause damage. 

\subsection*{Why fixes are Difficult}
Imposing checks meant to reduce mistakes creates more bureaucracy. Oversight roles carry prestige in an organization, so checks intended to limit mistakes propagate up the chain of command.


\subsection*{Why fixes Fail}
The path of information from the person with a problem to the person who can address the situation may pass through many people. 

The consequence is familiar to people who have played the \href{https://en.wikipedia.org/wiki/Chinese_whispers\%23Game}{game of telephone}.
\marginpar{[Tag] Story Time}
\index{story time!game of telelphone}
\begin{mdframed}
The person with problem explains the problem to the helpdesk staff member. The helpdesk staff member hears the problem with incomplete context and from their own frame, then relays it to their manager, who talks to the manager of the engineering team, who delegates the responsibility to the engineer. 
\end{mdframed}
This is a solvable problem: the engineer could talk with the customer. 
% source: 
% https://news.ycombinator.com/item?id=30480083

At each hand-off there is a significant chance the information is unintentionally mangled. 

\subsection*{Where Process Mistakes come from}
Mistakes in a process can originate from the bureaucrat who inflicts the process, in the hand off between participants, or from the \gls{subject} of the bureaucratic process. Typically there are more checks on subjects and fewer on bureaucrats. 

\subsection*{Choices about Mistakes}
As a process participant, I may cause a mistake or observe a mistake. My options are to can ignore it, fix it (which incurs extra work), or report it. If I chose to report it, I can do so anonymously (which decreases the risk of harm to my reputation, but also eliminates feedback) or with my name so that discussion can occur and I get feedback.
    % process change
Any Nash equilibrium is constantly being upset by the change in conditions and change in people (who have varying motives).


Processes change over time because the conditions change. Processes are implemented differently than intended because the people implementing them are not the same people who came up with and designed them


Changing complex processes is hard because individuals involved in the process may depend on steps that are not visible to other people.
Processes do not exist in isolation
\footnote{https://www.hyrumslaw.com/} % https://news.ycombinator.com/item?id=29848295

    \subsection{Deployment of Processes, aka policy update}

Design processes to prevent bad people from joining the process, detect them when they are in the process, and remove them when they do not change their behavior. Simultaneously, enable and detect good ideas and promote them and give them resources.
Promote efficiency, have enough justification so that something goes wrong there's a clear track of why without burdening too much documentation.


% https://graphthinking.blogspot.com/2016/04/if-you-want-boulder-to-roll-place-it-at.html
When designing a process, look for ways to minimize the impact on the people involved. Minimizing work (both physical and mental) for the people involved means less sensitivity to entropy in a bureaucracy. Minimizing work requires significant situational awareness on the part of the designer. In order to minimize impact on existing bureaucracy, it's vital to minimize the requirements to achieve success.


Similar to product deployment in that similar dilemmas are faced

Deployment of processes and products need to account for 
normal users, Power users, malicious users, and edge cases

\href{https://en.m.wikipedia.org/wiki/The_Innovator's_Dilemma}{Innovator's Dilemma} applies within bureaucracies to policies and processes

\href{https://en.wikipedia.org/wiki/Diffusion_of_innovations}{Diffusion of Innovation}

\subsubsection{external to the org}


\subsubsection{Internal product development and deployment\label{sec:internal_product}}

Teams in a bureaucratic organization 
\begin{itemize}
    \item consume from outside-the-organization
    \item consume from within-the-organization
    \item produce to outside-the-organization
    \item produce to within-the-organization
\end{itemize}

This section focuses on the relation between teams that produce to within-the-organization and teams that consume from within-the-organization. Tools or products that are created internally and consumed internally.

Feedback mechanisms and incentives in a non-profit monopoly. The claim of success that a team created a product that met all design requirements on-time may have no actual benefit to users. Or maybe a product that benefited internal customers was created and there were happy users, but the originating team has to quickly move onto the next project to create another success and thus has no attention to on-going support. Determining the metric of success is tricky. do you want to aim for highest average happiness of stakeholders, or are some more important?

captive users who have little leverage 




    \subsection{Processes involving two Organizations}

When two independently-governed processes interact, there can be friction. Process friction, especially two uncoordinated organizations, results in either exceptions to the process or lying or change in process. 
Example: receipts from on process submitted for reimbursement to another process.
Submitting the receipt shifts the accountability and justification burden to the accepting bureaucrat.



% https://graphthinking.blogspot.com/2017/02/financial-motivation-in-bureaucracy.html

I have dental insurance. I visited the dentist December 2 and one of the procedures during a routine cleaning was ``bitewing x-rays." This procedure is covered by my insurance, so I was surprised when I received a bill for it from my dentist.

I called my dentist and they explained that the dental insurance had declined to pay for the procedure. I called the insurance company Jan 3 and they confirmed that the procedure was covered by my policy. 
Every time I call the insurance provider I have to provide my SSN, DoB, and Zip twice -- once to the automated system, a second time to the person I talk with. The insurance company apparently had made a mistake and said they would cover the cost of the procedure. I followed up with my dentist and explained the situation.

I called the dentist to see if they had received payment yet. They had not, so I called the dental insurance provider again Feb 6 and 8. Feb 8 the insurance company said they would process the payment within 7-10 business days. I called again Feb 20 and the claim hadn't been initiated within the insurance company. I spoke to the supervisor and she said she would personally visit the claims office within the insurance company.

From the perspective of the dentist, they are seeking money for the service they provided me.

From the perspective of the insurance company, delaying payment on a claim makes good financial sense -- the policy holder is likely to just pay the balance in order to avoid going to court with the dentist.

From my perspective, the question is whether chasing this issue makes financial sense. I think of my hourly rate as $40, so after an hour the charge of $38 would have been better to pay out of pocket. Effectively I'm devaluing my time. The emotional stress and thought-cycles expended are also relevant, though harder to quantify.

    \subsection*{Tactic: Approval, Forgiveness, Opposition\label{sec:approval-forgiveness-opposition}}
% https://graphthinking.blogspot.com/2017/10/flipping-approval-mentatlity.html

Bureaucracy involves distributed decision-making. 
A common bureaucratic task is seeking consensus regarding action or spending resources. There are distinct options for how to get that consensus:
\begin{itemize}
    \item Seek approval before taking action. This approach incurs both providing justification and waiting.
    \item Ask forgiveness after taking action. Often viewed as being in contrast to seeking approval. Less delay, and usually works if things go well or if no one notices. 
    \item Notification of Intent with a deadline for response. The window for a response should be sufficient to  allow feedback. If no response is provided default is for action to be taken.
    \item Solicit opposition before taking action~\cite{2009_Perr}. This is a different framing from approval or forgiveness. It decreases the risk the approver has to take on.
\end{itemize}
The best way to proceed depends on the personalities of the people involved in building consensus and their relationships. 

Most organizations default to approval-based  processes. Each new idea needs to be signed off as approved by a sequential list of \glspl{bureaucrat}. The sequential (not concurrent) process may be known in advance, or it may be ad hoc if the request is novel.

Relying on approval is harmful to innovation because sign-off by each bureaucrat is interpreted as ``I am 100\% in agreement with this.'' Each stakeholder has to bless innovation and tie their reputation to the outcome.

Soliciting opposition and a response of ``I won't stop this'' is a more useful paradigm. With the consensus process language changed to ``I won't stop this," then the bureaucrat reviewing the idea can avoid taking responsibility for the idea and therefore is not tying their reputation to the result.

% https://news.ycombinator.com/item?id=15407757

%\ \\

% TRANSITION to next section: process_deployment
%The section above described tactics for changing processes when the changes are driven by changes to resources or staffing while assuming the underlying policy remained the same. The next section describes the case where changes to a policy require that processes be revised.  % subsection
    \subsection*{Static and Dynamic Process\label{sec:static-dynamic-processes}}

% https://graphthinking.blogspot.com/2017/04/static-versus-dynamic-processes.html

Change within an \gls{organization} \iftoggle{glossaryinmargin}{\marginpar{[Glossary]}}{} is to be expected since the external environment the organization exists in is not static. 
Sources of external change include improving technology,  changes to the \gls{shared resource}, \iftoggle{glossaryinmargin}{\marginpar{[Glossary]}}{} or shifting expectations of subjects of the bureaucracy.
Change is also driven internally to the organization by \hyperref[sec:turnover]{staff turnover}.
\marginpar{See page~\pageref{sec:turnover}.}
Since change is expected, why are static processes that are not robust to change created in the first place? Because static processes are easier to design and appear initially to require less maintenance.

Creating robust processes that are dynamic takes more effort to create. First, the process must be documented so that it can be analyzed. What is expected to happen? Who are the stakeholders? These conditions are likely to change, making the process fragile. Second, document assumptions used in the process. If the assumptions are invalidated, then the process is broken and needs to be discarded or at least revised. 

A challenge is that even when the process is documented and assumptions enumerated, there may not be an incentive to check to see if revision is necessary. Measurements (which are costly and disruptive) need to be periodically taken to see if the assumptions are still applicable. To force periodic validation of assumptions, one approach is to use \href{https://en.wikipedia.org/wiki/Sunset_provision}{sunset provisions}. 
\index{Wikipedia!\href{https://en.wikipedia.org/wiki/Sunset_provision}{sunset provisions}}
\iftoggle{WPinmargin}{\marginpar{$>$Wikipedia: sunset provisions}}{}

A more quantitative approach is to tie a process to a cost-benefit model. Implementing a process provides a benefit and comes at some cost. If the assumptions of the process can be tied to a cost-benefit model, then we can determine whether the process is worth implementing. Periodic measurements are needed to update the cost-benefit model and determine whether the process is effective.

Summarizing the steps for creating a robust process,
\begin{enumerate}
    %\item If a process already exists, document the process that is fragile.
    \item List assumptions used in the process. Who are the stakeholders, what are the goals, and what are the constraints?
    \item Relate the assumptions to a \href{https://en.wikipedia.org/wiki/Cost\%E2\%80\%93benefit_analysis}{cost-benefit model}.
    \index{Wikipedia!\href{https://en.wikipedia.org/wiki/Cost\%E2\%80\%93benefit_analysis}{cost-benefit model}}
    \iftoggle{WPinmargin}{\marginpar{$>$Wikipedia: cost-benefit model}}{}
    \item Determine the measurable parameters of the cost-benefit model. 
    \item Collect recurring measurements to verify the assumptions. 
    \item If the assumptions are broken, revise the process. 
\end{enumerate}
A robust process is just a fragile process with a feedback loop informed by ongoing measurements. Robust processes require extra work by bureaucrats compared to static processes. A static process shifts the burden to subjects -- bureaucrats have externalized the burden.
In practice, ignoring exceptions and reacting to problems is common because then there's less work for the bureaucrats enacting the process. 

\ \\

\noindent\hrulefill

\ \\

The above description of robust dynamic processes and fragile static processes characterizes workflows in isolation from the history of a team or organization. Typically processes are evolved from previous processes.  That evolution induces another source of bureaucratic friction. 
    \subsection*{Decisions and Processes Create Bureaucratic Debt\label{sec:bureaucratic-debt}}

% https://graphthinking.blogspot.com/2017/09/bureaucratic-debt-and-what-to-do-about.html

Suppose a \gls{process} is enacted and later found to be ineffective. Some work is needed to revise the process and hopefully improve effectiveness (or an \hyperref[sec:exceptions-to-process]{exception} is needed).
\ifsectionref
; see section~\ref{sec:exceptions-to-process}).
\fi
\gls{bureaucratic debt} is 
\marginpar{[Tag] Definition} the work needed to change a process\footnote{Similar to \href{https://en.wikipedia.org/wiki/Technical_debt}{technical debt}}. The bureaucratic debt is caused by choosing an easy solution now (with limited information or insufficient resources) instead of using an approach that would take longer to design and enact but be more robust.


Decisions made by \glspl{bureaucrat} occur in a resource constrained environment.
Getting information (measurement) and analysis are costly in terms of money, time, skill, and labor.
Each decision made results in options that are not explored. These missed opportunities are associated with short-term versus long-term trade-offs of costs.

The opportunity costs (options the organization doesn't take) alters which future decisions become available.

\ \\

The purpose of defining bureaucratic debt as a concept is to capture the work resulting from decisions that would otherwise be unaccounted for.
Once the concept of bureaucratic debt is understood it can be tracked.

To document bureaucratic debt, you need to record aspects of decisions as they are made:
\begin{itemize}
    \item What is the decision to be made?
    \item When the decision was identified?
    \item When the decision was made?
    \item Who made the decision?
    \item What options were identified?
    \item Which option was chosen?
    \item Why was that option was chosen over the other options?
\end{itemize}
The purpose of documenting decisions is to enable both aversion to bad decisions and attraction to good decisions. That may sound strange, but the default of decision makers is to apply the same behavior in future decisions. 
Without documenting decisions, there is no transparency, accountability, historical measure of progress, or ability to track dependencies. 

Creating a record of decisions is necessary but not sufficient. The documentation of decisions needs to be disseminated to stakeholders to enable accountability. This should occur as promptly as possible. 

Every bureaucrat exercises policies that apply to subjects, even if the subjects are other bureaucrats. What I'm describing above is beyond merely documenting what the processes and policies are for subjects. Documenting bureaucratic debt is for use internal to the team or organization.  

The scale of decision impact determines the level of documentation. ``Do I choose pencil or pen?" incurs negligible bureaucratic debt; therefore the documentation needed is also negligible. Projecting impact of decisions is a subjective prediction. 

%Similarities of tech debt and bureaucratic debt.
%In developing software, there are three artifacts: the software, documentation on how to use the software, and documentation on why to use the software. The two distinct types of documentation are typically combined in one document. Each of these three artifacts are independent. The ramification of this is that each artifact can be created independently, and it takes work to maintain synchronization of the artifacts. 

 %subsection
    \subsection*{Design Processes for Turnover of Staff\label{sec:turnover}}

% https://graphthinking.blogspot.com/2020/02/design-for-turnover-rather-than-rely-on.html

When designing a \gls{process}, there are a few goals to optimize for: time-to-first-result, average latency, initial financial cost, total financial cost, flexibility to input conditions, throughput, scalability. In a theoretical sense all these factors should inform decision making. An often neglected aspect which is harder to predict and harder to measure is the importance of employee turnover. On a long time scale, turnover of \glspl{bureaucrat} is a significant source of risk for any team or project. 

Designing processes relies on roles: participants are treated as interchangeable with other people who have similar skills. As a bureaucrat coming up with a novel process, accounting for differences in enthusiasm or communication among participants is difficult. 
Designing processes that are robust to turnover does not mean ignoring the unique talents of participants. 
Emphasize documentation (how and why) and training of new participants. 
\marginpar{[Tag] Actionable Advice}
\index{actionable advice}

\ \\

Making processes resilient to change requires  bureaucrats be educated beyond the requirements of the immediate task. That education occurs during on-boarding of the bureaucrat, while carrying out the process, and as bureaucrats exit participation in the process. 

On-boarding new bureaucrats involves technical training, explanation of norms, learning the processes, and creating a professional network of coworkers. During this on-boarding the new bureaucrat should be documenting their observations. Postponing the creation of documentation until the new person has experience results in a skewed and incomplete capture of the challenges.

Trained and experienced bureaucrats executing a process are responsible for coordinating with other bureaucrats. Bureaucrats should be cross-trained in other roles to gain familiarity with other parts of a process. 

Bureaucrats exiting the team or organization or changing roles have a responsibility to document their knowledge for team members. An exit interview can inform how the team improves. 

In each phase documentation is the mechanism for spreading knowledge. Documenting the why (in addition to the how) is critical for the reader. Access to documentation should be measured to determine whether there is value to the investment. 
  \clearpage

\clearpage
\chapter{Conclusion\label{sec:last-chapter}}

You could operate within a bureaucracy and entirely focus on just doing the job you were hired for and ignore administrative distractions. Or you can be more effective in your job by understanding the role of engaging with other people -- peer bureaucrats, supervisors, and subjects. 

Learning bureaucracy as a skill doesn't mean you can ignore personalities of individuals. Bureaucracy as a skill separate from and in addition to being a good person, being an effective member of a team, being a good project manager, being a good product owner, skillful writing, excellent verbal communication skills, applying technical skills, etc. The distinction from those aspects is that as a bureaucrat you understand the complications and constraints of your environment and then can more effectively operate within those conditions.

If you are new to being a bureaucrat, then this book armed you with understanding your environment. On your first day of employment you are poised to frame incoming information constructively.
If you are an experienced bureaucrat, it is not too late to improve. Even on your last day of employment you can learn and be more effective.

In this book I described bureaucracy and provided options for action. If you are able to apply these generalized perspectives to your specific situation, you are applying the paradigm developed throughout this book. Concepts like learning the history of your situation, identifying and engaging stakeholders to learn their perspective. Brainstorming the incentives of the individuals involved, and listing what levers they have for action. What are the dilemmas? Are there feedback loops?


\section{Skills of an Effective Bureaucrat}

The narrow scope of your role (which hopefully leverages your education/training) does not capture all relevant aspects of your job. 
Based on the definition of bureaucracy (distributed knowledge and distributed decision making), the critical aspects of success are having knowledge, sharing your knowledge, leveraging the knowledge of others, effective communication, working well with other people, understanding the role of both processes and social influence, and how those interplay. 

Because bureaucracy is a system of distributed knowledge and distributed decision making, coordination is critical. 
\begin{itemize}
    \item You are able to facilitate meetings. This means agenda, providing rules on interaction (raising hands), taking notes, sharing notes, and following up after the meeting.
    \item You participate in meetings, whether that means actively contributing or intentionally supporting other attendees. You leverage relationships with other attendees. 
    \item You focus your educational investment in written communication (emails, text-based chats, reports). Your writing empathizes with readers, captures relevant context, is concise, and is clearly worded.
    \item You communicate verbally concisely and precisely. You listen and you teach. You seek shared definitions. You are professionally vulnerable when appropriate. 
    \item If in-person with colleagues, you walk around and talk with people one-on-one.  
    \item In your role as bureaucrat you leverage project management skills: you have a vision, you make and share plans, all while building consensus with stakeholders.
    \item You apply your negotiation skills~\cite{1982_Cohen} that improve your interactions and outcomes.
\end{itemize}

The attitude of the effective bureaucrat  is that of realistic optimism. A realistic optimist will occasionally be wrong but makes progress, whereas a pessimist will be right (because the prediction is self-fulfilling) but not make progress.


% If bureaucracy is a distributed knowledge distributed decision system, compare with paxos and Byzantine generals
% https://en.wikipedia.org/wiki/Paxos_(computer_science)
%https://blog.the-pans.com/understanding-paxos/
%https://martinfowler.com/articles/patterns-of-distributed-systems/paxos.html

% https://en.wikipedia.org/wiki/Complexity_theory_and_organizations


\appendix

\chapter{Assessments for Publication}
  \section{Reader's Report\label{sec:reader_report}}
An view from inside bureaucracy as a \href{https://en.wikipedia.org/wiki/Complexity_theory_and_organizations}{complex adaptive system}. 

Central claims:
\begin{itemize}
    \item Everyone in modern society participates in bureaucracy. Therefore learning to be a skilled bureaucrat is useful.
    \item Individuals already have experience with bureaucracy from conventional roles in modern society. 
    \item Emergent behavior and being a \href{https://en.wikipedia.org/wiki/Wicked_problem}{wicked problem} make the complexity of bureaucracy irreducible to a simplistic model.
    \item The techniques of bureaucracy are meetings, processes, communication. These facilitate the coordination of distributed knowledge for distributed decision making.  
    \item Bureaucracy occurs when there is a lack of common quantitative feedback mechanism for individuals.
\end{itemize}
The consequences of thinking like a bureaucrat include
\begin{itemize}
    \item You are not limited to only the direct personal interactions that you have with other people. Bureaucratic processes that exceed your direct visibility and experience extend your influence.
    \item Help you recognize options beyond the naive defaults associated with fallacies in section~\ref{sec:fallacies}. What's feasible? What's negotiable? Once you realize rules and processes are subjectively created and enforced, negotiation (and identifying who to negotiate with) is more obvious.
    \item Decreased surprise when thing is not the way you naively expect. See the dilemmas in section~\ref{sec:dilemma_trilemma} and known hazards in section~\ref{sec:unavoidable_hazards}.
\end{itemize}


The effects of bureaucracy are not attributable to one person, but each person can improve bureaucracy.\clearpage % section
  \section{Audience Analysis}

Primary audience for this book is self-identified novice bureaucrats and people planning to work in a bureaucracy. Value of this book is advice beyond ``be a good person'' and specific to the role of a bureaucrat (and not project manager/team lead/software developer). 

Secondary audience is white collar workers who have experience working in a bureaucracy and are just realizing they are bureaucrats. This audience will find some of the information to be already known.  The value for this audience is a re-framing of the environment and problems. 

Tertiary audience is non-bureaucrats who want to better understand why bureaucratic systems are challenging. Americans consider themselves "individuals" and neglect the necessary integration of operating within a society. This book won't be a popular best seller, but it is intended for a lay audience and does not assume prior academic exposure to the topic.

Lastly, researchers of bureaucracy. This book is not a theoretical systemic analysis. \clearpage  % section

\chapter{Interview questions}

When I interview candidates for a bureaucratic position, I am seeking specific skills. Communication that facilitates coordination, working with experts from other domains, self-reflection, an ability to professionally disagree, and negotiation skills. 

\section{New Hire}

%\subsection{Questions for a Generic Bureaucrat}

\begin{itemize}
    \item Tell me about a time you had competing priorities. \\
    \textit{What to evaluate for}:
    \begin{itemize}
        \item Did the interviewee coordinate a response with other people? Teamwork is an essential skill in a bureaucratic organization.
        \item Did the interviewee clarify a misconception?
    \end{itemize}
    \item How does this position fit into your career goal?\\
    \textit{What to evaluate for}:
    \begin{itemize}
        \item Can the organization provide the candidate what they are seeking?
    \end{itemize}
    \item In previous jobs, how frequently did you check in with your supervisor (formally or informally)?\\
    \textit{What to evaluate for}:
    \begin{itemize}
        \item Does the candidate communicate with their supervisor? 
        \item Who initiates the discussion?
    \end{itemize}
    \item In previous jobs, how frequently did you check in with your coworkers (formally or informally)?\\
    \textit{What to evaluate for}:
    \begin{itemize}
        \item Does the candidate communicate with their peers? 
        \item Who initiates the discussion?
    \end{itemize}
    \item In previous jobs, in what situations did you turn to your boss for help? Or did your boss help you?\\
    \textit{What to evaluate for}:
    \begin{itemize}
        \item When the candidate's boundary is tested, how do they engage?
        \item Who initiates the discussion?
    \end{itemize}
    \item In previous jobs, in what situations did you turn to your coworkers for help? Or did your coworkers help you?\\
    \textit{What to evaluate for}:
    \begin{itemize}
        \item When the candidate's boundary is tested, how do they engage?
        \item Who initiates the discussion?
    \end{itemize}
    \item How do you engage with coworkers who know something work-related that you want to learn?
    \item Tell me about a time you were given the wrong scope (or an ill-defined scope) for a task. 
    \item In previous jobs, how did you evaluate the success and short-comings of your organization?
    \item What are the leading indicators you look for that the team you are on is headed in a bad direction?
    \item Tell me about a time you initiated a collaboration with a coworker to leverage their expertise?
    \item Tell me about a time you worked with someone who had a distinct perspective from yours. How did you collaborate?
    \item Tell me about a time when project requirements were unclear, ill-defined, or constantly shifting.
    \item Describe how you responded to a mistake you made at work.
    \item Tell me about a time the organization you were in changed and how you navigated that. 
    \item Tell me about a time when you, in your role as an expert, collaborated with people who didn't share your expertise. How does that compare with interacting with fellow experts?
    \item When you have learned a new topic in the past, what strategies did you use to balance understanding the theory and doing the practice?
    \item How have you assessed your progress when you were learning a new topic in the past?
    \item How have you structured your time when faced with a novel challenge?
    \item How do you distinguish between accidental complexity and essential complexity? \footnote{from Brook's \href{https://en.wikipedia.org/wiki/No_Silver_Bullet}{No Silver Bullet}}
\end{itemize}

%\subsection{Questions for a Generic Manager}

\section{New Team}

Prior to leaving the team you're on, you can reflect on your current position:
\begin{itemize}
    \item What are my current care-abouts?
    \item What would I miss from my current role? Are there projects, tasks, responsibilities, or people that I will miss? If yes, can I change my current role to have more of what I like?
\end{itemize}
You can negotiate with your management and coworkers to address your needs. Sometimes other people are not aware of what you want. 

Once you are interviewing for a position you should be considering the interactions as a chance to learn more about the the work and the people. When interviewing for a new position, ask the following questions\footnote{See also \href{https://www.themuse.com/advice/51-interview-questions-you-should-be-asking}{51 Great Questions to Ask in an Interview}}:
\begin{itemize}
    \item Is the team growing or is the focus shifting? \\
    (Teams hire to replace or grow. If replacing, why did the person leave? If growing, is it an increase in capacity or expansion of scope or change of focus?)
    \item What's the turn-over rate of the team I would be joining? Has anyone quit? Has anyone recently joined?
    \item How many members are their on the team I would be joining?\\
    (Is the team size stable, or does it fluctuate according to changes in tasking?)
    \item Would I be on a single team or would my work span multiple teams?
    \item How many different roles would a person in my position have? 
    
    \item What would my day-to-day routine be like?
    \item How many meetings would I have per day? Per week?

    \item What indicates success in my new team? 
    \item What indicates failure in my new team? 
    \item How am I evaluated as an employee?
    \item How is my team evaluated?

    \item Who do I report to? 
    \item How many bosses do I have?

    \item What balance of independence and collaboration is expected of me?
    \item Do I identify what to work on or is a task assigned to me? If a task is assigned, who assigns the work? (the customer or a manager or another team?)
    \item What is the typical duration of a task I would work on? \\
    (An hour or a day or a week or a month or a year.)
    \item How does my new team support my autonomy?
    \item What support does the organization provide a person in my role? (Training)
    \item In the past year, how quickly have new employees been trained in this position?
    \item What computer resources are available? Are other equipment relevant to the job? 

    \item How frequently would I interact with customers? 
    \item Who are my customers? What do my customers want?
    \item How many customers are there?
    \item How frequently would I interact with coworkers? 
    \item How frequently would I interact with boss(es)? 
    \item What are the goals for the organization? For the team?

    \item What is your favorite part of working in this organization?

    \item What are the technical challenges?
    \item How does the team document shared knowledge?
    \item Does the team have current onboarding materials to reference? Or is the process more social?
    \item What are the biggest challenges someone in this role faces? \\
    (Sometimes the challenges are political or personality rather than technical.)
    \item How diverse are the educational and experiential backgrounds of team members?
    \item What skills do other team members have?
    \item What gaps in skills does the team have that you are looking to fill?
    
    \item Where is the office I would be working in located?
    \item Is the location of the office expected to change in the next six to twelve months?
    \item Would I expected to travel?
    
    \item How many different tasks would I be expected to work on in a day? In a week?
    
    \item Do you expect the answers to any of the above questions to change in the next six to twelve months?

    \item What are the next steps in the hiring process?
    \item When will a hiring decision be made?
    \item When would I be starting in this position?
    
    % https://graphthinking.blogspot.com/2019/06/questions-to-ask-in-interview-of-your.html
    \item What are ethical issues or moral challenges should I expect to face if I were employed on this team in this organization?
    \item What are the known conflicts of interest that exist?

    \item Who is working against the progress of the team?
    
\end{itemize}

That's a long list of questions, and it doesn't even include questions specific to the role your interviewing for or the organization you're interviewing. Bringing a printed out list demonstrates thinking ahead and preparing. 

\section{Exit Interview Questions} \clearpage% chapter

% https://tex.stackexchange.com/questions/5894/latex-conditional-expression
\newif\ifbooknotes

\booknotesfalse

\ifbooknotes
\chapter{Book notes}
%\section{``Dynamics of Bureaucracy'' by Blau}

\cite{1955_Blau}

Intended audience:

Ben Payne has read this book: in progress\\
Ben Payne has a copy: yes\\
Ben Payne's assessment:

\clearpage
\section{``Values of Bureaucracy'' by Du~Gay\label{review:dugay_values}}

\cite{2005_DuGay}

Intended audience:

Ben Payne has read this book: no\\
Ben Payne has a copy: yes, electronic\\
Ben Payne's assessment:


 \clearpage
\input{notes_on_books/farazmand_handbook_of_bureaucracy}\clearpage
\input{notes_on_books/graeber_utopia_of_rules}\clearpage
\input{notes_on_books/hupe_research_handbook_on_street-level_bureaucracy}\clearpage
\section{``Street-Level Bureaucracy'' by Lipsky\label{review:lipsky_street}}

\cite{1983_Lipsky}

Intended audience:

author has read this book: yes\\
author has a copy: yes, physical\\
author's assessment:


pages 192 to the end.

Roles: customer, bureaucrat, boss of the bureaucrat

\subsection{Summary of claims}
From \cite{2015_Cooper}
\begin{quote}
Street-level bureaucracy (SLB) is a sociological theory that seeks to explain the working practices and beliefs of front-line workers in public services and the ways in which they enact public policy in their routine work. Developed by an American, Michael Lipsky,1,2 it examines the workplace in terms of systematic and practical dilemmas that must be overcome by employees, with a particular focus on public services such as welfare, policing, and education. The theory is based on the notion that public services represent ‘the coal mines of welfare where the “hard, dirty and dangerous work” of the state’ is done.’3 According to Lipsky,1,2 that is because:

\begin{itemize}
\item demand from clients will always outstrip supply due to finite resources (cost, time, or service access). Most clients are unable to obtain similar services elsewhere (such as private alternatives to state organisations). As a result, employees must resort to ‘mass processing’2 of excessive client caseloads.

\item extensive personal discretion is a critical component of the work of many front-line public sector employees, particularly those who undertake private, face-to-face interaction with clients to assess the credibility of cases. Employees must use their personal discretion to become ‘inventive strategists’ by developing ways of working to resolve excessive workload, complex cases, and ambiguous performance targets.4

    \item employees compromise the quality of their work by ‘creaming off’2 cases that are likely to be straightforward or to have a positive outcome. Alternatively, workers may act as an ‘advocate’2 for clients who are perceived as being at the tip of an iceberg of social vulnerability. Because workers are unable to offer all services to every individual they may be forced to ‘deny the basic humanity’2 of other clients. These pragmatic micro choices ultimately become the de facto policy of the organisation, which may contrast starkly with its official stated aims.
\end{itemize}
This theory has implications not just for the individual employee but also the overall system. In particular, Lipsky suggests that the extensive unmet demand from clients means that even substantial expansion of staff and budgets are unlikely to decrease workload pressures. Instead, he predicted that increased capacity would result in ongoing expansion of the same level of service quality at a higher volume.
\end{quote}


\subsection{Summary of remedies}
% https://graphthinking.blogspot.com/2018/08/how-to-decrease-bureacracy.html
In isolation, none of these ideas should be surprising.
\begin{itemize}
    \item automate recurring decision processes. This decreases bureaucracy by removing subjective influence of bureaucrats
\item where automation is infeasible, make customer advocates with end-to-end authority available
This decreases bureaucracy by improving customer's navigation of processes
\item make processes transparent to participants
This decreases bureaucracy by improving customer's understanding of processes
\item make information discoverable (e.g., via search engine) 
This decreases bureaucracy by 
\item make information directly available, rather than mediated by a person
\item after an interaction is completed, summarize the steps and outcome for the participant
\item when a process fails the needs of a participant, investigate the failure and improve the process
\item make the goals and priorities of the organization clear to all
\item define measurable standards of performance, both for individuals and teams
\item train bureaucrats how to engage participants effectively; these interactions determine the culture
\item train bureaucrats by addressing their immediate problems (e.g., through mentorship)
\item make employment desirable to people who have desirable characteristics (e.g., educated candidates)
\item enhance accountability to peers (e.g., peer review of actions and outcomes)
\item ensure that incentives for organizations and individuals encourage improvement rather than maintenance of status quo
\item decision making should be pushed down the hierarchy to the practitioner
\item when decision making requires cross-organization interaction, form a team of practitioners
\item bosses should share workload with their team in order to gain practical exposure to current challenges
\item seek feedback from process participants; then provide status updates on the implementation process
\end{itemize}
\clearpage
\input{notes_on_books/ohearn_bureaucrats_handbook}\clearpage
\section{``Bureaucracy: A Key Idea for Business and Society'' by Vine\label{review:vine_key}}

\cite{2020_Vine}

Intended audience:

Ben Payne has read this book: no\\
Ben Payne has a copy: no\\
Ben Payne's assessment:
\clearpage
\section{``Bureaucracy'' by von~Mises\label{review:vonMises_bur}}

\cite{1996_Mises}

Intended audience:

Ben Payne has read this book: no\\
Ben Payne has a copy: yes, electronic\\
Ben Payne's assessment:
\clearpage
\section{``Bureaucracy'' by Wilson}

\cite{1991_Wilson}

Intended audience: researchers

Ben Payne has read this book: no\\
Ben Payne has a copy: yes, physical\\
Ben Payne's assessment: Written from the perspective of an outside. Within that limitation it has useful analysis, though there are parts I disagree with. Well-written and easy to read.


Wilson defines ``operators'' as the street-level bureaucrats [Page 33].

Agencies typically have (ambiguous) goals, which are separated (subjectively) into tasks. [page 34]

Incentives (rewards and penalties) matter more than attitude.
[page 51]

Agencies operate under constraints set by Congress; businesses have more freedom to respond to clients.

What distinguishes business bureaucracy from government bureaucracy are feedback loops and self-determination of scope. -- From page 115.

Agencies are ``production" (ch8), or procedural/craft, or coping.
[page 245]

Procedural agencies have ambiguous outputs; Craft agencies have invisible operations
[page 250].

Coping or procedural agencies can discuss their activities but cannot verify their achievements
[page 252].

In chapter 5, agency environments were classified into four categories: majoritarian, entrepreneurial, clientist, and interest group.
[page 248].


\else

\fi

\clearpage

\printglossaries

\nocite{*} % causes LaTeX to include every entry in your .bib file.

% To investigate: bibliography per chapter
% https://stackoverflow.com/questions/2765209/latex-bibliography-per-chapter
% https://stackoverflow.com/questions/2503555/using-latex-how-can-i-have-a-list-of-references-at-the-end-of-each-section


% http://www-math.ucdenver.edu/~billups/courses/guides/annotated_bibliography.html
\bibliographystyle{plain-annote}
\bibliography{biblio_bureaucracy}%,biblio_meetings}

\printindex


\end{document}
